% !TEX program = xelatex
\documentclass[UTF8,a4paper,12pt,fontset=fandol,scheme=chinese]{ctexart}
% !TEX program = xelatex
% GMN1 中译 —— 主控导言区（改自数译 SY01-25 模板，保留原文字母宏以便逐式核对）
\usepackage[a4paper,margin=2.6cm,top=2.8cm,bottom=2.8cm]{geometry}
\usepackage{amsmath,amssymb,amsfonts,amsthm,mathtools}
\usepackage{mathrsfs}
\usepackage{bm}
\usepackage{graphicx}
\usepackage{booktabs}
\usepackage{array}
\usepackage{enumitem}
\usepackage{xcolor}
\usepackage{tcolorbox}
\usepackage{hyperref}
\usepackage{xspace}
\usepackage{tikz}
\usepackage{epigraph}

\graphicspath{{figures/}}

% ---------- 字体与行距 ----------
\linespread{1.18}
\setlength{\parskip}{0.25em}
\emergencystretch=1.2em
\sloppy

% 公式按节编号（与原刊排版一致：(1.1)、(2.18)…）
\numberwithin{equation}{section}

% ---------- 超链接 ----------
\hypersetup{
  colorlinks=true,
  linkcolor=blue!45!black,
  citecolor=green!40!black,
  urlcolor=magenta!60!black,
  pdftitle={四维穿墙：经由三维场论},
  pdfauthor={Gaiotto--Moore--Neitzke},
}

% ---------- 定理类环境（与数译体例一致：各类独立连续编号） ----------
\theoremstyle{plain}
\newtheorem{thm}{定理}
\newtheorem{lem}[thm]{引理}
\newtheorem{prop}[thm]{命题}
\newtheorem{cor}[thm]{推论}
\newtheorem{conj}[thm]{猜想}
\theoremstyle{definition}
\newtheorem{defn}{定义}
\newtheorem{exmp}{例}
\newtheorem{prob}{问题}
\theoremstyle{remark}
\newtheorem{rem}{注}
\newtheorem*{conv}{约定}
\newtheorem*{notn}{记号}

% ---------- 译注 ----------
\newcounter{transnote}
\newenvironment{yizhu}
  {\par\medskip\refstepcounter{transnote}%
   \begin{tcolorbox}[colback=orange!4,colframe=orange!55!black,boxrule=0.5pt,
                     left=4pt,right=4pt,top=3pt,bottom=3pt,arc=1pt]
   \noindent\textbf{译注 \thetransnote：}}
  {\end{tcolorbox}\par\medskip}
% 行内短译注：\zy{...} 生成脚注
\newcommand{\zy}[1]{\footnote{\textcolor{orange!70!black}{【译注】}#1}}

% ---------- 原文字母宏（与 wcf-paper.tex 对齐，便于逐式核对） ----------
\newcommand{\IC}{\mathbb{C}}
\newcommand{\IN}{\mathbb{N}}
\newcommand{\IZ}{\mathbb{Z}}
\newcommand{\IR}{\mathbb{R}}
\newcommand{\IP}{\mathbb{P}}
\newcommand{\IQ}{\mathbb{Q}}

\newcommand{\CA}{\mathcal{A}}
\newcommand{\CB}{\mathcal{B}}
\newcommand{\CC}{\mathcal{C}}
\newcommand{\CD}{\mathcal{D}}
\newcommand{\CE}{\mathcal{E}}
\newcommand{\CF}{\mathcal{F}}
\newcommand{\CG}{\mathcal{G}}
\newcommand{\CH}{\mathcal{H}}
\newcommand{\CI}{\mathcal{I}}
\newcommand{\CJ}{\mathcal{J}}
\newcommand{\CK}{\mathcal{K}}
\newcommand{\CL}{\mathcal{L}}
\newcommand{\CM}{\mathcal{M}}
\newcommand{\CN}{\mathcal{N}}
\newcommand{\CO}{\mathcal{O}}
\newcommand{\CP}{\mathcal{P}}
\newcommand{\CQ}{\mathcal{Q}}
\newcommand{\CR}{\mathcal{R}}
\newcommand{\CS}{\mathcal{S}}
\newcommand{\CT}{\mathcal{T}}
\newcommand{\CU}{\mathcal{U}}
\newcommand{\CV}{\mathcal{V}}
\newcommand{\CW}{\mathcal{W}}
\newcommand{\CX}{\mathcal{X}}
\newcommand{\CY}{\mathcal{Y}}
\newcommand{\CZ}{\mathcal{Z}}

\newcommand{\N}{\mathcal{N}}
\newcommand{\la}{\langle}
\newcommand{\ra}{\rangle}
\newcommand{\half}{\tfrac{1}{2}}
\newcommand{\abs}[1]{\lvert #1 \rvert}
\newcommand{\norm}[1]{\lVert #1 \rVert}
\newcommand{\inprod}[1]{\langle #1 \rangle}
\newcommand{\pz}{(\zeta)}
\newcommand{\ti}[1]{\textit{#1}}
\newcommand{\re}{\operatorname{Re}}
\newcommand{\im}{\operatorname{Im}}
\newcommand{\tr}{\operatorname{tr}}
\newcommand{\sgn}{\operatorname{sgn}}
\newcommand{\vol}{\operatorname{vol}}
\newcommand{\Res}{\operatorname{Res}}
\newcommand{\Tr}{\operatorname{Tr}}
\newcommand{\sfm}{\mathrm{sf}}
\newcommand{\hk}{超凯勒}
\newcommand{\qk}{四元数凯勒}
\newcommand{\kahler}{凯勒}
\newcommand{\elec}{\mathrm{e}}
\newcommand{\magn}{\mathrm{m}}
\newcommand{\degen}{\Omega}
\newcommand{\hsymp}{\varpi}
\newcommand{\clat}{\Gamma}
\newcommand{\vm}{\CB}
\newcommand{\p}{\partial}
\newcommand{\pb}{\bar\partial}
\newcommand{\inst}{\mathrm{inst}}
\newcommand{\elecm}{\mathrm{e}}
\newcommand{\cwarrow}{\text{\Large$\curvearrowright$}}
\newcommand{\ccwarrow}{\text{\Large$\curvearrowleft$}}

% —— 原文自定义宏（与原文同名，便于逐式核对；勿在各 chunk 里重复定义）——
\renewcommand\sf{{\mathrm{sf}}}
\def\tp{\tilde\phi}
\def\n{\vec{n}}
\def\nd{\vec{n}^{\dagger}}
\newcommand{\plog}[1]{#1\partial_{#1}}
\newcommand{\eps}{\epsilon}

% 译者常用的几家译名（统一在此，便于全系列一致）
\newcommand{\KS}{Kontsevich--Soibelman\xspace}
\newcommand{\SW}{Seiberg--Witten\xspace}
\newcommand{\CFT}{CFT\xspace}


\begin{document}

\pagestyle{plain}
\thispagestyle{empty}

\begin{center}
{\LARGE\bfseries 四维穿墙：经由三维场论}\\[0.35em]
{\large Four-dimensional wall-crossing via three-dimensional field theory}\\[1.0em]
{\large Davide Gaiotto$^{1}$,\quad Gregory W. Moore$^{2}$,\quad Andrew Neitzke$^{1}$}\\[0.5em]
{\small $^{1}$ 高等研究院自然科学学院，普林斯顿；\quad $^{2}$ 罗格斯大学 NHETC 与物理与天文系}\\[0.6em]
{\small 原文：arXiv:0807.4723v3\quad 66 页\quad 2009 年}\\[0.3em]
{\small 译者：数译（DSH agent）\quad 版本：v1.0\quad 2026-11-21}
\end{center}

\vspace{1.2em}
\begin{center}
\rule{0.85\textwidth}{0.4pt}
\end{center}
\vspace{0.4em}

\section*{摘要}
\addcontentsline{toc}{section}{摘要}
我们为 \SW 理论中 BPS 谱的 \KS 穿墙公式给出一个物理学解释。在此过程中，我们给出该理论在 $\IR^3 \times S^1$ 上模空间的 \hk 度量所受 BPS 瞬子修正的精确描述。穿墙公式归结为这样一个陈述：该度量是连续的。我们对度量的构造用到了二维 $tt^*$ 方程的四维类比。

\vspace{0.6em}
\begin{center}
\rule{0.85\textwidth}{0.4pt}
\end{center}

\newpage
\tableofcontents
\newpage

% ch01 —— §1 Introduction and Summary（原文 wcf-paper.tex 199–555 行）
\section{引言与摘要}
\label{sec:Introduction}

本文的主题是 $d=4$、$\CN=2$ 超对称量子场论中 BPS 态简并度的一个\textbf{穿墙公式}（wall-crossing formula，下称 WCF）。我们的约定与相关定义的汇总见第 2 节。电荷为 $\gamma$ 的 BPS 态空间 $\CH_{\gamma,\mathrm{BPS}}$，指的是单粒子 Hilbert 空间中的这样一个空间：其中各态的电磁荷为 $\gamma$，且使 BPS 界 $M \geq \abs{Z_\gamma(u)}$ 取等号。这里 $u$ 记矢量多重态模空间 $\vm$（即真空模空间的库仑分支）中的一点。

对 $d=4$、$\CN=2$ 超对称而言，唯一可用的指标是第二螺旋度超迹：
\begin{equation} \label{eq:index}
\degen(\gamma;u) := -\half\,{\Tr}_{\CH_{\mathrm{BPS},\gamma}}(-1)^{2J_3}(2J_3)^2,
\end{equation}
其中 $J_3$ 是有质小群转动子群的任一生成元。人们早已知道，这类指标一般并不独立于 $u$，而只是分段常值 \cite{Cecotti:1992qh}。事实上，$\degen(\gamma;u)$ 会在边缘稳定墙处跳变，在那里 $\gamma=\gamma_1+\gamma_2$ 且 $\arg Z_{\gamma_1}(u)=\arg Z_{\gamma_2}(u)$。这一事实在 \SW 理论的发展中起过重要作用 \cite{Seiberg:1994rs,Seiberg:1994aj}。

近些年来，人们对这一指标如何依赖于 $u$ 开始有了更为系统的理解。当衰变 $\gamma\to\gamma_1+\gamma_2$ 中至少有一个组分是原始的（primitive）时，\cite{Denef:2007vg} 给出了跨过边缘稳定墙时变化量 $\Delta\degen$ 的公式。这些原始与半原始穿墙公式，是由基于超引力多中心解的物理图像导出的 \cite{Denef:2000nb,Denef:2000ar}。然而，当两个组分都带非原始电荷时，\cite{Denef:2007vg} 的方法便难以施展。

Kontsevich 与 Soibelman \cite{ks1} 提出了一个适用于一切可能衰变情形的、关于 $\Delta\degen$ 的漂亮穿墙公式。我们将在第 2.2 节回顾他们的公式，下文有时称之为 KS 公式。

一方面，Kontsevich--Soibelman 的「Donaldson--Thomas 不变量」$\hat\degen(\gamma;u)$ 与物理上所关心的 $\degen(\gamma;u)$ 显然并不相同，他们导出公式所用的技巧也似乎与标准的物理考量相去较远。另一方面，与 \cite{Denef:2007vg} 中半原始穿墙公式的表述相比，他们的 WCF 牵涉到一些引人注目的新概念。特别是，WCF 是用环面的一族辛同胚的某个乘积来表达的（见下文 (2.18)），该乘积依赖于诸 $\hat\degen(\gamma;u)$，从而\emph{先验地}依赖于 $u$。WCF 的陈述是：这个乘积实际上与 $u$ 无关。这反过来就决定了 $\hat\degen(\gamma;u)$ 对 $u$ 的依赖方式。这一进展提出了 KS 公式的物理推导与物理诠释问题，也预示着其中含有某些本质性的新物理思想。事实果然如此。

在本文中，我们针对 $d=4$、$\CN=2$ 的\emph{场论}给出 KS 公式的物理解释与证明。推广到超引力是一个有趣而重要的问题，留待将来。

以下是主要思想与基本策略的梗概。考虑该规范理论处在空间 $\IR^3\times S^1$ 上，其中 $S^1$ 的半径为 $R$。低能下该理论由以超凯勒流形 $(\CM,g)$ 为靶空间的 $d=3$ sigma 模型描述。该 sigma 模型受到 BPS 瞬子的修正：所谓 BPS 瞬子，就是 $d=4$ 理论中某个 BPS 粒子的世界线绕 $S^1$ 缠绕一圈所形成的组态。于是，把度量 $g$ 在大 $R$ 处展开，便可读出 BPS 粒子的简并度 $\degen(\gamma;u)$。这立刻引出一个困惑：我们知道诸 $\degen(\gamma;u)$ 是不连续的，而 $g$ 却应当是连续的！度量的连续性基于这样一条物理原则（它在 \cite{Seiberg:1994rs,Seiberg:1994aj} 中至关重要）：低能有效场论拉氏量中的奇异性，只能来自在特殊模处出现的无质量粒子——这些粒子本就不该被从有效理论中积分掉。

物理上，这一困惑的化解类似于 \cite{GarciaEtxebarria:2007zv} 中最近讨论的情形。精确的度量 $g$ 确实是光滑的，但它既受单粒子态修正，也受多粒子态修正。当某个粒子衰变时，其单瞬子贡献的消失，正好被它的衰变产物所带来的多瞬子贡献中的不连续性所补偿。同样地，消失的 $n$-瞬子贡献被 $m>n$ 的 $m$-瞬子贡献中的不连续性所补偿。更精确地说，$n$-瞬子修正具有如下形式
\[
\sum \prod_{i=1}^{n}\degen(\gamma_i;u)\,F^{(n)}_{(\gamma_i)}(u),
\]
其中求和遍历电荷的一切 $n$ 元组 $\{\gamma_i\}$，而 $F^{(n)}$ 本质上是 $R$ 与诸 $Z_{\gamma_i}$ 的普适函数。跨过墙时，每个 $\degen(\gamma;u)$ 都有一个不连续量，它正比于满足 $\sum\gamma_j=\gamma$ 的诸 $\degen(\gamma_j;u)$ 之乘积的和；与此同时，诸函数 $F^{(n)}$ 也带有不连续量，正比于下标 $n'<n$ 的诸 $F^{(n')}$。我们将会看到，Kontsevich--Soibelman 穿墙公式所表达的，正是这一连串抵消的自洽性。

因此，理解 WCF 的主要技术障碍就在于：对 BPS 瞬子给 $g$ 带来的修正给出一种高效的描述。超凯勒度量是个复杂的对象，若去研究它各个分量的修正，几乎难有进展；一般说来，也没有像凯勒势那样简单的可加对象可用。为克服这一困难，我们借用扭量理论的一些想法。回忆：一个超凯勒流形相对于整整一个 $\IC\IP^1$ 的复结构都是复辛的。基本想法是：研究 $g$ 等价于研究 $\CM$ 上的全纯达布坐标——只要我们同时考虑\emph{全部}这些复结构。

在本文主体中，我们假定 Kontsevich--Soibelman 穿墙公式对 $\degen(\gamma;u)$ 成立。在这一假定下，我们构造 $\CM$ 上的度量，办法是给出一族标准函数 $\CX_\gamma(u,\theta;\zeta)$，它们定义在 $\CM\times\IC^\times$ 上，并由电磁荷 $\gamma$ 标记。这里 $(u,\theta)$ 标定 $\CM$ 中一点，参数 $\zeta$ 标记 $\CM$ 上的复结构。每个 $\CX_\gamma$ 关于 $\zeta$ 是\emph{分段}全纯的；BPS 瞬子的效应就是在 $\zeta$ 平面上沿某些射线 $\ell$ 制造出 $\CX_\gamma(u,\theta;\zeta)$ 的不连续性。这些不连续性被等同于 Kontsevich 与 Soibelman 引入的那些辛同胚。在这一进路下，度量的连续性成为 WCF 的推论。最后一节我们把论证倒过来运行：利用超对称规范理论的一般原理，推出度量 $g$ 的若干性质，它们足以证明 WCF。

\subsection*{小结}

我们首先在第 2 节回顾 $d=4$、$\CN=2$ 规范理论的 \SW 解以及 Kontsevich--Soibelman 穿墙公式，然后讨论该理论在 $\IR^3\times S^1$ 上的表述。它是一个映入流形 $\CM$ 的 sigma 模型；$\CM$ 在拓扑上是 $d=4$ 模空间 $\vm$ 上的 \SW 环面丛，并装备一个超凯勒度量 $g$。该度量依赖于 $S^1$ 的半径 $R$。当 $R\to\infty$ 时它趋于一个简单形式，该形式可由 $d=4$ 理论的朴素量纲约化得到；我们称这个简单度量为 $g^{\sfm}$（sf 取自 ``semiflat''）。

在第 3 节我们解释构造超凯勒度量的「扭量」方法：给定 $\CM$ 上一族函数 $\CX_\gamma(u,\theta;\zeta)$，它们随 $\zeta\in\IC^\times$ 全纯变化并满足若干附加条件，就存在一个超凯勒度量，使得这些 $\CX_\gamma(u,\theta;\zeta)$ 成为它的全纯达布坐标。特别地，我们给出对应于半平坦度量 $g^{\sfm}$ 的函数 $\CX^{\sfm}_\gamma(u,\theta;\zeta)$。

有了这些背景，我们就可以考虑瞬子修正。第 4 节从最简单的情形开始：$U(1)$ 规范理论耦合一个电荷为 $q$ 的物质超多重态。在这一理论中，修正后的度量 $g$ 是精确已知的 \cite{Ooguri:1996me,Seiberg:1996ns}。我们说明如何通过纳入瞬子修正、把函数 $\CX^{\sfm}_\gamma(u,\theta;\zeta)$ 改造为新的 $\CX_\gamma(u,\theta;\zeta)$ 来得到这一修正度量。在这一构造中，Kontsevich--Soibelman 公式的构件已经现身：我们的 $\CX_\gamma(u,\theta;\zeta)$ 自然而然地带有 $\zeta$ 平面中的不连续性，而这些不连续性恰恰就是对应于电荷 $\pm q$ 的初等 Kontsevich--Soibelman 辛同胚。

随后，第 5 节转向更有意思的情形：存在多种 BPS 瞬子修正，它们来自 $d=4$ 中相互非局域的 BPS 粒子。此时我们为 $\CX_\gamma(u,\theta;\zeta)$ 找到一个自然的拟设（ansatz）：本质上，只要求每个 BPS 粒子各自独立地贡献一个不连续性，与我们为单个粒子找到的那一个同型。这一不连续性最自然地落在 $\zeta$ 平面中由该 BPS 粒子中心荷的相位所决定的一条射线上。相互非局域粒子的 Kontsevich--Soibelman 因子并不对易；但一般说来这不构成问题，因为这些粒子的中心荷互不共线，它们的不连续性于是出现在 $\zeta$ 平面中彼此不同的射线上。射线的分离恰在边缘稳定墙上消失；在那里，诸不连续性堆积成 Kontsevich--Soibelman 因子的乘积。WCF 说的正是：从墙的两侧分别趋近时，这一乘积相同。这一要求对我们至关重要：它蕴含着由 $\CX_\gamma(\zeta)$ 构造出的度量是连续的。

更确切地说，为确定 $\CX_\gamma(u,\theta;\zeta)$，我们既要指定它们在 $\zeta$ 平面中的不连续性，又要指定它们在 $\zeta\to0,\infty$ 时的渐近行为。换言之，我们提出一个无穷维的「黎曼--希尔伯特问题」，其解就是诸 $\CX_\gamma(u,\theta;\zeta)$。我们并不精确地构造它的解；我们采用的策略与 Cecotti 和 Vafa 所用者极为接近——他们在研究 $d=2$、$\CN=(2,2)$ 超对称理论时遇到过类似的（但为有限维的）黎曼--希尔伯特问题 \cite{Cecotti:1993rm}。把他们论证稍作变形，就足以证明我们的问题至少在 $R$ 充分大时存在解。事实上，在大 $R$ 极限下，所需的 $\CX_\gamma$ 可以逐次逼近得到：第零阶近似就是 $\CX^{\sfm}_\gamma$，第 $n$ 阶近似则把至多 $n$ 个瞬子的多瞬子效应纳入进来。

在构造出诸 $\CX_\gamma(u,\theta;\zeta)$、从而构造出度量 $g$ 之后，我们核验 $g$ 具有若干基于一般场论理由所能预期的性质；它通过了全部检验，因此我们论证它应当就是 $\CM$ 上正确的物理度量，这推广了 \cite{Seiberg:1996nz} 中的一个类似论证。

如前所述，我们对 $\CX_\gamma(u,\theta;\zeta)$ 的构造与 $d=2$ 情形 \cite{Cecotti:1993rm} 中出现的构造有着惊人的相似。在那一情形，BPS 畴壁简并度的穿墙公式是利用 $d=2$ 理论真空丛上平坦的「$tt^*$ 联络」证明的。该联络的两个分量给出表达 $d=2$ 理论 $R$ 对称性与标度不变性的微分方程。我们的构造同样可以表述为：在 $\CM_u$（即 $\CM$ 的环面纤维）上实解析函数的无穷维丛中，$\vm\times\IC\IP^1\times\IR_+$ 上的一个平坦联络 $\CA$。黎曼--希尔伯特构造保证了这一 $\CA$ 的存在性。每个 $\CX_\gamma$ 定义一个平坦截面。特别地，这一平坦性给出关于 $\CX_\gamma$ 对 $\zeta$ 与 $R$ 的依赖性的—对微分方程，它们在物理上源自 $d=4$ 理论的 $R$ 对称性反常与标度不变性。

正如我们在第 7 节所述，这个类 $tt^*$ 的平坦联络，仅利用超对称规范理论的一般原理就能被发现。而且，它仅仅\emph{存在}这一事实，就足以先验地证成我们对度量的拟设。特别地，穿墙公式原本是在该拟设内部作为一条自洽性要求而出现的，现在它可以理解为由 $\CA$ 构造出的一个「等单值形变」的存在性。这就给出了穿墙公式的一个物理证明。

为方便起见，本文大部分篇幅使用不包含味对称信息的简单形式的穿墙公式，相应地也把所有味质量取为零。第 6 节将说明如何恢复味荷与质量的信息。

我们附上若干包含更多细节的附录。附录 A 说明如何直接验证穿墙公式在纯 $SU(2)$ 理论情形给出正确的 BPS 简并度。附录 B 描述 $\CM$ 上的柯西--黎曼方程，其表述与我们对超凯勒度量的构造、以及 \cite{Cecotti:1991me} 的 $tt^*$ 方程相衔接。附录 C 给出从我们的黎曼--希尔伯特问题中提取度量大 $R$ 修正所需的渐近分析。附录 D 讨论从黎曼--希尔伯特问题的解中提取微分方程的若干细节。最后，附录 E 解释我们的主要结果之一，即方程 (5.13)，与热力学贝特拟设之间的一个有趣联系。这一联系还有更多可说的，但留待另一次。

文中有若干小节专门讨论与 KS 公式中一个细微却重要的符号相关的整体性问题。初读时跳过这些讨论是合理的。选择这样做的读者，尽可允许自己在正文中把 $T$ 与 $\tilde T$、以及 $\CM$ 与 $\widetilde{\CM}$ 混为一谈。

\subsection*{讨论}

我们来点评几个特别有意思的地方。

\begin{itemize}
\item 物理上，我们对 $\CM$ 上度量的构造，相当于给出了一条在 $d=4$ 中把相互非局域的粒子「积分掉」的规则。这个问题\emph{先验地}看来很难，因为找不到任何一个使所有粒子都带电荷的对偶标架，从而很难写出包含全部相关场的拉氏量。这里我们绕开了这一困难。

\item 我们对度量的构造用的是它的扭量描述。扭量空间出现的最自然的物理语境是射影超空间 \cite{Karlhede:1984vr,Lindstrom:1987ks,Lindstrom:2008gs}，其中参数 $\zeta$ 是一个玻色型超空间坐标。$g$ 的修正只来自 BPS 瞬子、并且定域在 $\zeta$ 平面中特定的射线上，这一事实在射影超空间的语言中应当有自然的解释。

\item Kontsevich--Soibelman WCF 的灵感之一，是他们早年的工作 \cite{MR2181810}：他们通过「修正」半平坦 K3 上的函数层，给出了 K3 曲面上全纯函数层的一个构造。修正是用某些辛同胚的乘积来表述的，与穿墙公式中出现的那些类似。这一构造与我们的构造密切相关，只需把 K3 换成 $\CM$。我们工作中的关键新要素是同时考虑所有复结构，从而引入参数 $\zeta\in\IC\IP^1$；一旦这样做了，我们就能写出那个至关重要的黎曼--希尔伯特问题。这个想法或许在原来的 K3 语境中也有用。

\item $g$ 的多瞬子展开是一些基本构件之和，各构件以 BPS 简并度 $\degen(\gamma;u)$ 的乘积为权重。这些基本构件在边缘稳定墙处有精细的不连续行为，它与 $\degen(\gamma;u)$ 的跳变相互配合，使 $g$ 关于 $u$ 连续。这一切令人想起 Joyce 关于穿墙的近期工作 \cite{Joyce:2006pf}。此外，Joyce 的工作被 Bridgeland 与 Toledano Laredo 在 \cite{stab-stokes} 中解释为 $\IC\IP^1$ 上某个联络的等单值形变；该联络与我们此处考虑的联络有些相像，但形式略有不同：它只在 $t=0$ 处有非正则奇点，而我们的联络在 $\zeta=0$ 与 $\zeta=\infty$ 两处都有非正则奇点。我们的联络有一个有趣的标度极限：$R\to0$、$\zeta\to0$ 而 $\zeta/R=t$ 固定，它把我们的联络带到 \cite{stab-stokes} 中联络的形式（尽管结构群不同）。这一极限保留了关于 BPS 简并度及其穿墙的信息。看看是否存在某种意义把我们的联络与 \cite{stab-stokes} 中的联系起来，将会很有意思。

\item 我们的讨论中研究的是定义在矢量多重态模空间 $\vm$ 上的结构。然而 Kontsevich--Soibelman 与 Joyce 都把他们的不变量表述在一个更大的空间上，即「Bridgeland 稳定性条件」的空间 \cite{MR2373143}。我们尚不理解，把这些构造推广到该更大空间时的含义为何。

\item Kontsevich--Soibelman 所表述的穿墙公式，不仅对 $\CN=2$ 场论、而且对超引力都有意义，事实上后者正是 \cite{Denef:2007vg} 的主要关注点。此时理论在 $\IR^3\times S^1$ 上的模空间 $\CM$ 是四元数凯勒流形而非超凯勒流形。尽管如此，经由适当修改，我们的多数考量在该语境中似乎仍然说得通。例如，Hitchin 定理被 LeBrun 定理取代，后者用全纯接触结构来刻画四元数凯勒流形的扭量空间。特别地，仍然有「全纯」函数 $\CX_\gamma(x,\zeta)$ 的自然概念（即 $\CM$ 的扭量空间上的全纯函数），而 $\CX^{\sfm}_\gamma$ 的四元数凯勒类比已在 \cite{Neitzke:2007ke} 中求出。我们预期 $\CM$ 上受瞬子修正的度量 $g$，可以用与本文平行的方法得到：为 $\CX_\gamma$ 写出黎曼--希尔伯特问题，用 $\CX^{\sfm}_\gamma$ 定渐近行为，用 Kontsevich--Soibelman 因子定不连续性。需要克服的一个重要困难是：在引力情形，简并度 $\degen(\gamma;u)$ 随 $\gamma$ 增长得极快，这使得迭代解 $\CX_\gamma$ 的收敛性在此情形下不那么显然。与超凯勒情形一样，WCF 应当作为保证 $g$ 光滑的一条自洽性条件而出现。

\item 丛 $\CM\to\vm$ 的超凯勒几何与 \cite{Cecotti:1991me,Cecotti:1992qh,Cecotti:1993rm} 的 $tt^*$ 几何之间的类似是惊人的：两种结构非常相像，尽管一个关乎 $d=4$ 的场论，另一个关乎 $d=2$。两者之间有没有直接的联系？一种可能是仅通过紧致化把它们联系起来，例如在 $S^2$ 上紧致化。$S^2$ 上 $U(1)$ 磁通的不同取值于是对应于 $d=2$ 理论的不同真空，而 $d=4$ 理论的 BPS 态则可被等同于 $d=2$ 中在这些真空之间插值的畴壁。\footnote{这一图像是 Cumrun Vafa 向我们倡导的。}相关的想法此前已在文献中出现过——特别见 \cite{Losev:1997tp,Losev:1999nt,Ooguri:2005vr}。另见 \cite{Dorey:1999zk,Ritz:2006rt}，其中给出了 $d=2$ 与 $d=4$ 中 BPS 谱之间一种略有不同的联系。

\item 包含半平坦几何在内的一类超凯勒流形的无穷小形变，最近在 \cite{Alexandrov:2008ds} 中得到了研究。用他们的语言描述半平坦几何的首阶修正将会很有意思。我们的方程 (4.33) 与他们的 (3.38) 相像，只要适当选取 $H$ 与积分围道。
\end{itemize}

\begin{yizhu}
本章是全书译稿的样板，以下术语与记号在此一并交代，后文沿用。
\begin{enumerate}[leftmargin=2em,itemsep=2pt,topsep=2pt]
\item \textbf{墙与穿墙}：wall 译「墙」，wall-crossing 与 wall-crossing formula 依数译既有译法作「穿墙」「穿墙公式」（不用「墙穿越」）；wall of marginal stability 作「边缘稳定墙」（不用「边际稳定」）。BPS 简并度指标 $\degen(\gamma;u)$ 即原文的 \verb|\degen|，为第二螺旋度超迹，沿用 $\Omega$ 这一符号；Kontsevich--Soibelman 的不变量则记为 $\hat\degen(\gamma;u)$，正文中称 $\hat\Omega$ 或「KS 不变量」。
\item \textbf{情形与语境}：原文以「$d=4$ 理论放在 $\IR^3\times S^1$ 上」的方式过渡到三维，故 instanton 一词在此指「BPS 粒子沿 $S^1$ 缠绕」所形成的组态，译「BPS 瞬子」，与四维瞬子区别开；multi-instanton 作「多瞬子」，$n$-instanton 作「$n$-瞬子」。
\item \textbf{人名与缩写}：Kontsevich--Soibelman 首次出现用全名，其后作「KS 公式」「KS 因子」「KS 辛同胚」；Seiberg--Witten 作「\SW 理论」；人名之间用 en dash（--），与中文连接号区分。
\item \textbf{交叉引用}：本章提及的公式 (2.18)、(4.33)、(5.13) 与各节号均按原著排版给出；合稿时统一改为自动编号，以免因增删文字而错位。
\end{enumerate}
\end{yizhu}

% ch02 —— §2 Preliminaries（原文 wcf-paper.tex 559–1164 行）
\section{预备知识}
\label{sec:preliminaries}

\subsection{$d=4$、$\CN=2$ 规范理论}

我们考虑一个 $d=4$、$\CN=2$ 超对称的规范理论：规范群 $G$ 的秩为 $r$，并带有一个特征（复）质量标度 $\Lambda$。\SW 理论（由 \cite{Seiberg:1994rs,Seiberg:1994aj} 开创，更一般的综述见如 \cite{Lerche:1997sm,Freed:1997dp,ias-volumes}）对这类规范理论在其库仑分支上、能量 $\mu \ll \Lambda$ 时的行为给出了相当完整的描述，具体如下。

库仑分支是一个复维数为 $r$ 的复流形 $\vm$（4d 模空间），由矢量多重态标量的真空期望值参数化。我们把 $\vm$ 上的一般坐标系记作 $(u^1, \dots, u^r)$。在每一点 $u \in \vm$，规范群破缺为极大环面 $U(1)^r$。存在一个电磁荷构成的格 $\clat_u \simeq \IZ^{2r}$，其上装备着取整数值的辛配对 $\inprod{,}$。这个格是 $\vm$ 上一个局部系统 $\clat$ 的纤维。也就是说，在 $u \in \vm$ 之上有一个以 $\clat_u$ 为纤维的格的丛，它绕 $\vm$ 中某些奇异轨迹具有非平凡的单值性——那些轨迹的复余维为 $1$，其上有些 BPS 粒子变成无质量。我们有时不严谨地写「$\gamma \in \clat$」，意思是说 $\gamma$ 是 $\clat$ 的一个局部截面。

存在一个「周期」向量 $Z(u) \in \clat_u^* \otimes_\IZ \IC$，它随 $u$ 全纯地变化。对任意 $\gamma \in \clat$，我们如下定义中心荷 $Z_\gamma(u)$：
\begin{equation}
Z_\gamma(u) = Z(u) \cdot \gamma.
\end{equation}
无论是无质量部分还是有质量部分，$Z(u)$ 的描述都起根本的作用。

我们从无质量部分开始。在 $\vm$ 上局部地，可以把 $\clat$ 分裂为 $\clat = \clat^m \oplus \clat^e$，其中两个子格分别由「磁」荷与「电」荷构成，且都是拉格朗日子格。利用 $\clat$ 上的配对，$\clat^m$ 与 $\clat^e$ 互为对偶。这样的分裂称为一个电磁对偶标架。具体地，我们可以为 $\clat^m$ 取一组基 $\{\alpha_1, \dots, \alpha_r\}$，为 $\clat^e$ 取一组基 $\{\beta^1, \dots, \beta^r\}$，使得
\begin{equation}
\inprod{\alpha_I, \alpha_J} = 0, \quad \inprod{\beta^I, \beta^J} = 0, \quad \inprod{\alpha_I, \beta^J} = \delta_I^J
\end{equation}
其中 $I, J = 1, \dots, r$。选定这样的标架后，我们在 $\vm$ 上得到一组「特殊坐标」$a^I$；它们不是别的，正是电中心荷，即
\begin{equation} \label{eq:special-coords}
Z_{\beta^I} = a^I.
\end{equation}
磁中心荷于是是诸 $a^I$ 的全纯函数。它们由单个函数 $\CF(a^I)$（依赖于所选标架）——即 $\CN=2$ 预势——决定：
\begin{equation} \label{eq:prepot}
Z_{\alpha_I} = \frac{\partial \CF}{\partial a^I}.
\end{equation}
这特别蕴含 $Z$ 并非任意：由混合偏导数的对称性可得
\begin{equation} \label{eq:lag-constraint}
\inprod{dZ, dZ} = 0.
\end{equation}
在 \eqref{eq:lag-constraint} 的左端，我们既用了反对称配对 $\inprod{,}$，又用了 $\vm$ 上 $1$-形式的反对称楔积；两者复合而成的配对是对称的，所以这个条件并非空话。事实上，\eqref{eq:lag-constraint} 说的是：在 $u$ 附近，$\CB$ 可以局部地等同于 $\Gamma^*_u \otimes_\IZ \IC$ 的一个复拉格朗日子流形。

预势完全决定了二导数阶的有效拉氏量，后者用电源矢量多重态写出。为写出这个拉氏量，我们引入由
\begin{equation}
\tau_{IJ}(a) = \frac{\partial^2 \CF}{\partial a^I \partial a^J},
\end{equation}
定义的对称矩阵 $\tau$，然后采用一种略去规范指标的记号，例如用 $\tau \abs{da}^2$ 表示 $\tau_{IJ} da^{I} \wedge \star d\bar{a}^J$。于是拉氏量的玻色部分为
\begin{equation} \label{eq:sw-lagrangian}
\CL^{(4)} = \frac{\im \tau}{4\pi} \left(- \abs{da}^2 - F^2 \right) + \frac{\re \tau}{4\pi} F \wedge F.
\end{equation}

中心荷 $Z_\gamma$ 对有质量谱同样具有根本的重要性。事实上，任何电荷为 $\gamma$ 的单粒子态，其质量都满足
\begin{equation} \label{eq:bps-bound}
M \ge \abs{Z_\gamma}
\end{equation}
且等号成立当且仅当该态是 BPS 态。BPS 态属于超庞加莱对称的有质量短多重态；在小群 $SU(2)$ 下，这样一个多重态中的静止态变换为
\begin{equation}
[j] \otimes \left(\left[1/2\right] + 2[0]\right).
\end{equation}
取 $j=0$ 给出有质量超多重态，而 $j = \half$ 则是有质量矢量多重态。

有一个标准的指标可以「数」短多重态，即第二螺旋度超迹 $\degen(\gamma;u)$。对 $u \in \vm$ 处理论谱中每一个电荷为 $\gamma$ 的有质量超多重态，这个超迹贡献 $+1$；类似地，对每一个有质量矢量多重态贡献 $-2$。在任何使单粒子态不与多粒子态连续谱混合的形变下，$\degen(\gamma;u)$ 不变。由 \eqref{eq:bps-bound} 可知，这种混合受到很强的限制：一个 BPS 粒子只能衰变为其他 BPS 粒子，而且只有当它们的中心荷全都具有同一相位时才行。因此，在 $\vm$ 中远离\emph{边缘稳定墙}处，$\degen(\gamma;u)$ 关于 $u$ 局部常值。这些边缘稳定墙的实余维为 $1$；对一对线性无关的电荷 $\gamma, \gamma'$，它们定义为使 $Z_\gamma$ 与 $Z_{\gamma'}$ 非零且具有同一相位的诸 $u\in \vm$ 所构成的轨迹。

理解 $\degen(\gamma;u)$ 在 $u$ 穿过边缘稳定墙时的跳变行为，是本文的主要动机之一。下面我们就转向它。

\subsection{Kontsevich--Soibelman 穿墙公式}
\label{sec:ReviewKS}

本节回顾 Kontsevich--Soibelman 穿墙公式。如 \cite{ks1} 最初所提出的，该公式决定了「广义 Donaldson--Thomas 不变量」$\hat\degen(\gamma;u)$ 的跳变行为。下文将会看到，若把 Donaldson--Thomas 不变量与螺旋度超迹等同起来，即 $\hat\degen(\gamma;u) = \degen(\gamma;u)$，那么穿墙公式在若干非平凡的例子中给出物理上所预期的答案：特别地，它重现了 \cite{Denef:2007vg} 的「原始穿墙公式」，也重现了规范群为 $SU(2)$ 的 \SW 理论之 BPS 谱的穿墙行为。

有一个技术要点：为使 KS 公式有意义，诸 $\degen(\gamma;u)$ 不能完全任意。在 $\clat$ 上引入一个正定范数后，必须要求存在某个 $K > 0$，使得
\begin{equation}
\frac{\abs{Z_\gamma}}{\norm{\gamma}} > K
\end{equation}
对所有满足 $\hat\degen(\gamma;u) \neq 0$ 的 $\gamma$ 成立。本文通篇假定这条性质——称为「支撑性质」（Support Property）——成立。

\subsubsection*{Kontsevich--Soibelman 代数}

穿墙公式是用一个李代数给出的；该李代数由生成元 $e_\gamma$（$\gamma \in \clat$）以及一条基本对易关系定义：
\begin{equation} \label{eq:ks-commutator}
[e_{\gamma_1}, e_{\gamma_2}] = (-1)^{\langle
\gamma_1,\gamma_2\rangle} \langle \gamma_1,\gamma_2\rangle
e_{\gamma_1+\gamma_2}.
\end{equation}
本文将把这个抽象的李代数实现为一个复化环面上的复辛同胚代数，这一点很重要。除去本节末将出现的一个细微之处，该环面就是局部系统 $\tilde{T} := \clat^* \otimes_{\IZ} \IC^\times$ 的纤维 $\tilde{T}_u$。

任意 $\gamma \in \clat$ 给出 $\tilde{T}_u$ 上一个相应的函数 $X_\gamma$，满足 $X_\gamma X_{\gamma'} = X_{\gamma + \gamma'}$。取定 $\clat$ 的一组基 $\{\gamma^1, \dots, \gamma^{2r}\}$ 后，我们可以取 $X^i := X_{\gamma^i}$ 作为 $\tilde{T}_u$ 的坐标。$\clat^*$ 上的辛配对给出 $\tilde{T}_u$ 上一个全纯辛形式 $\hsymp^{\tilde{T}}$：若 $\epsilon^{ij} = \inprod{\gamma^i, \gamma^j}$，而 $\epsilon_{ij}$ 是它的逆，则
\begin{equation}
\hsymp^{\tilde{T}} = \frac{1}{2}\epsilon_{ij} \frac{dX^i}{X^i} \wedge \frac{dX^j}{X^j}.
\end{equation}

我们希望把 $e_\gamma$ 等同于 $\tilde{T}_u$ 上由哈密顿量 $X_\gamma$ 生成的无穷小辛同胚。这差不多给出代数 \eqref{eq:ks-commutator}，但漏掉了额外的符号 $(-1)^{\langle \gamma_1,\gamma_2\rangle}$。这个符号在下文与物理上已知的诸穿墙公式作比较时很重要；在那种语境下，它与下面的事实有关：由两个电荷为 $\gamma_1,\gamma_2$ 的粒子组成的束缚态，其费米子数会平移 $\langle \gamma_1,\gamma_2\rangle$。

在 $\vm$ 的一个局部片上，我们可以通过引入 $\IZ_2$ 值二次型 $(-1)^{\langle \gamma_1,\gamma_2\rangle}$ 的一个「二次加细」（quadratic refinement）来吸收这个符号：所谓二次加细，是指一个满足
\begin{equation}
\sigma(\gamma_1)\sigma(\gamma_2) = (-1)^{\langle \gamma_1,\gamma_2\rangle} \sigma(\gamma_1 + \gamma_2)
\end{equation}
的映射 $\sigma: \clat \to \IZ_2$。得到这样的 $\sigma$ 的一个办法是：取一个局部电磁对偶标架 $\clat \cong \clat^e \oplus \clat^m$，写 $\gamma = \gamma^e + \gamma^m$，并置 $\sigma(\gamma)=(-1)^{\inprod{\gamma^e,\gamma^m}}$。注意
\begin{equation}
\sigma(\gamma_1)\sigma(\gamma_2) = \sigma(\gamma_1 + \gamma_2)
(-1)^{\inprod{\gamma_1^e,\gamma_2^m} + \inprod{\gamma_2^e,\gamma_1^m}} =
\sigma(\gamma_1 + \gamma_2) (-1)^{\langle \gamma_1,\gamma_2\rangle}
\end{equation}
正如所需。无论如何，选定了某个 $\sigma(\gamma)$ 之后，我们便可以把 $e_\gamma$ 等同于由哈密顿量 $\sigma(\gamma) X_\gamma$ 生成的辛同胚。

任意两个加细 $\sigma$、$\sigma'$ 都满足 $\sigma(\gamma) \sigma'(\gamma) = (-1)^{c(\sigma, \sigma') \cdot \gamma}$，其中 $c(\sigma, \sigma') \in \clat^* / 2 \clat^*$ 是某个固定的元素。与这两个加细相联系的哈密顿量 $\sigma(\gamma) X_\gamma$ 与 $\sigma'(\gamma) X_\gamma$，则由 $\tilde{T}_u$ 的一个自同构相联系，该自同构把 $X_\gamma$ 送到 $(-1)^{c(\sigma, \sigma') \cdot \gamma} X_\gamma$。

\subsubsection*{穿墙公式}

现在我们准备表述穿墙公式。它的基本构件是群元素
\begin{equation}
\CK_\gamma := \exp \sum_{n=1}^{\infty} \frac{1}{n^2} e_{n \gamma}.
\end{equation}
在我们的等同之下，$\CK_\gamma$ 成为一个作用在 $\tilde{T}_u$ 上的辛同胚，由
\begin{equation} \label{eq:ks-symplect}
\CK_\gamma : X_{\gamma'} \to X_{\gamma'} (1- \sigma(\gamma)
X_\gamma)^{\langle \gamma',\gamma \rangle}
\end{equation}
给出。

对每个电荷为 $\gamma$ 的 BPS 粒子，联系一条复 $\zeta$ 平面中的射线，它由中心荷决定：
\begin{equation}
\ell_\gamma := \{\zeta: Z_\gamma(u) / \zeta \in \IR_- \}.
\end{equation}
当 $u \in \vm$ 变动时，这些射线在 $\zeta$ 平面中转动。射线的循环次序只在 $u$ 到达边缘稳定墙时才改变。在这样一面墙上，一组 BPS 射线 $\ell_\gamma$ 汇聚到一起，它们对应一组使诸 $Z_\gamma$ 变得共线的电荷。在边缘稳定墙上的一般点处，这组电荷可以参数化为\footnote{为确立这些 $\gamma_1$、$\gamma_2$ 的存在性，我们需要用到支撑性质：否则很容易设想诸共线的 $Z_\gamma$ 在原点附近堆积的情形。} $\{n \gamma_1 + m \gamma_2: m, n > 0\}$，其中 $\gamma_1, \gamma_2$ 是某些原始向量，且 $Z_{\gamma_1} / Z_{\gamma_2} \in \IR_+$。

现在，对每个电荷为 $\gamma$ 的 BPS 粒子联系群元素 $\CK_\gamma$，并对在墙上变得共线的诸态取乘积：
\begin{equation} \label{eq:ks-product}
A := \prod^\cwarrow_{\substack{\gamma = n \gamma_1 + m \gamma_2 \\ m>0,\,n>0}}
\CK_\gamma^{\degen(\gamma;u)},
\end{equation}
其中各因子的次序对应于诸射线 $\ell_\gamma$ 的顺时针次序。我们可以对墙两侧的 $u$ 分别考虑这个乘积。当 $u$ 穿过墙时，因子的次序反转，而诸 $\degen(\gamma;u)$ 跳变。穿墙公式的陈述是：整个乘积 $A$ 不变。

这个条件强得足以由诸 $\degen(\gamma;u_-)$ 定出诸 $\degen(\gamma;u_+)$，其中 $u_\pm$ 是在墙的两侧作无穷小位移所得的点。为理解实践中如何做到这一点，我们首先得处理一个重要而细微之处：由于 BPS 态的谱通常是无限的，乘积 \eqref{eq:ks-product} 一般涉及无穷多个因子。按照 \cite{ks1}，我们可以这样理解它。该乘积只涉及诸生成元 $e_{n \gamma_1 + m \gamma_2}$，其中 $m,n>0$。它们生成的李代数可以通过固定某个整数 $L$、并令 $e_{n \gamma_1 + m \gamma_2} = 0$（只要 $n+m > L$）来一致地截断。（也就是说，该李代数被一族满足 $n+m>L$ 的李子代数过滤，而我们可以依次模去 $L$ 越来越大的那些子代数。）经过这样的截断，\eqref{eq:ks-product} 只涉及有限多个非平凡项；无穷乘积则理解为这些截断乘积当 $L \to \infty$ 时的极限。

本着同样的精神，考虑辛同胚 $A$ 的幂级数展开
\begin{equation}
A : X_{\gamma'} \to  ( 1 + \sum_{m>0,n>0} c^{m,n}_{\gamma'} X_{n \gamma_1 + m \gamma_2}) X_{\gamma'},
\end{equation}
并将它截断到 $n+m \le L$。我们可以在边缘稳定墙的一侧算出这一展开，然后在墙的另一侧递归地定出诸 $\degen(\gamma;u)$。具体地说，先取 $L=1$；此时 $\degen(\gamma_1;u)$ 与 $\degen(\gamma_2;u)$ 由「正确地重现 $c^{1,0}_{\gamma'}$ 与 $c^{0,1}_{\gamma'}$」这一要求定出。接着取 $L=2$，考虑 $A \CK_{\gamma_1}^{-\degen(\gamma_1;u)} \CK_{\gamma_2}^{-\degen(\gamma_2;u)}$ 的展开，以提取下一组简并度。这一迭代可以直接地继续下去，定出所有的 $\degen(n \gamma_1 + m \gamma_2 ; u)$。远非显然——但 \cite{ks1} 已作猜想——的是：这样算出的诸 $\degen(\gamma ; u)$ 都是整数！

\subsubsection*{几个例子}

在上述对 Kontsevich--Soibelman 公式的诠释中，我们把他们的广义 Donaldson--Thomas 不变量与物理上定义的 $\degen(\gamma;u)$ 等同了起来。为给这一等同提供动机，下面我们描述几个例子。

如上所述，在边缘稳定墙的一般点上，进入 WCF 的诸辛同胚由一个二维的电荷格生成，即 $\gamma = (p,q) \in \IZ^2$，其典范辛形式为 $\inprod{(p,q), (p',q')} = pq' - q p'$。相应地，我们写 $X_{1,0}=x$，$X_{0,1}=y$。于是辛同胚 $\CK_{p,q}$ 由它们在 $x$ 与 $y$ 上的作用完全确定，具体为
\begin{equation}
\CK_{p,q}: (x,y) \to \biggl(  \bigl(1-(-1)^{pq} x^p y^q \bigr)^q x,
\bigl(1-(-1)^{pq} x^p y^q\bigr)^{-p} y \biggr).
\end{equation}

考虑这样一面边缘稳定墙：一个原始电荷为 $(1,0)$ 的单个 BPS 粒子与一个原始电荷为 $(0,1)$ 的单个粒子，其中心荷在这面墙上汇合。Kontsevich 与 Soibelman 注意到一个漂亮的「五边形恒等式」：
\begin{equation} \label{eq:product-formula-primitive}
\CK_{1,0} \CK_{0,1} = \CK_{0,1} \CK_{1,1} \CK_{1,0}.
\end{equation}
于是 WCF 预言：穿过该墙时只会多产生一个粒子，即由一个荷电粒子与一个荷磁粒子组成的双荷束缚态。确实，来自超引力的「原始穿墙公式」（它在场论中也成立）\cite{Denef:2007vg} 预言：这对粒子会形成一个处于超多重态表示中的单一束缚态。它还预言：一个电荷为 $(1,0)$ 的单个粒子不能与多于一个电荷为 $(0,1)$ 的粒子束缚在一起。要数清由若干个不同类型的粒子组成的更一般的束缚态是相当困难的。它们的不存在，已经是 KS 穿墙公式的一个非平凡预言。

与原始穿墙公式的进一步比较，有助于我们理解 $e_\gamma$ 的对易关系 \eqref{eq:ks-commutator} 中那个符号的作用。考虑乘积 $\CK_{\gamma_1} \CK_{\gamma_2}$，试着把它改写成沿相反方向排列（即让诸 $Z_{\gamma_i}$ 的斜率递增而非递减）的乘积，形如 $\CK_{\gamma_2} \cdots \CK_{\gamma_1}$。设 $\gamma_1,\gamma_2$ 都是原始向量，并考虑由诸 $e_{n\gamma_1 + m \gamma_2}$ 生成的子代数，再模去其中 $n\geq 2, m\geq 2$ 的部分。结果是一个海森堡代数。截断后的海森堡群中的 KS 公式为
\begin{equation}
\CK_{\gamma_1}^{\degen(\gamma_1;u_+)}\CK_{\gamma_1+\gamma_2}^{\degen(\gamma_1+\gamma_2;u_+)}
\CK_{\gamma_2}^{\degen(\gamma_2;u_+)}=
\CK_{\gamma_2}^{\degen(\gamma_2;u_-)}\CK_{\gamma_1+\gamma_2}^{\degen(\gamma_1+\gamma_2;u_-)}
\CK_{\gamma_1}^{\degen(\gamma_1;u_-)}
\end{equation}
其中 $u_\pm$ 是在墙两侧作无穷小位移所得的点。现在，在边缘稳定墙上的一般点处，对 $i=1,2$ 有 $\degen(\gamma_i; u_+) = \degen(\gamma_i; u_-)$。此外，$\CK_{\gamma_1+\gamma_2}$ 在海森堡群中是中心的，因此计算群对易子就重现了原始穿墙公式的如下推论：
\begin{equation}
\Delta \degen = (-1)^{\langle \gamma_1,\gamma_2\rangle -1} \langle
\gamma_1,\gamma_2\rangle \degen(\gamma_1;u) \degen(\gamma_2;u).
\end{equation}
这一论证的一个更精细的版本，可以从 KS 公式提取出 \cite{Denef:2007vg} 的半原始穿墙公式。\footnote{这是与 Wu-yen Chuang 合作的未发表工作中证明的。}

\eqref{eq:product-formula-primitive} 中的例子是特殊的：它的两端都只含有限多项。更典型的情形是遇到无穷乘积。Kontsevich 与 Soibelman 给出的第二个漂亮例子是
\begin{equation}
\CK_{1,0}^2 \CK_{0,1}^2 = \left( \CK_{0,1}^2 \CK_{1,2}^2 \CK_{2,3}^2 \cdots \right)
\CK_{1,1}^{4} \CK_{2,2}^{-2} \left( \cdots \CK_{3,2}^2 \CK_{2,1}^2 \CK_{1,0}^2 \right).
\end{equation}
我们在附录 A 中给出这个恒等式的一个有启发性的证明。% TODO 合稿改 \ref{app:ks-product}
通过一次基变换，我们得到一个物理上很有意思的公式，
\begin{equation} \label{eq:nf2}
\CK_{1,-1}^2 \CK_{0,1}^2 = \left( \CK_{0,1}^2 \CK_{1,1}^2 \CK_{2,1}^2 \cdots \right) \CK_{1,0}^4 \CK_{2,0}^{-2}
\left( \cdots \CK_{3,-1}^2 \CK_{2,-1}^2 \CK_{1,-1}^2 \right),
\end{equation}
它刻画了带两个无质量味的 $SU(2)$ Seiberg--Witten 理论的谱（更确切地说，是超多重态在 $SO(4)=SU(2)_A \times SU(2)_B$ 味对称的矢量表示下变换的情形），如 \cite{Bilal:1996sk} 所述。\footnote{恒等式 \eqref{eq:nf2} 与 Seiberg--Witten 理论的关系，最早由 Frederik Denef 提示；\eqref{eq:nf2} 与 $N_f=2$ 理论的确切关系，则是与 Wu-yen Chuang 合作搞清的。}在右端我们看到完整的弱耦合谱：一个电荷为 $(2,0)$ 的 $W$ 玻色子（它对螺旋度超迹贡献 $-2$）、四个电荷为 $(1,0)$ 的超多重态，以及一组电荷为 $(n,\pm 1)$、重数为 $2$ 的双荷子。（事实上，这些双荷子处于 $SU(2)_A$ 或 $SU(2)_B$ 的双重态中，视 $n$ 的奇偶性而定。）在左端我们看到强耦合谱：一个重数为 $2$ 的单极子（$SU(2)_A$ 的双重态），以及一个重数为 $2$ 的双荷子（$SU(2)_B$ 的双重态）。

变量的一次小小改变 $y \to -y^2$ 把乘积公式 \eqref{eq:nf2} 变成
\begin{equation} \label{eq:nf0}
\CK_{2,-1} \CK_{0,1} = \left( \CK_{0,1} \CK_{2,1} \CK_{4,1} \cdots \right)
\CK_{2,0}^{-2} \left( \cdots \CK_{6,-1} \CK_{4,-1} \CK_{2,-1} \right).
\end{equation}
这个公式刻画了纯 $SU(2)$ Seiberg--Witten 理论的穿墙行为。\footnote{\eqref{eq:nf2} 与 \eqref{eq:nf0} 之所以极为相似，是因为零质量 $N_f=2$ 理论的 SW 曲线是 $N_f=0$ 理论 SW 曲线的二重覆盖。}左端包含强耦合处存在的两个粒子 \cite{Ferrari:1996sv}：一个电荷为 $(0,1)$ 的单极子与一个电荷为 $(2,-1)$ 的双荷子。右端包含弱耦合处的无穷多个双荷子，以及 $W$ 玻色子的贡献 $\CK_{2,0}^{-2}$。

\subsubsection*{加入味信息}

乘积 \eqref{eq:nf2} 描述了 $N_f = 2$ 的 $SU(2)$ Seiberg--Witten 理论的 BPS 谱，但不携带 BPS 粒子的味荷信息。下面我们描述 KS 公式的一个猜想性变体，它把味荷的信息包含进来。（这一公式的物理动机将在第 6 节看到。）% TODO 合稿改 \ref{sec:masses}
引入一个新的格，即味荷格 $\clat^f$，以及一个新参数 $\log \mu \in (\clat^f)^* \otimes_{\IZ} \IC^\times$。然后把诸 $X_\gamma$ 推广为以 $(\gamma,\gamma^f) \in \clat \oplus \clat^f$ 标记的新函数：\footnote{严格说来，完整的电荷局部系统 $\hat\Gamma$ 整体上并不分裂为 $\Gamma \oplus \Gamma^f$；我们实际上有一个扩张 $0 \to \Gamma^f \to \hat\Gamma \to \Gamma \to 0$。不过，这个扩张总可以局部地分裂，而这对我们的目的已经足够。}令 $a$ 取遍 $\clat^f$ 的一组基，
\begin{equation}
X_{\gamma,\gamma^f} := \prod_i (X^i)^{\gamma_i} \prod_a
(\mu^a)^{\gamma^f_a} = X_\gamma \prod_a (\mu^a)^{\gamma^f_a}.
\end{equation}
再定义携带味信息的加细辛同胚：
\begin{equation}
\CK_{\gamma,\gamma^f} : X_{\gamma'} \to X_{\gamma'} (1- \sigma(\gamma)
X_{\gamma,\gamma^f})^{\langle \gamma',\gamma \rangle}.
\end{equation}
中心荷现在依赖于诸质量 $m^a$，即 $Z_{\gamma,\gamma^f}(u) = Z_\gamma(u) + \gamma^f_a m^a$，并确定新的边缘稳定墙。（诸 $\mu^a$ 是诸 $m^a$ 的函数。见下文第 6 节。）% TODO 合稿改 \ref{sec:masses}
我们引入一个与 \eqref{eq:ks-product} 类似的乘积，
\begin{equation} \label{eq:ks-product-masses}
A := \prod^\cwarrow_{\substack{(\gamma,\gamma^f) = n \gamma_1 + m \gamma_2 \\ m>0,\,n>0}}
\CK_{\gamma,\gamma^f}^{\degen(\gamma,\gamma^f;u)},
\end{equation}
推广后的 WCF 断言：$A$ 穿过墙时连续。我们在附录 A 中导出无穷乘积 \eqref{eq:nf2} 的一个含味荷的加细版本；把它与推广后的 WCF 结合起来，就得到 $N_f = 2$ 的 $SU(2)$ 理论正确的穿墙行为。% TODO 合稿改 \ref{app:ks-product}

\subsubsection*{整体性问题}

到目前为止，本节都在 $\vm$ 的一个局部片上工作，并固定了一个二次加细 $\sigma$，以便把 Kontsevich--Soibelman 代数等同于复化环面 $\tilde{T}_u$ 的辛同胚代数——该环面是局部系统 $\tilde{T}$ 的一个纤维。一般而言，不可能在 $\vm$ 上整体地选取这样的加细，因为局部系统 $\clat$ 有单值性。因此，「KS 代数就是作用在 $\tilde{T}$ 上的辛同胚代数」这一说法整体上并不成立。

不过，通过对 $\tilde{T}$ 作适当的扭曲，我们可以定义一个与它密切相关的复化环面丛 $T$，KS 代数确实作用在它上面。$T$ 的定义使得：局部选取一个二次加细就给出一个等同 $T \simeq \tilde{T}$；而给定两个不同的加细 $\sigma, \sigma'$，相应的两个等同相差 $\tilde{T}$ 上的映射 $X_\gamma \to (-1)^{c(\sigma, \sigma') \cdot \gamma} X_\gamma$。纤维方向的辛形式 $\hsymp^{\tilde T}$ 诱导出 $T$ 上相应的纤维方向辛形式 $\hsymp^T$。

具有上述性质的扭曲丛 $T$ 可以这样构造。设 $R$ 记 $\vm$ 上以诸加细 $\sigma$ 为局部截面的那个局部系统。$R$ 是 $\clat^*/2\clat^*$ 的一个挠子（torsor），而 $T$ 是与之相伴的丛：
\begin{equation}
T := \left( \tilde{T} \times R \right) / \left( (X_\gamma, \sigma) \sim
((-1)^{c(\sigma, \sigma') \cdot \gamma} X_\gamma, \sigma') \right).
\end{equation}

\subsection{$\IR^3 \times S^1$ 上的低能有效理论}
\label{sec:reduction}

我们的目标是把 Kontsevich--Soibelman 的 WCF 解释为一个关于 $\IR^3 \times S^1$ 上规范理论的陈述，其中 $S^1$ 的半径为 $R$。我们在一个比其他一切标度都低的能标 $\mu$ 处研究该理论，即 $\mu \ll \Lambda$，且 $\mu \ll 1/R$。在这个能量下，理论看起来有效地是三维的。本节描述它的一些基本性质。

在大半径极限 $R \gg 1/\Lambda$ 下，我们可以用红外拉氏量 \eqref{eq:sw-lagrangian} 确定三维动力学。动力学自由度不过是四维场中不依赖于 $x^4$ 的那些模式。它们当然包括诸标量 $a^I$。此外，从规范扇区我们得到「电」Wilson 线
\begin{equation}
\theta_\elec^I := \oint_{S^1} A_4^I dx^4,
\end{equation}
以及另一组周期标量 $\theta_{\magn,I}$，它们由对偶化 $d=3$ 规范场 $A^I_\alpha dx^\alpha$ 得到。我们常常把它们看作「磁」Wilson 线，
\begin{equation}
\theta_{\magn,I} := \oint_{S^1} (A_{D,4})_I dx^4.
\end{equation}
我们可以用两种方式定义 $\theta_{\magn,I}$：或者在一个把规范场当作自对偶的表述中工作，或者在固定的磁量子数 $P^I$ 下工作、把 $\theta_{\magn,I}$ 作为它们的傅里叶对偶引入。

所有这些周期标量在任一固定的 $u \in \vm$ 处构成一个 $2r$-环面 $\CM_u$ 的坐标。让 $u$ 变动，我们就得到一个环面丛 $\CM$（3d 模空间）。纤维 $\CM_u$ 在 $\vm$ 的奇异轨迹之上退化。$\IR^3$ 上的低能理论是一个以 $\CM$ 为靶空间的 sigma 模型。

更确切地说，$\theta = (\theta_\elec^I, \theta_{\magn,I})$ 是 $2r$-环面的一个局部系统的纤维中的元素，该局部系统记作 $\widetilde\CM := \clat^*\otimes_{\IZ} (\IR/2\pi \IZ)$。对任意 $\gamma \in \clat$，我们得到 $\widetilde\CM$ 上一个角坐标，记作 $\theta_\gamma := \gamma \cdot \theta$。$\widetilde\CM$ 与 $\CM$ 并不完全相同；这里存在一个我们上文略过的整体扭曲，本节末将讨论它。

总之，三维理论是一个映入黎曼流形 $\CM$ 的 sigma 模型；$\CM$ 的实维数为 $4r$，拓扑上是 $\vm$ 上的 $2r$-环面丛。该理论具有 $\CN=4$ 超对称（$8$ 个实超荷），这蕴含 $\CM$ 上的度量是超凯勒的。这个度量是本文研究的主要对象。它此前已在 \cite{Seiberg:1996nz} 中研究过，那里特别把纯 $SU(2)$ 规范理论的 $R \to 0$ 极限识别为 Atiyah--Hitchin 流形。本文更关心相反的极限 $R \to \infty$，因为在这个极限下可以读出 $d=4$ 理论全部 BPS 谱的印记。下一节我们先考虑该极限下的首项行为。

\subsubsection*{整体性问题}

在上面的描述中，我们对 Wilson 线的确切定义稍嫌天真。我们的描述在 $\vm$ 的一个局部片上够用，但正如我们将在第 4 节看到的，它在整体上不可能完全正确。% TODO 合稿改 \ref{sec:periodicnut}
事实上，为使 $\CM$ 上的度量光滑，我们将会看到：绕 $\vm$ 中路径的单值性一般必须伴随 Wilson 线平移 $\pi$。这与我们的天真描述相矛盾，因为环面丛 $\widetilde\CM$ 自带一个零截面。

我们提出，正确的整体图像如下：在任一固定的 $u \in \vm$ 处，Wilson 线位于一个环面 $\CM_u$ 之中，它与 $\widetilde\CM_u$ 同构，但并非典范同构。选定 $\clat_u$ 上二次型 $(-1)^{\inprod{\gamma_1,\gamma_2}}$ 的一个加细 $\sigma$ 后，就得到一个同构 $\CM_u \simeq \widetilde\CM_u$。\footnote{这类二次加细在自对偶规范理论的精确表述中频繁出现 \cite{Witten:1996hc,Freed:2000ta,Hopkins:2002rd,Belov:2006jd}。这里的 $\sigma$ 的来源，很可能可以用这种方式解释。}这样的加细一般只局部存在，所以丛 $\CM$ 与 $\widetilde\CM$ 整体上并不相同。给定两个局部加细 $\sigma$、$\sigma'$，相应的两个局部同构 $\CM \simeq \widetilde\CM$ 相差 $\widetilde\CM$ 上的平移 $\theta \to \theta + \pi c(\sigma,\sigma')$。当然，这一讨论与我们在第 \ref{sec:ReviewKS} 节末所述环面丛 $T$ 与 $\tilde{T}$ 之间的关系是密切平行的。

\subsection{半平坦几何}
\label{sec:semiflat}

在 $R \to \infty$ 极限下，$\CM$ 上度量的首项行为由 $d=3$ 有效作用量支配，后者通过把 \eqref{eq:sw-lagrangian} 简单地截断到不依赖于 $x^4$ 的扇区而得到。这给出
\begin{equation}
\CL^{(3)} = (\im \tau) \left(- \frac{R}{2} \abs{da}^2 - \frac{R}{2}
F^{(3)} \wedge \star F^{(3)} -  \frac{1}{8 \pi^2 R} d\theta_\elec^2
\right) +  (\re \tau) \left(\frac{1}{2\pi} d \theta_\elec \wedge
F^{(3)}\right).
\end{equation}
然后，把 $d=3$ 规范场 $A^I$ 对偶化为标量 $\theta_{\magn,I}$，经过少许整理便得
\begin{equation}
\CL^{(3)}_{dual} = - \frac{R}{2} (\im \tau) \abs{da}^2 - \frac{1}{8 \pi^2 R} (\im \tau)^{-1} \abs{d \theta_\magn - \tau d\theta_\elec}^2.
\end{equation}
这是映入 $\CM$ 的 sigma 模型的拉氏量，其度量局部地由下式给出
\begin{equation} \label{eq:semiflatmetric}
g^\sfm = R (\im \tau) \abs{da}^2 + \frac{1}{4 \pi^2 R} (\im \tau)^{-1} \abs{dz}^2,
\end{equation}
其中我们引入了
\begin{equation}
dz_I = d \theta_{\magn,I} - \tau_{IJ} d \theta_\elec^J.
\end{equation}
（这个记号非常方便，但我们要强调：「$dz_I$」这种形式在整个 $\CM$ 上并不闭合：它只在限制到每个环面纤维 $\CM_u$ 上时才闭合。）

我们称 $g^\sfm$ 为 $\CM$ 上的「半平坦」度量，因为在这个度量下环面纤维是平坦的。表达式 \eqref{eq:semiflatmetric} 反映了这样一个事实：$g^\sfm$ 关于 $\CM$ 上的某个复结构是凯勒的，而在该复结构下 $da^I$ 与 $dz_I$ 构成 $\Omega^{1,0}$ 的一组基。在这个复结构下，$\CM$ 就是 $\vm$ 上由紧\emph{复}环面构成的 Seiberg--Witten 纤维化。（我们将把它与 $\CM$ 上的其他复结构相对照，那些复结构稍后就会遇到；在那些复结构下，诸环面 $\CM_u$ \emph{不是}复子流形。）诸纤维 $\CM_u$ 的体积都是
\begin{equation}
\vol(\CM_u) = \left(\frac{1}{R}\right)^r.
\end{equation}

表达式 \eqref{eq:semiflatmetric} 只在局部有效，因为它用到了一个对偶标架的选择。尽管如此，不同标架下的表达式仍然拼合成一个光滑的度量——除了在 $\vm$ 的奇异轨迹之上，那里 $g^\sfm$ 有奇点。从有效场论的观点看，这样的奇点本不该出现；我们将会看到，它在精确的量子修正度量 $g$ 中由 BPS 瞬子修正所消解。

\begin{yizhu}
$\tilde{T}$ 与 $T$、$\widetilde\CM$ 与 $\CM$ 是两对不同的对象，后文会反复用到，读者务必分清。带波浪号者（$\tilde{T}$、$\widetilde\CM$）是局部系统，只在 $\vm$ 的局部片上才有定义；不带波浪号者（$T$、$\CM$）是整体定义的扭曲丛。二者之间的等同依赖二次加细 $\sigma$ 的局部选取：换一个加细，等同就差一个平移 $X_\gamma \to (-1)^{c(\sigma,\sigma')\cdot\gamma}X_\gamma$（在环面情形为 $\theta \to \theta+\pi c(\sigma,\sigma')$）。本块 §2.2 末与 §2.3 末的两段「整体性问题」讨论的正是这件事。
\end{yizhu}

% ch03 —— §3 超凯勒度量的扭量构造（原文 wcf-paper.tex 1165–1388 行）
\section{超凯勒度量的扭量构造}
\label{sec:twistor}

本节先回顾超凯勒几何的一些一般事实，然后说明我们对 $g$ 的描述背后的基本想法。

\subsection{来自超凯勒流形的全纯数据}
\label{sec:hk-general}

我们首先回顾附着在任一超凯勒流形上的一些全纯数据。按定义，超凯勒流形 $(\CM,g)$ 关于复结构的三元组 $\vec{J}$ 是凯勒的，这些复结构满足关系
\begin{equation} \label{eq:quaternion-complex-structures}
J_1 J_2 = J_3, \quad J_2 J_3 = J_1, \quad J_3 J_1 = J_2, \quad
J_\alpha^2 = -1.
\end{equation}
记 $\omega_\alpha$ 为对应的三个凯勒形式。

事实上，任一超凯勒流形 $(\CM,g)$ 关于更一般的复结构也是凯勒的，即关于 $a^\alpha J_\alpha$ 是凯勒的，其中 $\sum_{\alpha=1}^3 a_\alpha^2 = 1$，对应的凯勒形式为 $a^\alpha \omega_\alpha$。于是我们手上有整整一个 $S^2$ 的复结构。扭量进路的关键洞见之一在于：把这个 $S^2$ 看作黎曼球面、并用一个复参数 $\zeta$ 来标记，是有益的。因此我们把一般的复结构与对应的凯勒形式写成
\begin{align} \label{eq:Jzeta}
J^\pz &= \frac{i(- \zeta + \bar\zeta)J_1 - (\zeta + \bar\zeta) J_2 + (1 - \abs{\zeta}^2) J_3}{1 + \abs{ \zeta}^2},\\
\omega^\pz &= \frac{i(- \zeta + \bar\zeta)\omega_1 -  (\zeta + \bar\zeta) \omega_2 + (1 - \abs{\zeta}^2)\omega _3}{1 + \abs{\zeta}^2}.
\end{align}
我们还将诸凯勒形式组织成第二个组合，
\begin{equation} \label{eq:Omegazeta}
\hsymp(\zeta) = - \frac{i}{2\zeta} \omega_+ + \omega_3 - \frac{i}{2} \zeta \omega_-,
\end{equation}
其中我们引入了记号
\begin{equation}
\omega_\pm = \omega_1 \pm i \omega_2.
\end{equation}
$\hsymp(\zeta)$ 的本质性质是：对任一固定的 $\zeta \in \IC\IP^1$，它是 $\CM$ 上关于复结构 $J^\pz$ 的全纯辛形式。（要对 $\zeta = 0, \infty$ 理解这句话，须分别用 $\zeta$、$1/\zeta$ 去重标度 $\hsymp(\zeta)$。整体地，可以说 $\hsymp(\zeta)$ 被 $\IC\IP^1$ 上的线丛 $\CO(2)$ 扭转。）

\subsection{$g$ 的扭量构造}
\label{sec:twistor-general}

现在我们描述由 $\CM$ 上的全纯数据确定 $g$ 的方法；本文余下部分将一直使用它。

首先我们明确所作假设。回忆 $\CM$ 在拓扑上是 $\vm$ 上的环面丛。对 $\vm$ 中局部片、二次加细（quadratic refinement）、以及电荷格 $\clat$ 的局部截面 $\gamma$ 的每一种选取，我们都假定已给定一个局部定义的 $\IC^\times$ 值函数 $\CX_\gamma(u,\theta;\zeta)$，其自变量为 $(u,\theta) \in \CM$ 与 $\zeta \in \IC^\times$，并具有以下性质：
\begin{itemize}

\item 诸 $\CX_\gamma$ 是乘性的，
\begin{equation}
\CX_\gamma \CX_{\gamma'} = \CX_{\gamma + \gamma'}.
\end{equation}

\item 诸 $\CX_\gamma$ 满足一个实性条件，
\begin{equation} \label{eq:X-real}
\CX_\gamma(\zeta) = \overline{\CX_{-\gamma}(-1 / \bar\zeta)}.
\end{equation}

\item 所有 $\CX_\gamma$ 都是同一组微分方程的解，其形式为
\begin{align}
\frac{\partial}{\partial u^i} \CX &= \left( \frac{1}{\zeta} \CA^{(-1)}_{u^i} + \CA^{(0)}_{u^i} \right)  \CX , \label{eq:cr-1} \\
\frac{\partial}{\partial \bar{u}^{\bar i}} \CX &= \left( \CA^{(0)}_{\bar{u}^{\bar{i}}} + \zeta \CA^{(1)}_{\bar{u}^{\bar{i}}} \right) \CX, \label{eq:cr-2}
\end{align}
其中算子 $\CA^{(n)}_{u^i}$、$\CA^{(n)}_{\bar{u}^{\bar{i}}}$ 是环面纤维 $\CM_u$ 上的复竖直向量场，诸 $\CA^{(-1)}_{u^i}$ 在每一点线性无关，$\CA^{(1)}_{\bar{u}^{\bar{i}}}$ 亦然。（为说明这些方程的动机，注意我们在附录 B 中证明：$(\CM, g)$ 上的柯西--黎曼方程具有这种形式。）% TODO 合稿改 \ref{app:holomorphy}

\item 对每个固定的 $x \in \CM$，$\CX_\gamma(x; \zeta)$ 作为 $\zeta$ 的函数在 $\IC^\times$ 的一个稠密子集上全纯。（在下文的应用中，$\CX_\gamma(x; \zeta)$ 将在可数条直线之并以外全纯。）

\end{itemize}
为陈述关于诸函数 $\CX_\gamma$ 的最后三条假设，我们先定义
\begin{equation} \label{eq:darboux}
\hsymp(\zeta) := \frac{1}{8 \pi^2 R} \epsilon_{ij} \frac{d \CX_{\gamma^i}}{\CX_{\gamma^i}} \wedge \frac{d \CX_{\gamma^j}}{\CX_{\gamma^j}},
\end{equation}
其中 $d$ 指沿纤维的微分，即我们把 $\zeta$ 视为固定参数。我们假设：
\begin{itemize}

\item $\hsymp(\zeta)$ 是整体定义的（特别地，在 $\vm$ 的不同局部片上定义的 $\hsymp(\zeta)$ 彼此相符），并且关于 $\zeta \in \IC^\times$ 全纯。（注意这并不意味着诸 $\CX_\gamma$ 关于 $\zeta$ 全纯；在我们的应用中它们只是分段全纯的。）

\item $\hsymp(\zeta)$ 按全纯辛形式应有的意义非退化，即 $\ker \hsymp(\zeta)$ 是 $4r$ 维空间 $T_\IC \CM$ 中一个维数减半的子空间。

\item $\hsymp(\zeta)$ 当 $\zeta \to 0$ 或 $\zeta \to \infty$ 时只有单极点。

\end{itemize}
在本节余下部分，我们说明如何在 $\CM$ 上定义一个超凯勒度量 $g$，使得诸 $\CX_\gamma(\zeta)$ 是关于复结构 $J^\pz$ 的全纯函数，而 $\hsymp(\zeta)$ 就是 \eqref{eq:Omegazeta} 中那样的全纯辛形式。

我们考虑流形 $\CZ := \CM \times \IC\IP^1$。它具有以下性质：

\begin{enumerate}

\item  \emph{$\CZ$ 是复流形。}
在任一点 $(x, \zeta)$，$2r$ 个方程 \eqref{eq:cr-1}、\eqref{eq:cr-2} 定义了 $T_\IC \CM$ 的一个维数减半的子空间（若 $\zeta = 0$ 或 $\zeta = \infty$，把其中一个方程乘以因子 $\zeta$ 后此事仍成立）。该子空间与 $\partial / \partial \bar{\zeta}$ 所生成子空间的直和，是 $T_\IC \CZ$ 的一个维数减半的子空间。我们把这子空间定义为 $T^{0,1} \CZ$。这\emph{先验地}只定义了 $\CZ$ 上的一个近复结构。然而，诸函数 $\CX_\gamma$ 的存在保证了该近复结构实际上可积。（当然，诸 $\CX_\gamma$ 并非处处关于 $\zeta$ 全纯；但它们在稠密集上全纯，这已足以保证 Nijenhuis 张量消没。特别地，由此可知即使在 $\zeta = 0$ 或 $\zeta = \infty$ 附近，$\CZ$ 上也有复坐标。）

\item  \emph{$\CZ$ 是 $\IC\IP^1$ 上的全纯纤维化。}
投影就是 $p(x, \zeta) = \zeta$。

\item  \emph{存在 $\Omega^2_{\CZ/\IC \IP^1} \otimes \CO(2)$ 的一个全纯截面，它在每条纤维 $p^{-1}(\zeta)$ 上给出全纯辛形式。} 这就是整体定义的 $\hsymp(\zeta)$。

\item  \emph{存在一族全纯截面 $s: \IC\IP^1 \to \CZ$，每个的法丛都是 $N \simeq \CO(1)^{\oplus 2r}$。}
事实上，对每个 $x \in \CM$，可以定义截面 $s_x: \IC \IP^1 \to \CZ$ 为 $s_x(\zeta) = (x, \zeta)$。为看出它是全纯截面，先注意它至少在 $\zeta = 0,\infty$ 之外全纯，理由仅仅是 $\CZ$ 的局部复坐标 $\CX_{\gamma^i}(x, \zeta)$ 在固定 $x$ 时关于 $\zeta$ 全纯；但它连续延拓到 $\zeta = 0,\infty$，故由黎曼可去奇点定理，它在彼处也必全纯。为证明法丛 $N(s_x) \simeq \CO(1)^{\oplus 2r}$，先注意 $N^*(s_x)$ 的全纯截面与 $s_x$ 的一阶无穷小邻域上、在 $s_x$ 上取零的全纯函数之间有一一对应。但这样的函数由其在 $x$ 附近的一阶泰勒展开决定，即它们对应于平凡丛 $p^*((T^*_\IC)_x \CM)$ 的全纯截面，而这些截面零化由方程 \eqref{eq:cr-1}、\eqref{eq:cr-2} 定义的子丛 $B \subset p^*((T_\IC)_x \CM)$。取对偶，得 $N(s_x) \simeq p^*((T_\IC)_x \CM) / B$。另一方面，\eqref{eq:cr-1}、\eqref{eq:cr-2} 给出 $B \otimes \CO(1)$ 的 $2r$ 个平凡化截面。于是我们断定 $N(s_x) \otimes \CO(-1)$ 平凡。

\item  \emph{存在反全纯对合 $\sigma: \CZ \to \CZ$，它覆盖 $\IC \IP^1$ 上的对径映射，并按 $\sigma^* \hsymp = \overline{ \hsymp}$ 的意义保持 $\hsymp$。} 这个对合就是 $\sigma(x, \zeta) = (x, - 1 / \bar\zeta)$。利用实性条件 \eqref{eq:X-real}，可以验证它是反全纯的，且保持 $\hsymp$。
\end{enumerate}

这些正是一个超凯勒流形的\emph{扭量空间}的特征性质，如 \cite{Hitchin:1986ea,MR1206066} 所述。特别地，按 \cite{Hitchin:1986ea,MR1206066} 的配方，可以由 $\CZ$ 重构出 $\CM$ 上的超凯勒度量 $g$。$g$ 可以具体描述如下：注意由 $\hsymp(\zeta)^{r+1} = 0$ 可得 $\omega_+^r \wedge \omega_3 = 0$，后者蕴含实 2-形式 $\omega_3$ 在复结构 $J_3$ 下是 $(1,1)$ 型。因此我们可以用 $J_3$ 与 $\omega_3$ 在 $\CM$ 上造出一个凯勒度量 $g$。这个 $g$ 与扭量构造所保证的超凯勒度量一致。在随后几节中我们将采用这一进路。

\subsection{半平坦几何的扭量构造}
\label{sec:semiflat-twistor}

上述对超凯勒度量的描述，在半平坦度量 $g^\sfm$ 的情形尤为方便；该度量是我们在第 2.4 节引入的。% TODO 合稿改 \ref{sec:semiflat}
与前面一样，我们在 $\vm$ 的一个局部片上工作，并对二次加细作一个局部选取。于是对任一 $\gamma \in \clat$，我们写出局部定义的函数\footnote{这个公式最初是与鲍里斯·皮奥利纳（Boris Pioline）合作得到的；它本质上也是 \cite{Neitzke:2007ke} 中四元数凯勒情形的公式的刚性极限。它为发现本文所描述的构造提供了重要线索。}
\begin{equation} \label{eq:X-sf}
\CX_\gamma^\sfm(\zeta) := \exp\biggl[ \pi  R \zeta^{-1}  Z_\gamma + i
\theta_\gamma + \pi R \zeta \bar Z_\gamma \biggr].
\end{equation}
这些函数满足形如 \eqref{eq:cr-1}、\eqref{eq:cr-2} 的「柯西--黎曼方程」，其中
\begin{gather}
\CA^{(-1)}_{u^i} = -i \pi R \frac{\partial Z}{\partial u^i} \cdot \frac{\partial}{\partial \theta}, \quad \CA^{(1)}_{\bar{u}^{\bar{i}}} = -i \pi R \frac{\partial \bar{Z}}{\partial \bar{u}^{\bar{i}}} \cdot \frac{\partial}{\partial \theta}, \\
\CA^{(0)}_{u^i} = 0, \quad \CA^{(0)}_{\bar{u}^{\bar{i}}} = 0
\end{gather}
而 $Z$ 代表周期向量。于是 $\hsymp^\sfm(\zeta)$ 为
\begin{align} \label{eq:Omega-sf}
\hsymp^\sfm(\zeta) & := \frac{1}{8 \pi^2 R} \epsilon_{ij} \frac{d \CX^\sfm_{\gamma^i}}{\CX^\sfm_{\gamma^i}} \wedge \frac{d \CX^\sfm_{\gamma^j}}{\CX^\sfm_{\gamma^j}} \\
&= \frac{1}{4\pi} \left[ \frac{i}{\zeta} \left\langle d Z, d \theta \right\rangle + \left( \pi R \langle d Z, d \bar{Z} \rangle - \frac{1}{2 \pi R} \langle d \theta , d \theta \rangle \right) + i \zeta \langle d \bar{Z}, d \theta \rangle \right]. \label{eq:Omega-sf-2}
\end{align}
（注意：消没条件式 (2.5) 保证 $\hsymp^\sfm(\zeta)$ 不含 $\zeta^{-2}$ 阶与 $\zeta^2$ 阶的项。）% TODO 合稿改 \eqref{eq:lag-constraint}
$\hsymp^\sfm(\zeta)$ 与 $\CX^\sfm_\gamma(\zeta)$ 满足我们在第 \ref{sec:twistor-general} 节所述构造所需的诸条件，因此它们就是 $\CM$ 上某个超凯勒度量的全纯辛形式与复坐标。正如我们下面要验证的，这个度量正是所求的 $g^\sfm$。

先注意，比较 \eqref{eq:Omegazeta} 与 \eqref{eq:Omega-sf-2} 中的领头项，给出
\begin{equation} \label{eq:omegaplus-sf}
\omega_+^\sfm = - \frac{1}{2\pi} \inprod{d Z, d \theta}.
\end{equation}
由 $\omega_+^\sfm$ 可以确定复结构 $J^\sfm_3$：事实上，选定一个电磁对偶标架后，我们可以把 \eqref{eq:omegaplus-sf} 改写为
\begin{equation}
\omega_+^\sfm = \frac{1}{2\pi} da^I \wedge dz_I.
\end{equation}
这就清楚地表明，$\CM$ 在复结构 $J^\sfm_3$ 之下正是由复环面构成的 \SW 纤维化。这正是我们在第 2.4 节已经描述过的复结构。% TODO 合稿改 \ref{sec:semiflat}

类似地，比较 \eqref{eq:Omegazeta} 与 \eqref{eq:Omega-sf-2} 中与 $\zeta$ 无关的项，给出
\begin{equation} \label{eq:omega3-sf}
\omega_3^\sfm = \frac{R}{4} \langle d Z , d \bar{Z} \rangle - \frac{1}{8 \pi^2 R} \langle d\theta, d \theta \rangle,
\end{equation}
它可改写为
\begin{equation}
\omega_3^\sfm = \frac{i}{2} \left( R (\im \tau)_{IJ} d a^I \wedge d \bar{a}^J + \frac{1}{4 \pi^2 R} ((\im \tau)^{-1})^{IJ} dz_I \wedge d\bar{z}_J \right).
\end{equation}
将此式与式 (2.35) 比较可知，$g^\sfm$ 确实关于复结构 $J^\sfm_3$ 是凯勒的、且凯勒形式为 $\omega_3^\sfm$，因而是由 $\hsymp^\sfm(\zeta)$ 出发的扭量构造所保证的那个超凯勒度量。% TODO 合稿改 \eqref{eq:semiflatmetric}

本节我们看到，$\CM$ 上的半平坦度量及其超凯勒结构，都可以由 \eqref{eq:X-sf} 中定义的诸函数 $\CX_\gamma^\sfm$ 构造出来。这些函数对下文具有根本的重要性。

\begin{yizhu}
$\hsymp(\zeta)$ 在原文中出现两次定义：式~\eqref{eq:Omegazeta} 用它指超凯勒结构给出的全纯辛形式，式~\eqref{eq:darboux} 又以「$:=$」把它定义为由诸 $\CX_\gamma$ 组装出的量。两者是同一个对象：本节的构造正是要求由式~\eqref{eq:darboux} 得到的 $\hsymp(\zeta)$ 恰好等于式~\eqref{eq:Omegazeta} 中的那一个（见本节末的陈述）。
\end{yizhu}

% ch04 —— §4 相互局域修正（原文 wcf-paper.tex 1389–2128 行，对应 PDF 第 22–34 页）
% 原文宏 \plog、\eps 不在本工程 preamble 中，此处本地补齐（定义与原文 wcf-paper.tex 159、166 行相同）。
% 原文 (4.42) 中一处 \biggr) 疑为 \right) 之笔误（原文 1720 行），转录保持原文，编译无碍。
\providecommand{\plog}[1]{#1\partial_{#1}}
\providecommand{\eps}{\epsilon}

\section{相互局域修正}
\label{sec:periodicnut}

如果我们只考虑无质量部分的朴素量纲约化，那么半平坦度量 $g^\sfm$ 就是全部答案了。然而，$d=4$ 的理论中还有有质量的 BPS 粒子。当这些有质量粒子中的一个或多个绕 $S^1$ 转一圈时，相应的「瞬子」组态会给度量带来修正。由于 $d=4$ 理论中状态的能量满足 $M \ge \abs{Z}$，这些修正的权重至少是 $e^{-2 \pi R \abs{Z}}$。

本节所研究的例子中，所有 BPS 粒子都相互局域。这比一般情形容易处理得多，因为我们可以选一个对偶标架，使这些粒子全都带电荷，于是完全可以只在有效拉格朗日描述之内工作。

本节大部分篇幅里，我们进一步特化到耦合单个带电荷超多重态的自由 $U(1)$ 规范理论。这个理论既是最简单的例子，也在物理上切题：它描述的是 $\vm$ 中一般奇点附近的物理——在那里，有一个 BPS 粒子变得比其余粒子轻得多。

\subsection{精确的单粒子度量}

我们考虑 $\IR^3 \times S^1$ 上耦合单个电荷为 $q>0$ 的超多重态（连同其电荷为 $-q$ 的 CPT 共轭）的 $U(1)$ 规范理论。我们将要描述的度量，此前已在 \cite{Ooguri:1996me,Seiberg:1996ns} 中考虑过。

$d=4$ 理论的模空间 $\vm$ 由矢量多重态标量 $a \in \IC$ 坐标化。更确切地说，$\vm$ 只是 $\IC$ 中的一个开片，因为 $d=4$ 理论不是渐近自由的：在 $\abs{a} \sim \abs{\Lambda}$ 处存在截断。

正如我们在第 2.3 节中所解释的，$d=3$ 理论的模空间 $\CM$ 是 $\vm$ 上的 2-环面丛。环面纤维 $\CM_a$ 由电 Wilson 线 $\theta_\elec$ 与磁 Wilson 线 $\theta_\magn$ 坐标化（这里暂时忽略二次加细（quadratic refinement）的微妙之处），两者的周期都是 $2 \pi$。% TODO 合稿改 \ref{sec:reduction}

半平坦度量 $g^\sfm$ 具有 $U(1)^2$ 的等距作用，因为 $\theta_\elec$ 与 $\theta_\magn$ 的平移都是精确对称性。带电荷的超多重态与 $\theta_\elec$ 耦合，因而破坏了平移 $\theta_\elec$ 的那个等距。不过理论中没有带磁荷的 BPS 态，所以 $\theta_\magn$ 的平移仍是精确等距。因此修正后的度量 $g$ 具有 Gibbons–Hawking 形式。

为与 Gibbons–Hawking 拟设作比较，我们引入矢量 $\vec{x}$：
\begin{equation}
a = x^1 + i x^2, \quad \theta_\elec = 2 \pi R x^3.
\end{equation}
$\theta_\magn$ 是由 $\vec{x}$ 参数化的 $\IR^2 \times S^1$ 开子集上某个 $U(1)$ 丛的局部坐标。度量为
\begin{equation}
g = V(\vec{x})^{-1} \left(\frac{d \theta_\magn}{2\pi} + A(\vec{x})\right)^2 + V(\vec{x}) d\vec{x}^2,
\end{equation}
其中 $V$ 是正的调和函数（下文将算出它），$A$ 是曲率为
\begin{equation} \label{eq:FdV}
F = \star dV.
\end{equation}
的 $U(1)$ 联络。这是标准 Gibbons–Hawking 拟设的一个轻微推广：标准拟设中取 $\vec{x}$ 落在 $\IR^3$（的某个开子集）里。（我们可以先在 $\IR^3$ 的合适子集上工作，再对总空间上平移 $x^3$ 的 $\IZ$-作用取商。）在标准拟设中，所有满足 \eqref{eq:FdV} 的 $A$ 都互相规范等价，因而定义同一个度量。在我们的情形中这并不完全成立：与绕 $S^1$ 的 holonomy 相关，还有一个额外的规范不变自由度。这一选择与 $d=4$ 中 $\theta$ 角的选择有关。

在我们的情形中，$V(\vec{x})$ 可以通过一圈计算积分掉带电荷超多重态来得到。文献 \cite{Seiberg:1996ns} 断言了一个非重整化定理，它意味着这一计算是精确的。所得的 $V$ 是 $\IR^2 \times S^1$ 上带 $q$ 个奇点的调和函数。$\theta_\elec$ 的周期性来自对带电荷超多重态在 $S^1$ 上的 Kaluza–Klein 动量求和：
\begin{equation} \label{eq:V-sum}
V = \frac{q^2 R}{4 \pi} \sum_{n=-\infty}^\infty \left( \frac{1}{\sqrt{q^2 R^2 \abs{a}^2 + (q \frac{\theta_\elec}{2 \pi} + n)^2}} - \kappa_n \right)
\end{equation}
这里 $\kappa_n$ 是为使求和收敛而引入的正规化常数。\eqref{eq:V-sum} 的 Poisson 重求和表明
\begin{equation}
V = V^\sfm + V^\inst,
\end{equation}
其中
\begin{align}
V^\sfm &= -\frac{q^2 R}{4\pi} \left(\log  \frac{a}{\Lambda} + \log \frac{\bar{a}}{\bar{\Lambda}} \right), \\
V^\inst &= \frac{q^2 R}{2\pi} \sum_{n\neq 0} e^{i n q \theta_\elec} K_0(2 \pi R \abs{n q a}). \label{eq:V-inst}
\end{align}
这里 $\Lambda$ 是与 $\kappa_n$ 的选取有关的紫外截断。\footnote{例如，若取 $\kappa_n = (\vert \tilde \Lambda\vert^2 + n^2)^{-1/2}$，则可取 $\Lambda = (qR)^{-1} \tilde\Lambda \exp[-2 \sum_{m=1}^\infty K_0(2\pi m \vert \tilde \Lambda\vert )]$。}

要完全指定度量，还须给出满足 \eqref{eq:FdV} 的 $A(\vec{x})$：
\begin{equation}
A = A^\sfm + A^\inst,
\end{equation}
其中
\begin{align}
A^\sfm &= \frac{i q^2}{8 \pi^2} \left( \log \frac{a}{\Lambda} - \log \frac{\bar{a}}{\bar{\Lambda}} \right) d\theta_\elec, \label{eq:A0} \\
A^\inst &=  - \frac{q^2 R}{4 \pi} \left( \frac{da}{a} - \frac{d \bar a}{\bar a} \right) \sum_{n \neq 0} (\sgn n) e^{i n q \theta_\elec} \abs{a} K_1(2 \pi R \abs{n q a}).  \label{eq:A-inst}
\end{align}
在大 $R$ 时，$V$ 与 $A$ 中的领头项是 $V^\sfm$ 与 $A^\sfm$。只保留这些项，$g$ 就成为半平坦度量，其
\begin{equation}
\tau = \frac{q^2}{2 \pi i} \log \frac{a}{\Lambda}.
\end{equation}
这就是 $d=4$ 中积分掉超多重态所得的跑动耦合。

次领头项 $V^\inst$、$A^\inst$ 给出对半平坦度量的修正。由于 $x \to +\infty$ 时的渐近行为 $K_\nu(x) \sim \sqrt{\frac{\pi}{2x}} e^{-x}$，它们如我们所预期地呈瞬子展开的形式。它们也如预期地破坏了沿 $\theta_\elec$ 的平移不变性。最后，它们改善了奇性。回忆在 $g^\sfm$ 中 $a=0$ 处有一个奇点。由 \eqref{eq:V-sum} 可见，$g$ 唯一可能的奇点出现在 $a = 0$、$q \theta_\elec = 2 \pi n$ 处。研究这些点附近的度量可知，每一处都是 $A_{q-1}$ 型锥奇点。于是 $g^\sfm$ 中的一个奇点被 $g$ 中的 $q$ 个更高余维的奇点取代。在最简单的 $q=1$ 情形，奇性被完全光滑化。

\subsubsection*{整体性问题}

关于坐标 $\theta_\magn$ 的整体定义，有一个微妙的问题。我们选用的规范便于讨论 $\theta_\elec$ 的周期性。然而 $A^\sfm$ 中出现的对数表明，这个规范在 $a = 0$ 处是奇异的。此外，沿绕原点一周的延拓 $a \to e^{2 \pi i} a$，$A^\sfm$ 平移 $- \frac{q^2}{2 \pi} d \theta_\elec$。这一平移必须由一个规范变换来补偿：
\begin{equation}
\theta_\magn \to \theta_\magn + q^2 \theta_\elec + C.
\end{equation}
为确定 $C$，我们作一个规范变换，换到新坐标 $\theta_\magn'$：
\begin{equation}
\theta_\magn' = \theta_\magn + \frac{i}{4 \pi} \left( \log \frac{a}{\Lambda} - \log \frac{\bar a}{\bar\Lambda} \right) (q^2 \theta_\elec + C).
\end{equation}
当 $a$ 绕原点一周时，变换后的 $\theta_\magn'$ 是单值的。规范变换后的 $A^\sfm$ 是
\begin{equation}
(A')^\sfm = - \frac{i}{4 \pi} \left( \frac{da}{a} - \frac{d \bar a}{\bar a} \right) (q^2 \theta_\elec + C).
\end{equation}
现在我们关注 $q \theta_\elec = \pi$ 处的行为。这里 $A^\inst = 0$，所以精确的规范场恰由 $(A')^\sfm$ 给出。另一方面，一旦纳入瞬子修正，该点上度量与 $U(1)$ 丛都不再奇异（回忆唯一的奇点出现在 $a=0$、$q \theta_\elec = 2 \pi n$ 处）。又因 $\theta_\magn'$ 是单值的，于是 $(A')^\sfm$ 在这里不可能有奇性——即使延伸到 $a=0$ 也不行（更确切地说，唯一允许的奇性是一条量子化的 Dirac 弦）——这意味着
\begin{equation}
C = - q \pi + 2 \pi  k
\end{equation}
对某整数 $k$ 成立。于是我们的结论是：绕 $a = 0$ 一周时，角坐标的平移为
\begin{subequations}
\begin{align} \label{eq:angle-monodromy}
\theta_\elec &\to \theta_\elec, \\
\theta_\magn &\to \theta_\magn + q^2 \theta_\elec - q \pi. \label{angle-monodromy-varphi}
\end{align}
\end{subequations}
$q^2 \theta_\elec$ 的平移正如环面丛的单值性所预期。$-q \pi$ 的平移则更令人意外，但它正契合我们在第 2.3 节末的讨论：我们那时提出，Wilson 线只有在选定一个局部的二次加细 $\sigma$ 之后才有良好定义。本节至此一直选用「标准的」加细 $\sigma(\gamma_e,\gamma_m) = (-1)^{\gamma_e \gamma_m}$。单值性把 $\gamma_e \to \gamma_e + q^2 \gamma_m$，从而把 $\sigma$ 换成 $\sigma'(\gamma_e,\gamma_m) = (-1)^{q^2 \gamma_m^2} \sigma(\gamma_e,\gamma_m)$。这一加细的改变由 $\theta_\magn$ 平移 $-q \pi$ 来补偿。% TODO 合稿改 \ref{sec:reduction}

\subsection{超凯勒结构}

接下来我们要把 $\CM$ 描述为超凯勒流形。任何 Gibbons–Hawking 度量的超凯勒结构，都由如下三重辛形式决定：
\begin{equation}
\omega^\alpha = d x^\alpha \wedge \left(\frac{d \theta_\magn}{2\pi} + A(\vec{x})\right) + \half \eps^{\alpha \beta \gamma} V dx^\beta \wedge dx^\gamma.
\end{equation}
全纯辛形式 式 (3.3) 则为
\begin{equation}
\hsymp(\zeta) = \frac{1}{4 \pi^2 R} \xi_\magn \wedge \xi_\elec
\end{equation}
其中
\begin{align}
\xi_\magn &= i d\theta_\magn + 2 \pi i A(\vec{x}) + \pi i V \left(\frac{1}{\zeta} da - \zeta d\bar{a}\right), \label{eq:xi-magn} \\
\xi_\elec &= i d\theta_\elec + \pi R \left(\frac{1}{\zeta} da + \zeta d\bar{a}\right). \label{eq:xi-elec}
\end{align}
特别地，由此可知 $\xi_\elec$ 与 $\xi_\magn$ 都是 $(1,0)$ 型的。% TODO 合稿改 \eqref{eq:Omegazeta}

此外，$\xi_\elec$ 可以写成
\begin{equation}
\xi_\elec = \frac{d\CX_\elec}{\CX_\elec}
\end{equation}
其中
\begin{equation} \label{eq:X1-def}
\CX_\elec = \exp \left[ \pi R \frac{a}{\zeta} + i \theta_\elec + \pi R \zeta \bar{a} \right].
\end{equation}
所以 $\CX_\elec$ 是 $\CM$ 上复结构 $J^\pz$ 下的全纯函数。注意：若取 $\gamma$ 为单位电荷，则它与 式 (3.10) 给出的半平坦坐标 $\CX^\sfm_{\gamma}$ 一致，因为此时 $Z_\gamma = a$、$\theta_\gamma = \theta_\elec$。换言之，「电」复坐标不受带电荷粒子的瞬子修正影响：
\begin{equation}
\CX_\elec = \CX^\sfm_\elec.
\end{equation}
% TODO 合稿改 \eqref{eq:X-sf}

为完成 $\CM$ 复几何的描述，还应构造第二个「磁」复坐标 $\CX_\magn$，使得
\begin{equation} \label{eq:second-coord}
\hsymp(\zeta) =- \frac{1}{4 \pi^2 R}\frac{d \CX_\elec}{\CX_\elec} \wedge \frac{d \CX_\magn}{\CX_\magn}.
\end{equation}
这样的 $\CX_\magn$ 必形如
\begin{equation} \label{eq:mag-form}
\CX_\magn = e^{i \theta_\magn + \Phi(a, \bar{a},\theta_\elec, \zeta)}.
\end{equation}

构造 $\CX_\magn$ 最显然的做法，是写出 $\CM$ 上的 Cauchy–Riemann 方程，然后寻找形如 \eqref{eq:mag-form} 的特解。下一节我们改用另一种做法：直接给出 $\CX_\magn$ 的一个特解，其形式对后文特别方便；然后不去验证 Cauchy–Riemann 方程，而是直接验证 \eqref{eq:second-coord}。

\subsection{$\CX_\magn$ 的解}
\label{sec:solving-xmagn}

现在特化到我们的 $\CM$。此时
\begin{equation}
\hsymp(\zeta) = \hsymp^\sfm(\zeta) + \hsymp^\inst(\zeta)
\end{equation}
其中
\begin{align}
\hsymp^\sfm(\zeta) &= - \frac{1}{4 \pi^2 R} \xi_\elec \wedge \left[ i d \theta_\magn + 2 \pi i A^\sfm + \pi i V^\sfm \left(\frac{1}{ \zeta} da -  \zeta d\bar{a}\right) \right],  \\
\hsymp^\inst(\zeta) &= - \frac{1}{4 \pi^2 R} \xi_\elec \wedge \left[ 2 \pi i A^\inst +
\pi i V^\inst \left(\frac{1}{\zeta} da -  \zeta d\bar{a}\right)\right]. \label{eq:instanton-corrections-Omega}
\end{align}
若忽略瞬子修正，所需的磁坐标为
\begin{equation} \CX_\magn^\sfm(\zeta)= \exp\left[-i \frac{R q^2}{2 \zeta} \left(a
\log\frac{a}{\Lambda} -a\right) + i \theta_\magn + i \frac{\zeta R q^2}{2} \left(\bar a
\log\frac{\bar a}{\bar \Lambda} - \bar a\right)\right].
\end{equation}
若取 $\gamma$ 为单位磁荷、$Z_\gamma = \frac{q^2}{2 \pi i} (a \log \frac{a}{\Lambda} - a)$、$\theta_\gamma = \theta_\magn$，它便与半平坦几何中全纯坐标 $\CX^\sfm_\gamma$ 的表达式 式 (3.10) 一致。直接计算可以验证
\begin{equation}
\frac{d \CX^\sfm_\magn}{\CX^\sfm_\magn} = \left[ i d \theta_\magn + 2 \pi i A^\sfm + \pi i V^\sfm \left(\frac{1}{\zeta} da -  \zeta d\bar{a}\right) \right] - \frac{i q^2}{4 \pi}\left( \log \frac{a}{\Lambda} - \log \frac{\bar{a}}{\bar{\Lambda}} \right) \frac{d \CX_\elec}{\CX_\elec},
\end{equation}
从而特别地
\begin{equation} \label{eq:Omega-sf-repeat}
\hsymp^\sfm(\zeta) = - \frac{1}{4 \pi^2 R} \frac{d\CX_\elec}{\CX_\elec} \wedge \frac{d\CX_\magn^\sfm}{\CX_\magn^\sfm},
\end{equation}
如所预期。% TODO 合稿改 \eqref{eq:X-sf}

注意 $\CX_\magn^\sfm$ 绕 $a=0$ 有非平凡的单值性：$\log a$ 与 $\log \bar a$ 的单值性，与 \eqref{angle-monodromy-varphi} 中给出的 $e^{i \theta_\magn}$ 的单值性合起来，给出
\begin{equation}\label{eq:mono2}
\CX_\magn^\sfm \to (-1)^q \CX_\elec^{q^2} \CX_\magn^\sfm.
\end{equation}

接下来纳入瞬子修正。如下文将要证明的，满足 \eqref{eq:second-coord} 的所需 $\CX_\magn$ 可由如下积分公式给出：
\begin{equation}\label{eq:FullX2}
\begin{split}
\CX_\magn = \CX_\magn^\sfm \exp\Biggl[ \frac{i q}{4\pi} & \int_{\ell_+} \frac{d \zeta'}{\zeta'}
\frac{\zeta' + \zeta}{\zeta' - \zeta} \log[1-\CX_\elec(\zeta')^q]  \\
- \frac{i q}{4\pi} & \int_{\ell_-}
 \frac{d \zeta'}{\zeta'} \frac{\zeta' + \zeta}{\zeta' - \zeta} \log[1-\CX_\elec(\zeta')^{-q}] \Biggr],
\end{split}
\end{equation}
其中围道 $\ell_\pm$ 取为位于两个半平面
\begin{equation}
\CU_\pm = \left\{ \zeta: \pm {\rm Re} \frac{a}{\zeta}<0 \right\}.
\end{equation}
内、连接 $0$ 与 $\infty$ 的任意路径。\eqref{eq:FullX2} 中的两个积分贡献分别来自绕 $S^1$ 正绕数与负绕数的瞬子修正。

本节余下的部分验证 \eqref{eq:FullX2} 确实正确。这归结为验证
\begin{equation} \label{eq:tomatch}
- \frac{1}{4 \pi^2 R}\frac{d \CX_\elec}{\CX_\elec} \wedge \frac{d \CX_\magn}{\CX_\magn} = \hsymp(\zeta).
\end{equation}
由 \eqref{eq:FullX2} 有
\begin{equation} \label{eq:dx2x2}
\frac{d \CX_\magn}{\CX_\magn} = \frac{d \CX^\sfm_\magn}{\CX^\sfm_\magn} + \CI_+ + \CI_-
\end{equation}
其中
\begin{equation} \label{eq:ipm}
\CI_\pm = -\frac{i q^2}{4\pi} \int_{\ell_\pm} \frac{d \zeta'}{\zeta'}
\frac{\zeta' + \zeta}{\zeta' - \zeta} \left[ \frac{\CX_\elec(\zeta')^{\pm q}}{1-\CX_\elec(\zeta')^{\pm q}} \frac{d\CX_\elec(\zeta')}{\CX_\elec(\zeta')}\right].
\end{equation}
（这里用到了这样的事实：\eqref{eq:FullX2} 中的积分只通过 $\CX_\elec(\zeta')$ 依赖于 $a,\bar a,\theta_\magn,\theta_\elec$，并且绝对收敛，因此我们可以自由地把微分 $d$ 移入积分号内。）综合 \eqref{eq:tomatch}、\eqref{eq:dx2x2} 与 \eqref{eq:Omega-sf-repeat} 可见，积分 $\CI_\pm$ 必须给出 \eqref{eq:tomatch} 右端的瞬子部分 $\hsymp^\inst(\zeta)$，也就是说我们需要
\begin{equation} \label{eq:intermediate}
\frac{d \CX_\elec(\zeta)}{\CX_\elec(\zeta)} \wedge (\CI_+ + \CI_-) = \frac{d \CX_\elec(\zeta)}{\CX_\elec(\zeta)} \wedge \left[ 2 \pi i A^\inst + \pi i V^\inst \left(\frac{1}{\zeta} da - \zeta d\bar{a}\right)\right].
\end{equation}
为验证这一点，先注意
\begin{equation}
\frac{d \CX_\elec(\zeta)}{\CX_\elec(\zeta)} \wedge \CI_\pm = -\frac{i q^2}{4 \pi} \int_{\ell_\pm} \frac{d \zeta'}{\zeta'}
 \left( \frac{\zeta' + \zeta}{\zeta' - \zeta} \frac{d \CX_\elec(\zeta)}{\CX_\elec(\zeta)} \wedge \frac{d\CX_\elec(\zeta')}{\CX_\elec(\zeta')} \right) \left[ \frac{\CX_\elec(\zeta')^{\pm q}}{1-\CX_\elec(\zeta')^{\pm q}} \right]
\end{equation}
而这里出现的二形式可以改写为
\begin{align}\label{eq:rearrange}
\frac{\zeta' + \zeta}{\zeta' - \zeta} \frac{d \CX_\elec(\zeta)}{\CX_\elec(\zeta)} \wedge \frac{d \CX_\elec(\zeta')}{\CX_\elec(\zeta')} &=
\frac{\zeta' + \zeta}{\zeta' - \zeta} \frac{d \CX_\elec(\zeta)}{\CX_\elec(\zeta)} \wedge \left[\frac{d \CX_\elec(\zeta')}{\CX_\elec(\zeta')}-\frac{d \CX_\elec(\zeta)}{\CX_\elec(\zeta)}\right]\\
&=  - \pi R \frac{d \CX_\elec(\zeta)}{\CX_\elec(\zeta)} \wedge \left[\left(\frac{1}{\zeta'} + \frac{1}{\zeta}\right) da - (\zeta'+\zeta)d\bar a\right],
\end{align}
这里用到了 $\CX_\elec$ 的显式形式 \eqref{eq:X1-def}。于是 \eqref{eq:intermediate} 的左端成为
\begin{equation} \label{eq:intermediate-2}
\begin{split}
\frac{i q^2 R}{4} \frac{d \CX_\elec(\zeta)}{\CX_\elec(\zeta)} \wedge
\biggl( & \int_{\ell_+}  \frac{d \zeta'}{\zeta'}  \left[\left(\frac{1}{\zeta'}+\frac{1}{\zeta}\right)da - ( \zeta'+\zeta)d\bar a\right]\frac{\CX_\elec(\zeta')^{q}}{1-\CX_\elec(\zeta')^{q}}  \\
+ & \int_{\ell_-}  \frac{d \zeta'}{\zeta'}  \left[\left(\frac{1}{\zeta'}+\frac{1}{\zeta}\biggr)da - ( \zeta'+\zeta)d\bar a\right]\frac{\CX_\elec(\zeta')^{-q}}{1-\CX_\elec(\zeta')^{-q}}
\right).
\end{split}
\end{equation}

现在可以计算这些积分了。方便起见，先把两个围道 $\ell_\pm$ 分别形变到恰好位于 $\CU_\pm$ 正中的标准位置，即取
\begin{equation}
\ell_\pm = \left\{ \zeta: \pm \frac{a}{\zeta} \in \IR_- \right\}.
\end{equation}
先考虑乘 $\zeta$ 或 $\frac{1}{\zeta}$ 的那些项。展开几何级数可得：
\begin{align}
\int_{\ell_+}  \frac{d \zeta'}{\zeta'} \frac{\CX_\elec(\zeta')^{q}}{1-\CX_\elec(\zeta')^{q}} &=
\sum_{n>0} \int_{\ell_+} \frac{d \zeta'}{\zeta'}
\exp \left[ \pi R q n\frac{a}{\zeta'} + i q n\theta_\elec + \pi R q n \zeta' \bar{a} \right] \\
&= \sum_{n>0} 2 e^{i q n\theta_\elec}K_0(2 \pi R q|na|).
\end{align}
\eqref{eq:intermediate-2} 中 $\ell_-$ 上的积分给出 $n<0$ 的类似求和。合计得到：\eqref{eq:intermediate-2} 中乘 $\zeta$ 或 $\frac{1}{\zeta}$ 的项等于
\begin{equation} \label{eq:final-part-1}
\frac{d\CX_\elec}{\CX_\elec} \wedge \left( \frac{i q^2 R}{2} \sum_{n \neq 0} e^{i q n\theta_\elec}K_0(2 \pi R q|na|) \right) \left(\frac{1}{\zeta} da - \zeta d\bar{a}\right) = \frac{d\CX_\elec}{\CX_\elec} \wedge i \pi V^\inst\left(\frac{1}{\zeta} da - \zeta d\bar{a}\right).
\end{equation}
对 \eqref{eq:intermediate-2} 中其余的项，类似地有
\begin{align}
\int_{\ell_+}  \frac{d \zeta'}{\zeta'} \zeta' \frac{\CX_\elec(\zeta')^{q}}{1-\CX_\elec(\zeta')^{q}} &=
\sum_{n>0} \int_{\ell_+} \frac{d \zeta'}{\zeta'}
\zeta' \exp \left[ \pi R q n\frac{a}{\zeta'} + i q n\theta_\elec + \pi R q n \zeta' \bar{a} \right] \\
&= - \sum_{n>0} 2 \frac{\abs{a}}{\bar{a}} e^{i q n\theta_\elec}K_1(2 \pi R q|na|)
\end{align}
以及
\begin{align}
\int_{\ell_+}  \frac{d \zeta'}{\zeta'} \frac{1}{\zeta'} \frac{\CX_\elec(\zeta')^{q}}{1-\CX_\elec(\zeta')^{q}} &=
\sum_{n>0} \int_{\ell_+} \frac{d \zeta'}{\zeta'}
\frac{1}{\zeta'} \exp \left[ \pi R q n\frac{a}{\zeta'} + i q n\theta_\elec + \pi R q n \zeta' \bar{a} \right] \\
&= - \sum_{n>0} 2 \frac{\abs{a}}{a} e^{i q n\theta_\elec}K_{1}(2 \pi R q|na|).
\end{align}
把这些与来自 $\ell_-$ 上积分的对应项（它们多带一个负号）合并，我们看到 \eqref{eq:intermediate-2} 中的这些项等于
\begin{equation} \label{eq:final-part-2}
\frac{d\CX_\elec}{\CX_\elec} \wedge \left( - \frac{i q^2 R}{2} \sum_{n \neq 0} e^{i q n\theta_\elec} (\sgn n) \abs{a} K_1(2 \pi R q|na|) \right) \left( \frac{da}{a} - \frac{d \bar a}{\bar a} \right) = \frac{d\CX_\elec}{\CX_\elec} \wedge 2 \pi i A^\inst.
\end{equation}
于是最后，把 \eqref{eq:final-part-1} 与 \eqref{eq:final-part-2} 相加，便如所愿得到 \eqref{eq:intermediate}：对 $\CX_\magn$ 中的围道积分求导，确实正确地产生了瞬子修正 $V^\inst$ 与 $A^\inst$。这也就完成了「$\CX_\magn$ 是 $\CM$ 上所需的『磁』复坐标」这一验证。

\bigskip
\textbf{注}：$\CX_\magn$ 与量子可积系统理论中所谓的「$Q$ 函数」密切相关。\footnote{我们感谢 S. Lukyanov 与我们分享他关于这些函数的笔记。}我们觉得这并非巧合，而是指向与可积场论之间的某种深刻联系。下文的关键方程 式 (5.14) 是热力学贝特拟设的一种形式（见附录 E），这一点更强化了上述感觉。% TODO 合稿改 \eqref{eq:X-integral-mult-explicit}
% TODO 合稿改 \ref{app:TBA}

\subsection{解析性质}

现在考察序对 $(\CX_\magn, \CX_\elec)$ 在 $\zeta$ 平面上的解析行为。

对 $\CX_\elec$ 来说事情很简单：它在 $\zeta \in \IC^\times$ 上解析，在 $\zeta=0,\infty$ 处有本性奇点。对 $\CX_\magn^\sfm$ 也一样，但对完整的 $\CX_\magn$，情况就更复杂：\eqref{eq:FullX2} 中的积分只在离开围道 $\ell_\pm$ 的地方对 $\zeta$ 解析。当 $\zeta$ 跨过这两条围道中的任何一条时，被积函数的极点会越过积分路径。因此，我们关于 $\CX_\magn$ 的表达式定义的是一个分段解析函数，其不连续量由极点的留数决定。引入记号 $(\CX_\magn)_{\ell_+}^{+}$、$(\CX_\magn)_{\ell_+}^{-}$，分别表示 $\zeta$ 沿顺时针或逆时针方向趋近 $\ell_+$ 时 $\CX_\magn$ 的极限，对 $\ell_-$ 也类似。于是不连续量为
\begin{subequations}  \label{eq:X2-disc}
\begin{align}
(\CX_\magn)_{\ell_+}^{+} &= (\CX_\magn)_{\ell_+}^{-} (1 - \CX_\elec^q)^{-q}, \label{eq:X2-disc-a} \\
(\CX_\magn)_{\ell_-}^{+} &= (\CX_\magn)_{\ell_-}^{-} (1 - \CX_\elec^{-q})^q.
\end{align}
\end{subequations}

这些不连续量在下文将起关键作用：我们确实要把它们等同于Kontsevich–Soibelman辛同胚，具体如下。我们把这对复函数 $(\CX_\magn, \CX_\elec)$ 看作给出一个映射
\begin{equation}
\CX: \CM_a \to T_a
\end{equation}
从由 $(\theta_\magn, \theta_\elec)$ 坐标化的实 2-环面 $\CM_a$ 到由 $(X_\magn, X_\elec)$ 坐标化的复化 2-环面 $T_a$。映射 $\CX$ 作为 $\zeta$（以及 $a, \bar{a}, R$）的函数而变化。在第 2.2 节中，我们把 KS 因子 $\CK_\gamma$ 作为 $T_a$ 的辛同胚引入。我们的不连续量 \eqref{eq:X2-disc} 说的是：在射线 $\ell_\pm$ 上，$\CX^{+}$ 与 $\CX^{-}$ 相差一个与 $\CK_{0,\pm q}$ 的复合。% TODO 合稿改 \ref{sec:ReviewKS}

这里出现了一个有趣的现象。考虑 $\CX_\magn$ 在 $a$ 平面中绕 $a=0$ 的单值性。这一单值性有两个来源：\eqref{eq:mono2} 给出的 $\CX_\magn^\sfm$ 的单值性，以及来自 \eqref{eq:X2-disc} 的贡献。这两个贡献实际上互相抵消！这一事实本质上与「半平坦度量在 $a=0$ 处的奇性已被光滑化」有关。另一方面，若把 $\CX_\magn$ 沿 $\zeta = 0$ 解析延拓，它并不回到自身。这一单值性并不带来任何问题。特别地，$\hsymp(\zeta)$ 在 $\zeta = 0$ 附近的行为良好：它只有单极点，正如第 3.1 节的讨论所预期。% TODO 合稿改 \ref{sec:hk-general}

现在考察 $\zeta \to 0,\infty$ 时 $\CX_\elec, \CX_\magn$ 的渐近。$\CX_\elec$ 的渐近可直接由 \eqref{eq:X1-def} 读出：
\begin{equation} \label{eq:X1-asymptotics}
\CX_\elec \sim \begin{cases}
               \exp \left[ \pi R \frac{a}{\zeta} + i \theta_\elec \right] & \text{as } \zeta \to 0, \\
               \exp \left[ \pi R \zeta\bar{a} + i \theta_\elec \right] & \text{as } \zeta \to \infty.
               \end{cases}
\end{equation}
对 $\CX_\magn$，渐近更有意思。当 $\zeta \to 0,\infty$ 时，\eqref{eq:FullX2} 的被积函数简化：有理函数只归结为 $\pm 1$。再展开对数并算出积分，便得
\begin{equation} \label{eq:X2-asymptotics}
\CX_\magn \sim \begin{cases}
               \exp \left[ - i \frac{R q^2}{2\zeta} (a \log (a/\Lambda) - a) + i \theta_\magn + \frac{q}{2\pi i} \sum_{s \neq 0} \frac{1}{s} e^{i s q \theta_\elec} K_0(2 \pi R q \abs{sa}) \right] & \text{as } \zeta \to 0, \\
               \exp \left[ i \frac{\zeta R q^2}{2} (\bar{a} \log (\bar{a}/\bar{\Lambda}) - \bar{a}) + i \theta_\magn - \frac{q}{2\pi i} \sum_{s \neq 0} \frac{1}{s} e^{i s q \theta_\elec} K_0(2 \pi R q \abs{sa}) \right]  & \text{as } \zeta \to \infty.
               \end{cases}
\end{equation}
这些渐近对 $\zeta$ 的一切相位都成立。\eqref{eq:X2-disc} 沿 $\ell_\pm$ 的不连续量不会导致渐近的不连续，因为沿 $\ell_\pm$ 令 $\zeta \to 0,\infty$ 时，跳变指数地接近 $1$：沿 $\ell_+$ 有 $\CX_\elec \to 0$ 指数地快，沿 $\ell_-$ 则 $\CX_\elec^{-1} \to 0$ 指数地快。

另一方面，我们也可以定义另一个函数 $\CX'_\magn$：把 $\CX_\magn$ 顺时针跨过 $\ell_+$ 作解析延拓。由 \eqref{eq:X2-disc-a} 可知，在 $\ell_+$ 的顺时针一侧有
\begin{equation} \label{eq:continued-vs-jumping}
\CX'_\magn = \CX_\magn (1 - \CX_\elec^q)^{q}.
\end{equation}
现在假设把 $\CX'_\magn$ 继续延拓：顺时针延拓到 $\CU_+$ 的边界，再跨入 $\CU_-$。在 $\CU_-$ 中，当 $\zeta \to 0$ 时 $\CX_\elec$ 指数地大。于是由 \eqref{eq:continued-vs-jumping} 可知，$\CX'_\magn$ 与 $\CX_\magn$ 在 $\zeta \to 0$ 时的渐近不同；特别地，$\CX'_\magn$ 并\emph{不}满足 \eqref{eq:X2-asymptotics}。\emph{所以，函数 $\CX_\magn$ 的解析延拓的渐近，并不等于渐近的解析延拓。}这正是斯托克斯现象的标志。

总而言之，我们被引向考虑一个映射 $\CX: \CM_a \to T_a$：它全纯地依赖于 $\zeta$，在 $\zeta \to 0, \infty$ 处呈现斯托克斯现象，且斯托克斯因子由作用在 $T_a$ 上的 KS 辛同胚的复合给出。本文的关键想法是：这一图景对一般的规范理论也成立，而不只是我们这里考虑的阿贝尔理论；事实上，正如我们将在第 5 节看到的，它自动纳入来自相互非局域粒子的多瞬子效应，并给出 $\CM$ 上的精确度量。% TODO 合稿改 \ref{sec:constructing-moduli}

\subsection{微分方程}

上文我们看到，$\CM$ 的超凯勒几何自然地由一个呈现斯托克斯现象的映射 $\CX$ 来描述。斯托克斯现象通常出现在具非正则奇点的线性常微分方程理论中。我们的情形中确实有这样一个微分方程
\begin{equation} \label{eq:de-zeta}
\zeta \partial_\zeta \CX = \CA_\zeta \CX,
\end{equation}
它带有非正则奇点。本节就来识别这个方程。事实上，我们同时还会找到一个伴随方程，它控制对 $S^1$ 半径的依赖：
\begin{equation} \label{eq:de-R}
R \partial_R \CX = \CA_R \CX.
\end{equation}

形如 \eqref{eq:de-zeta}、\eqref{eq:de-R} 的方程，在 $\CA_\zeta$、$\CA_R$、$\CX$ 为有限维矩阵时很常见。此时 $\CA_\zeta$ 与 $\CA_R$ 以从左作矩阵乘的方式作用于 $\CX$，而斯托克斯因子从右作用，所以这两种作用特别地互相对易。我们的情形中，解 $\CX$ 是映射 $\CM_a \to T_a$。斯托克斯因子作为 $T_a$ 的微分同胚作用。$\CA_\zeta$ 与 $\CA_R$ 则作为 $\CM_a$ 的无穷小微分同胚作用，即作为 $(\theta_\magn, \theta_\elec)$ 中的微分算子作用。这两种作用彼此对易，因为它们作用在不同的空间上。

那么，所需方程的来源是什么？它们应当与 $(\CM, g)$ 的某些对称性有关。乍看之下，$(\CM, g)$ 似乎有一个 $U(1)$ 对称性，它只是把 $a \mapsto e^{i \theta} a$。这样的对称性会有显然的物理来源：它来自 $d=4$ 理论的 $U(1)_R$ 对称性。然而我们知道，一旦纳入物质超多重态，这个对称性实际上是反常的。确实，\eqref{eq:A0} 中的 $A^\sfm$ 含有因子 $\log(a/\Lambda)$，而它只在 $a$ 与 $\Lambda$ 的\emph{同时}旋转下才不变。因此这一同时旋转保持度量不变。它并不保持超凯勒形式 $\vec{\omega}$，而是把 $\omega_1$ 与 $\omega_2$ 互相旋转；因此，若与 $\zeta \mapsto e^{i \theta} \zeta$ 作用合起来，它就保持 $\hsymp(\zeta)$ 不变。逐项检查可知，$\CX_\elec$ 与 $\CX_\magn$ 在 $a$、$\Lambda$ 与 $\zeta$ 的这一联合旋转下都不变，这就给出一个微分方程：
\begin{equation} \label{eq:R-symmetry}
\plog{\zeta} \CX = \left( -\plog{\Lambda} + \plog{\bar\Lambda} - \plog{a} + \plog{\bar a} \right) \CX.
\end{equation}
类似地，$d=4$ 理论的反常标度不变性给出一个把 $R$、$a$ 与 $\Lambda$ 同时标度的对称性：
\begin{equation} \label{eq:conformal-symmetry}
\plog{R} \CX = \left( \plog{\Lambda} + \plog{\bar\Lambda} + \plog{a} + \plog{\bar a} \right) \CX.
\end{equation}
这些方程还不是所要的形式 \eqref{eq:de-zeta}、\eqref{eq:de-R}，因为它们仍含有对参数 $(\Lambda, \bar\Lambda, a, \bar a)$ 的导数。所以让我们考察对这些参数的依赖。

$\CX$ 对 $(a, \bar{a})$ 的依赖，完全由它对 $(\theta_\elec, \theta_\magn)$ 的依赖决定，因为要求 $(\CX_\elec, \CX_\magn)$ 在复结构 $J^\pz$ 下全纯。的确，用 \eqref{eq:xi-magn}、\eqref{eq:xi-elec} 作为 $(T^*)^{1,0} \CM$ 的基，我们看到 $\CM$ 上的 Cauchy–Riemann 方程就是
\begin{align}
\partial_a \CX &= \CA_a \CX, \label{eq:de-cr-1} \\
\partial_{\bar a} \CX &= \CA_{\bar a} \CX, \label{eq:de-cr-2}
\end{align}
其中联络形式 $\CA$ 由
\begin{align}
\CA_a &= \frac{1}{\zeta} \left[ - i \pi R \partial_{\theta_\elec} + \pi (V + 2 \pi i R A_{\theta_\elec}) \partial_{\theta_\magn} \right] + 2 \pi A_a \partial_{\theta_\magn}, \\
\CA_{\bar a} &= 2 \pi A_{\bar a} \partial_{\theta_\magn} - \zeta \left[ i \pi R \partial_{\theta_\elec} + \pi (V - 2 \pi i R A_{\theta_\elec}) \partial_{\theta_\magn} \right].
\end{align}
定义。我们也可以类似地处理对 $(\Lambda, \bar\Lambda)$ 的依赖。先注意 $\CX_\elec$ 根本不依赖 $(\Lambda, \bar\Lambda)$。对 $\CX_\magn$ 有 $\Lambda \frac{\partial \CX_\magn}{\partial \Lambda} = \frac{i R q^2 a}{2 \zeta} \CX_\magn$，对 $\bar\Lambda$ 类似。于是，记
\begin{equation}
\CA_\Lambda = \frac{q^2 R a}{2 \zeta} \partial_{\theta_\magn}, \quad \CA_{\bar\Lambda} = \frac{ \zeta q^2 R \bar{a}}{2} \partial_{\theta_\magn},
\end{equation}
就有
\begin{align}
\Lambda \partial_\Lambda \CX &= \CA_\Lambda \CX, \\
\bar\Lambda \partial_{\bar\Lambda} \CX &= \CA_{\bar\Lambda} \CX.
\end{align}
现在可以把方程 \eqref{eq:R-symmetry}、\eqref{eq:conformal-symmetry} 改写为所要的形式 \eqref{eq:de-zeta}、\eqref{eq:de-R}，其中
\begin{align}
\CA_\zeta &= -a \CA_a + \bar{a} \CA_{\bar a} - \Lambda \CA_{\Lambda} + \bar\Lambda \CA_{\bar\Lambda}, \label{eq:zeta-connection-abelian} \\
\CA_R &= a \CA_a + \bar{a} \CA_{\bar a} + \Lambda \CA_{\Lambda} + \bar\Lambda \CA_{\bar\Lambda}.
\end{align}
现在到了关键之处：\eqref{eq:zeta-connection-abelian} 给出的 $\CA_\zeta$ 对 $\zeta$ 的依赖非常简单——它在 $\zeta = 0, \infty$ 处只有单极点：
\begin{equation}
\CA_\zeta = \frac{1}{\zeta} \CA^{(-1)}_\zeta + \CA^{(0)}_\zeta + \zeta \CA^{(1)}_\zeta.
\end{equation}
于是方程 \eqref{eq:de-zeta} 定义了 $\IC\IP^1$ 上的一个亚纯联络，带有两个秩 $1$ 的非正则奇点。这就解释了斯托克斯现象的出现，而我们在上一节已明确看到了它。

一族与 \eqref{eq:de-zeta}、\eqref{eq:de-R}、\eqref{eq:de-cr-1}、\eqref{eq:de-cr-2} 非常相似的微分方程——它们定义了所谓「$tt^*$ 联络」——出现在 \cite{Cecotti:1993rm,MR1213301} 中，背景是 $d=2$ 的 $\CN=(2,2)$ 场论的分析与分类。这种相似并不只是形式上的。特别地，他们对这些方程关于 $\zeta$ 与 $R$ 依赖的解释，同样基于底层场论的 $U(1)_R$ 对称性与标度变换。他们分析中的一个关键点，是联络 $\CA$ 在大 $R$ 时的渐近、斯托克斯因子的显式形式与 $d=2$ 理论中 BPS 态的简并度之间的直接关系。我们的问题中也有类似的关系。事实上，正是这一关系成为理解 KS 穿墙公式的关键。

看一看 \cite{Cecotti:1993rm,MR1213301} 还提示了一个非常有用的技术工具，可用于取得进一步进展：我们应当把这些微分方程转化为 $\CX$ 的黎曼–希尔伯特问题，直接用斯托克斯数据与 $\zeta \to 0, \infty$ 时的渐近来定义。借助这个工具，我们可以立即写下对多个相互非局域 BPS 瞬子的推广。我们将在第 5 节转向这一问题。% TODO 合稿改 \ref{sec:constructing-moduli}

\subsection{更高自旋多重态}

至此我们已相当详细地考察了积分掉单个带电荷超多重态对 $g$ 的修正。同样可以问：单个带电荷的更高自旋多重态——例如含有有质量 $W$ 玻色子的矢量多重态——会带来什么修正？

原则上，这些修正可以通过三维中细致的一圈计算定出。我们改用一个小技巧：考虑 $\CN=4$ 超对称理论的有质量矢量多重态。按 $\CN=2$ 超对称分解，这个多重态含有两个超多重态与一个矢量多重态。另一方面，由于 $\CN=4$ 理论具有更高的超对称，人们预期 $\CM$ 上的度量不会受到任何瞬子修正。理由是：按标准的非重整化定理，Higgs 分支不受修正 \cite{Argyres:1996eh}，而无反常的 $R$ 对称性会把 Higgs 分支与库仑分支混合起来。由此可知，来自 $\CN=2$ 矢量多重态的修正，必定恰好抵消来自两个 $\CN=2$ 超多重态的修正。换言之，至少就这两个 $\CN=2$ 多重态而言，修正是以螺旋度超迹 $\degen(\gamma;u)$ 为权重的。

更一般地，我们可以考虑积分掉任意自旋的 $\CN=2$ 多重态。设 $a_j$ 记第 $j$ 个 $\CN=2$ 多重态（$j=0$ 为超多重态，$j=1$ 为矢量多重态，\dots）对瞬子修正的权重，并归一化为 $a_0 = 1$。上文我们看到 $a_1 = -2$。此外，由「$\CN=4$ 超对称的\emph{任意}多重态的贡献都应当为零」这一事实，得到 $a_{j+2} + 2 a_{j+1} + a_j = 0$。这就定出 $a_j = (-1)^j (j+1)$，所以瞬子修正确实以第二螺旋度超迹为权重。

\subsection{更高秩的推广}
\label{sec:higher-rank}

我们的全部讨论都可以容易地推广到秩 $r$ 的阿贝尔规范理论耦合一组带电荷超多重态的情形。设电荷为 $q_I^{(s)}$，其中 $I=1, \dots, r$ 跑遍各电规范场，$s$ 标记这组超多重态。

Gibbons–Hawking 拟设有 $4r$ 维的推广：底为 $(\IR^3)^r$，纤维为 $(S^1)^r$。底的坐标记作 $(x^{\alpha I}) = \vec x^I$，纤维坐标记作 $\theta_{\magn,I}$，并写
\begin{equation}
g = [V(\vec{x})^{-1}]^{IJ} \left(\frac{d \theta_{\magn,I}}{2\pi} +
A_I(\vec{x})\right)\left(\frac{d \theta_{\magn,J}}{2\pi} +
A_J(\vec{x})\right) + V(\vec{x})_{IJ} d\vec{x}^I d\vec{x}^J,
\end{equation}
其中 $A$ 与 $V$ 由一组微分方程相联系，这组方程说的是
\begin{equation}
\omega^{\alpha} = dx^{\alpha I} \wedge \left( \frac{d\theta_{\magn,I}}{2\pi} + A_I
\right) + \half V_{IJ} \epsilon^{\alpha\beta\gamma} dx^{\beta I } \wedge
dx^{\gamma J}
\end{equation}
对 $\alpha = 1,2,3$ 都是闭 2-形式。

一圈积分给出自然的结果
\begin{equation}
\label{eq:V-sum-r} V_{IJ} =  \mathrm{Im} \tau^0_{IJ} + \sum_{s}
\frac{q^{(s)}_I q^{(s)}_J R}{4 \pi} \sum_{n=-\infty}^\infty \left(
\frac{1}{\sqrt{| q^{(s)}_K R a^K|^2 + (q^{(s)}_K \frac{\theta_\elec^K}{2
\pi} + n)^2}} - \kappa_n \right).
\end{equation}
\eqref{eq:V-sum} 的 Poisson 重求和表明
\begin{equation}
V_{IJ} = V_{IJ}^\sfm + V_{IJ}^\inst,
\end{equation}
其中
\begin{align}
V_{IJ}^\sfm &= \mathrm{Im} \tau^0_{IJ} - \sum_s \frac{q^{(s)}_I
q^{(s)}_J R}{4\pi} \left(\log  \frac{q^{(s)}_K a^K}{\Lambda} + \log \frac{q^{(s)}_K \bar{a}^K}{\bar{\Lambda}} \right), \\
V_{IJ}^\inst &= \sum_s \frac{q^{(s)}_I q^{(s)}_J R}{2\pi}
\sum_{n\neq 0} e^{i n q^{(s)}_K \theta_\elec^K} K_0(2 \pi R \abs{n
q^{(s)}_I a^I}). \label{eq:V-inst-r}
\end{align}
又
\begin{equation}
A_I = A_I^\sfm + A_I^\inst,
\end{equation}
其中
\begin{align}
A_I^\sfm &= \mathrm{Re} \tau^0_{IJ} \frac{d \theta_\elec^J}{2 \pi} + \sum_s
\frac{i q^{(s)}_I
q^{(s)}_J}{8 \pi^2} \left( \log \frac{q^{(s)}_K a^K}{\Lambda} - \log \frac{q^{(s)}_K \bar{a}^K}{\bar{\Lambda}} \right) d\theta_\elec^J, \label{eq:A0-rankr} \\
A_I^\inst &=  - \frac{q^{(s)}_I q^{(s)}_J R}{4 \pi} \left(
\frac{da^J}{q^{(s)}_K a^K} - \frac{d\bar a^J}{q^{(s)}_K \bar a^K}
\right) \sum_{n \neq 0} (\sgn n) e^{i n q^{(s)}_I \theta_\elec^I}
\abs{q^{(s)}_K a^K} K_1(2 \pi R \abs{n q^{(s)}_K a^K}).
\label{eq:A-inst-rankr}
\end{align}
在大 $R$ 时，$V$ 与 $A$ 中的领头项是 $V^\sfm$ 与 $A^\sfm$。只保留这些项，$g$ 就成为半平坦度量，其
\begin{equation}
\tau_{IJ} =\tau^0_{IJ} + \sum_s \frac{q^{(s)}_I q^{(s)}_J}{2 \pi i}
\log \frac{q^{(s)}_K a^K}{\Lambda}.
\end{equation}
这就是 $d=4$ 中积分掉超多重态所得的耦合。

全纯辛形式为
\begin{equation}
\hsymp(\zeta) = - \frac{1}{4 \pi^2 R} \xi^I_e \wedge \xi_{m,I}
\end{equation}
其中
\begin{align}
\xi_e^I &= i d \theta_\elec^I + \pi R \left(\frac{da^I}{\zeta} +
\zeta d \bar a^I \right), \\
\xi_{m,I} &= i d \theta_{\magn,I} + 2 \pi i A_I + i \pi V_{IJ} \left(
\frac{da^J}{\zeta} - \zeta d \bar a^J \right).
\end{align}
如前面一样，电坐标与其半平坦近似一致：
\begin{equation}
\CX^I_e = \exp \left[ \pi R \frac{a^I}{\zeta} + i \theta_\elec^I +
\pi R \zeta \bar a^I \right].
\end{equation}
磁坐标的半平坦近似为
\begin{multline}
\CX^\sfm_{m,I} = \exp \Biggl[ \frac{\pi R}{\zeta} \left(\tau^0_{IJ}
a^J  + \sum_s \frac{q^{(s)}_I}{2 \pi i} q^{(s)}_K a^K \log
\frac{q^{(s)}_K a^K}{e \Lambda}\right) + i \theta_{\magn,I}\,+ \\
\pi R \zeta
\left( \tau^0_{IJ} \bar a^J + \sum_s \frac{q^{(s)}_I}{2 \pi
i} q^{(s)}_K \bar a^K \log \frac{q^{(s)}_K \bar a^K}{e \bar
\Lambda}\right) \Biggr].
\end{multline}
完整的磁坐标由如下积分公式给出：
\begin{equation}\label{eq:FullX2-r}
\begin{split}
\CX_{\magn,I} = \CX_{\magn,I}^\sfm \exp\Biggl[ \sum_s \frac{i
q^{(s)}_I}{4\pi} & \int_{\ell^s_+} \frac{d \zeta'}{\zeta'}
\frac{\zeta' + \zeta}{\zeta' - \zeta} \log[1-\prod_J \CX^J_\elec(\zeta')^{q^{(s)}_J}]  \\
- \frac{i q^{(s)}_I}{4\pi} & \int_{\ell^s_-}
 \frac{d \zeta'}{\zeta'} \frac{\zeta' + \zeta}{\zeta' - \zeta} \log[1-\prod_J \CX^J_\elec(\zeta')^{-q^{(s)}_J}] \Biggr],
\end{split}
\end{equation}
其中 $\ell^s_\pm$ 是位于两个半平面
\begin{equation}
\CU^s_\pm = \left\{ \zeta: \pm {\rm Re} \frac{a^K q^{(s)}_K}{\zeta}<0 \right\}.
\end{equation}
内、连接 $0$ 与 $\infty$ 的任意路径。

注意 $\CX_{m,I}$ 对每个超多重态都有不连续量，如前面一样，它们由 KS 辛同胚 $\CK_{0,q_I^{(s)}}$ 给出。

\begin{yizhu}
本块沿用第 1 章译注已交代的两条约定，此处只作提示：其一，instanton 在本节一律指 $d=4$ 理论中 BPS 粒子的世界线绕 $S^1$ 缠绕所形成的组态，译「BPS 瞬子」，以便与四维瞬子相区别；其二，$\degen(\gamma;u)$ 即原文的 \verb|\degen|，为第二螺旋度超迹，沿用 $\Omega$ 这一符号（不可与 KS 不变量 $\hat\degen(\gamma;u)$ 混同）。
\end{yizhu}

% ch05a —— §5.1–5.3（原文 wcf-paper.tex 2129–2544 行）

\section{相互非局域修正} \label{sec:constructing-moduli}

正如我们刚刚在最简单的非平凡情形中所见，精确的超凯勒度量 $g$ 并不等于半平坦度量 $g^{\sfm}$，因为存在来自瞬子的量子修正，而这些瞬子对应于 $d=4$ 的 BPS 态。

一般地，对每一个支撑 BPS 态的电荷 $\gamma$，我们都预期有这样一项量子修正。这些修正应以 BPS 简并度 $\degen(\gamma;u)$ 为权重。然而我们知道，当 $u$ 穿过边缘稳定墙时，$\degen(\gamma;u)$ 会发生跳变！于是似乎出现了一个困惑：$g$ 的量子修正难道不也会不连续地跳变吗？这与场论中「$g$ 应当光滑」的预期如何相容？

本节将给出大 $R$ 下精确的超凯勒度量 $g$ 的一个显式构造。我们将看到，只要 Kontsevich--Soibelman 穿墙公式成立，$g$ 确实是光滑的。这就是我们对 WCF 的物理解释。

把 $g$ 在 $R\to\infty$ 附近展开，我们便找到了上述困惑的解答：除单个 BPS 粒子的贡献之外，还有多粒子贡献。单粒子贡献中的不连续性，被多粒子部分的不连续性所补偿。关于 $\CN=1$ 语境中的相关讨论，见 \cite{GarciaEtxebarria:2007zv}。

\subsection{黎曼--希尔伯特问题的定义} \label{sec:defining-rh}

我们从第 4 节所研究的阿贝尔理论中汲取灵感，通过在 $\zeta$ 平面中求解一个黎曼--希尔伯特问题来构造该度量。% TODO 合稿改 \ref{sec:periodicnut}
我们起初在固定的 $u\in\vm$（4d 模空间）处工作，远离各边缘稳定墙。我们还在 $u$ 处取定一个二次加细 $\sigma$。

黎曼--希尔伯特问题是用映射 $\CX$ 表述的：$\CX$ 从实环面 $\widetilde\CM_u$（$\CM$ 为 3d 模空间）映入复化辛 $2r$-环面 $\tilde T_u$，后者是我们在回顾 Kontsevich--Soibelman 公式时（第 2.2 节）引入的。% TODO 合稿改 \ref{sec:ReviewKS}
给定任何这样的 $\CX$，我们把 $\tilde T_u$ 上的坐标函数 $X_\gamma$ 拉回，得到 $\widetilde\CM_u$ 上的函数 $\CX_\gamma$，其定义为 $\CX_\gamma(\theta) = X_\gamma(\CX(\theta))$。特别地，第 3.3 节所给的诸 $\CX^{\sfm}_\gamma$ 来自一个参考映射 $\CX^{\sfm}$；它就是我们下面所构造的 $\CX$ 的第零阶近似。% TODO 合稿改 \ref{sec:semiflat-twistor}

为表述这个黎曼--希尔伯特问题，我们需要确定 $\CX$ 当 $\zeta\to0,\infty$ 时的\emph{渐近行为}，以及它在 $\zeta$ 平面中的\emph{不连续性}。

我们先从渐近行为说起。引入
\begin{equation} \label{eq:def-upsilon}
\Upsilon := \CX (\CX^{\sfm})^{-1}.
\end{equation}
在本节中，我们对映射的复合采用一种非传统的记号：$(fg)(x)$ 意为 $g(f(x))$。\footnote{这种记号的一个好处，从下面的复合即可看出：$A \stackrel{f}{\to} B \stackrel{g}{\to} C$ 复合为 $A \stackrel{fg}{\to} C$。第二个好处将在第 5.5 节显现。}% TODO 合稿改 \ref{sec:differential-problem}
于是 $\Upsilon$ 是从 $\widetilde\CM_u$ 到自身的映射（更确切地说，是到其复化）。具体地，$\Upsilon$ 把
\begin{equation} \label{eq:upsilon-concrete}
e^{i \theta_i} \mapsto \CX^i(\theta) \exp\biggl[  - \pi R
\frac{Z_i}{\zeta} - \pi R  \zeta \bar Z_i \biggr].
\end{equation}
我们要求 $\Upsilon$ 当 $\zeta\to0$ 与 $\zeta\to\infty$ 时的极限存在，
\begin{equation} \label{eq:RH-asymptotics}
\lim_{\zeta \to 0} \Upsilon = \Upsilon_0, \quad \lim_{\zeta \to \infty} \Upsilon = \Upsilon_\infty,
\end{equation}
并且满足
\begin{equation} \label{eq:upsilon-reality}
\Upsilon_0 = \bar \Upsilon_\infty.
\end{equation}

接下来我们需要指定 $\CX$ 的不连续性，这里把 $\CX$ 视为 $\zeta$ 的分段解析函数。暂时假定 $u$ 不在任何边缘稳定墙上。不连续性将由与 BPS 态相伴的 Kontsevich--Soibelman 辛同胚 $\CK_\gamma:\tilde T_u\to\tilde T_u$ 给出。我们给 $\zeta$ 平面中过原点的每一条射线 $\ell$ 配以 $\clat_u$ 的一个子集
\begin{equation} \label{eq:clat-ell}
(\clat_u)_\ell := \{ \gamma: Z_\gamma(u)/\zeta \in \IR_- \text{ 对 } \zeta \in \ell \},
\end{equation}
以及相应地，对 BPS 态取的一个乘积：
\begin{equation}
S_\ell := \prod_{\gamma \in (\clat_u)_\ell} \CK_{\gamma}^{\degen(\gamma;u)}.
\end{equation}
（由于 $u$ 不在墙上，$(\clat_u)_\ell$ 至多是一维的，而诸 $\gamma\in(\clat_u)_\ell$ 对应的 $\CK_\gamma$ 彼此对易；因此无须指定这个乘积中的次序。）由于电荷格 $\clat_u$ 是可数的，除可数条射线 $\ell$ 之外，我们都有 $(\clat_u)_\ell=\emptyset$，从而 $S_\ell=1$。我们把满足 $(\clat_u)_\ell\neq\emptyset$ 的射线称为「BPS 射线」。

最典范的不连续性取法是要求
\begin{equation} \label{eq:RH-jumps}
\CX^{+} = \CX^{-} S_\ell
\end{equation}
其中 $\CX^{+}$、$\CX^{-}$ 分别是 $\zeta$ 沿顺时针、逆时针方向趋近 $\ell$ 时 $\CX$ 的极限。这是对我们在第 4 节所得结果的最直截了当的推广：在那里，我们得到在 $\zeta$ 中分段解析的映射 $\CX=(\CX_\magn,\CX_\elec)$，它有两条 BPS 射线 $\ell_\pm$；跨过 $\ell_\pm$ 时 $\CX_\elec$ 连续，而 $\CX_\magn$ 按式 (4.52) 跳变。% TODO 合稿改 \eqref{eq:X2-disc}；「第 4 节」改 \ref{sec:periodicnut}
这两条 BPS 射线对应于电荷为 $(0,\pm q)$ 的单个超多重态，其不连续性恰为 \eqref{eq:RH-jumps}，其中 $S_{\pm\ell}=\CK_{0,\pm q}$。现在我们把许多 BPS 粒子都纳入进来，办法只是要求沿许多 BPS 射线发生跳变。在这种更一般的情形下，将不存在处处连续的 $\CX_\gamma$。

至此，我们的黎曼--希尔伯特问题已经表述完毕。它的解并不唯一：确切地说，它只确定到相差一个变换 $\CX\to b\CX$，其中 $b$ 是 $\widetilde\CM_u$ 的任意微分同胚。下一节求解该问题时，我们将用一种方便的方式固定这一不确定性。

这个黎曼--希尔伯特问题或许显得有些不合常规，因为它是用 $\CX$ 与 $S_\ell$ 表述的，而它们不是线性映射，而是更一般的流形之间的映射。有顾虑的读者尽可把问题「线性化」：不去考虑 $\CX$，而考虑拉回操作 $\CX^*: C^\infty(\tilde T_u)\to C^\infty(\widetilde\CM_u)$。这样做的代价是：并非每个映射 $Q: C^\infty(\tilde T_u)\to C^\infty(\widetilde\CM_u)$ 都能写成某个映射 $\CX$ 的 $\CX^*$；因此，即使我们为线性化后的黎曼--希尔伯特问题找到了一个解 $Q$，还要额外费心验证 $Q=\CX^*$ 对某个 $\CX$ 成立。所幸黎曼--希尔伯特问题的「函子性」会来救场：$Q$ 是某个 $\CX$ 的 $\CX^*$，当且仅当它保持乘法，即 $Q(fg)=Q(f)Q(g)$。由于定义该线性问题的全部数据都与这一结构相容，其解也与之相容。

最后我们讨论本问题的一个实性（reality）性质，它对我们构造超凯勒度量至关重要。由于关系 $\degen(\gamma;u)=\degen(-\gamma;u)$，我们的不连续性条件具有一种离散对称性：给定任一解 $\CX$，我们都能得到另一个解，记作 $\bar\CX$，其定义为
\begin{equation}
\bar\CX_\gamma(\zeta) = \overline{\CX_{-\gamma}(-1/\bar\zeta)}.
\end{equation}
我们断言，事实上我们的解在这一变换下不变，即
\begin{equation} \label{eq:reality-X}
\CX = \bar \CX.
\end{equation}
为看出这一点，考虑映射 $Y=\bar\CX\CX^{-1}$。由于 $\bar\CX$ 与 $\CX$ 具有相同的不连续性，$Y$ 实际上是 $\zeta$ 的\emph{解析}函数。另一方面，由渐近条件 \eqref{eq:upsilon-reality}，当 $\zeta\to0,\infty$ 时 $Y\to1$。于是由刘维尔定理得 $Y=1$。

\subsection{KS 公式的作用} \label{sec:continuity-rh}

现在到了关键之处，这也是我们撰写本文的主要理由。假设对远离边缘稳定墙的任一固定 $u$，我们都找到一个解 $\CX$。那么它作为 $\zeta$ 的函数的行为便完全确定了：除在 BPS 射线处按 \eqref{eq:RH-jumps} 跳变外，它处处连续。但关于它作为 $u$ 的函数的行为，我们能说些什么呢？

我们的黎曼--希尔伯特问题中的 $u$ 依赖性来自两处。一处是渐近边界条件 \eqref{eq:RH-asymptotics}；这一依赖当然是连续的。另一处是不连续性处方 \eqref{eq:RH-jumps}。只要 $u$ 远离边缘稳定墙，这里的依赖也是连续的。但当我们穿过墙时，会发生什么？设 $u_w$ 为墙上的一个一般点。当 $u$ 从墙的一侧趋于 $u_w$ 时，对应于电荷 $\gamma=n\gamma_1+m\gamma_2$ 的诸 BPS 射线彼此相撞，汇合成一条射线 $\ell$。设 $A$ 为 $\CX$ 跨过这一组射线的总不连续性，
\begin{equation} \label{eq:disc-product}
A = \prod^\ccwarrow_{\substack{\gamma = n \gamma_1 + m \gamma_2 \\ m>0,\,n>0}}
\CK_\gamma^{\degen(\gamma;u)}.
\end{equation}
假定从这一侧取的极限 $\lim_{u\to u_w}\CX$ 存在，则它是如下黎曼--希尔伯特问题的解：跨过 $\ell$ 的不连续性为 $A$，而沿其他所有 BPS 射线的不连续性仍如前指定。

另一方面，我们也可以考虑从墙的另一侧取的极限 $\lim_{u\to u_w}\CX$。要使这两个极限一致，充分必要条件是：它们是\emph{同一个}黎曼--希尔伯特问题的解。于是这就要求由 \eqref{eq:disc-product} 算出的 $A$ 在墙的两侧相同。正如我们在第 2.2 节所回顾的，这正是 KS 穿墙公式的内容！\footnote{\eqref{eq:disc-product} 中的乘积取逆时针次序，而第 2.2 节中取顺时针次序；这一差别来自我们在第 \ref{sec:constructing-moduli} 节对映射复合所采用的非传统约定。}% TODO 合稿改 \ref{sec:ReviewKS}

我们于是断言：\emph{假定 BPS 简并度服从 KS 公式}，则黎曼--希尔伯特问题的解 $\CX$ 作为 $u$ 与 $\zeta$ 的函数是连续的，只有在 BPS 射线处例外。此外，跨过 BPS 射线的不连续性由一个辛同胚给出。

\subsection{求解黎曼--希尔伯特问题} \label{sec:solving-rh}

表述完黎曼--希尔伯特问题之后，我们想看到它有解，并理解其大 $R$ 行为。与第 4 节所考虑的那些简单情形不同——那里所有的 $S_\ell$ 彼此对易——这里我们无法为所求的 $\CX$ 写出显式的积分公式，只能走一条更间接的路线。% TODO 合稿改 \ref{sec:periodicnut}
我们利用这样一件事实：本问题的结构与文献 \cite{MR1213301,Cecotti:1993rm} 中所考虑的一个问题极为相似。事实上，我们的问题正是那里所考虑问题的无穷维版本。

在 \cite{Cecotti:1993rm} 中，黎曼--希尔伯特问题被改写成关于 $\Upsilon(\zeta)$ 的类比物 $\Phi$ 的一个积分方程。当 $R$ 足够大时，该方程把 $\Phi$ 描述为单位矩阵的一个小修正。因此它可以迭代求解：这既证明了当 $R$ 足够大时解的存在性，也给出第零阶近似 $\Phi=1$ 的首阶修正的显式公式。这些首阶修正直接用不连续因子表达。

这正是我们希望获得的、关于映射 $\CX$ 的那类信息。一种直接的做法，是写下 \cite{Cecotti:1993rm} 中积分方程的无穷维类比。这一做法只直接适用于线性的黎曼--希尔伯特问题，因此人们必须先过渡到上一小节末尾提到的线性问题。积分方程的解随后将给出函数空间之间的一个线性映射；如我们所述，这个线性映射将是某个映射 $\CX:\widetilde\CM_u\to\tilde T_u$ 的 $\CX^*$。

一个小问题是：倘若严格照搬 \cite{Cecotti:1993rm} 的处方，我们得到的解将满足边界条件 $\Upsilon_0=1$。而我们的构造需要选择不同的边界条件，即 \eqref{eq:upsilon-reality}，它的好处是与实性条件 $\CX_\gamma(\zeta)=\overline{\CX_{-\gamma}(-1/\bar\zeta)}$ 相容。\footnote{本节中我们把 $\CX=\CX(\zeta)$ 显式写出，即把 $\CX$ 视为随 $\zeta$ 变化的\emph{映射}，从而略去它对 $\widetilde\CM_u$ 的坐标 $\theta$ 的依赖。}幸运的是，只需对积分核稍作修改，就很容易写出一个顾及这一不同边界条件的积分方程变体。

这一策略似乎足以证明解的存在性，但它有一个重要的缺点：迭代求解的中间步骤不必具有 $\CX^*$ 的形式。若能有一种表述，使逼近格式的每一步本身都是映射 $\widetilde\CM_u\to\tilde T_u$，那将很有用。利用 $\tilde T_u$ 上的阿贝尔群结构，只要写出如下积分方程即可做到：
\begin{equation} \label{eq:X-integral-mult}
\CX_\gamma(\zeta) = \CX^{\sfm}_\gamma(\zeta) \exp \left[ \frac{1}{4 \pi
i}\sum_{\ell} \int_{\ell } \frac{d\zeta'}{\zeta'} \frac{\zeta' +
\zeta}{\zeta' - \zeta} \log \frac {\CX_\gamma(\zeta')}{(\CX S_\ell)_\gamma(\zeta')}\right]
\end{equation}
这里求和遍历各 BPS 射线 $\ell$。\eqref{eq:X-integral-mult} 的任何解都满足不连续性条件 \eqref{eq:RH-jumps}。此外，我们对积分核的选择保证该解还满足实性条件 \eqref{eq:upsilon-reality}。因此，\eqref{eq:X-integral-mult} 的解就是黎曼--希尔伯特问题的解。\footnote{注意，尽管黎曼--希尔伯特问题在 $\widetilde\CM_u$ 的微分同胚下不变，方程 \eqref{eq:X-integral-mult} 却并非如此；它的解是唯一的，而不是唯一到相差一个微分同胚。}

利用由式 (2.16) 给出的 Kontsevich--Soibelman 因子的显式形式，我们有 % TODO 合稿改 \eqref{eq:ks-symplect}
\begin{equation}
(\CX S_\ell)_\gamma = \CX_\gamma \prod_{\gamma' \in (\clat_u)_\ell} (1-
\sigma(\gamma') \CX_{\gamma'})^{\degen(\gamma';u)\langle
\gamma,\gamma' \rangle}
\end{equation}
（其中 $(\clat_u)_\ell$ 由式 \eqref{eq:clat-ell} 定义）。把它代入 \eqref{eq:X-integral-mult}，便得到关于 $\CX$ 的最终积分方程：
\begin{equation} \label{eq:X-integral-mult-explicit}
\CX_\gamma(\zeta) = \CX^{\sfm}_\gamma(\zeta) \exp \left[ -\frac{1}{4
\pi i} \sum_{\gamma'} \degen(\gamma';u) \langle \gamma,\gamma'
\rangle \int_{\ell_{\gamma'}} \frac{d\zeta'}{\zeta'} \frac{\zeta' +
\zeta}{\zeta' - \zeta} \log (1- \sigma(\gamma')
\CX_{\gamma'}(\zeta'))\right].
\end{equation}
正如我们已经提到的，方程 \eqref{eq:X-integral-mult-explicit} 是热力学贝特拟设的一种形式。见附录 E。% TODO 合稿改 \ref{app:TBA}

在附录 C 中，我们论证 \eqref{eq:X-integral-mult-explicit} 当 $R$ 足够大时有解，并描述 $u$ 远离各墙时其 $R\to\infty$ 的展开。% TODO 合稿改 \ref{app:integral-asymptotics}
第一个非平凡近似是
\begin{equation} \label{eq:X-first-approx}
\CX_\gamma(\zeta) \sim \CX^{\sfm}_\gamma(\zeta) \exp \left[ -\frac{1}{4
\pi i}  \sum_{\gamma'} \degen(\gamma';u) \langle \gamma,\gamma'
\rangle \int_{\ell_{\gamma'}} \frac{d\zeta'}{\zeta'} \frac{\zeta' +
\zeta}{\zeta' - \zeta} \log (1- \sigma(\gamma')
\CX^{\sfm}_{\gamma'}(\zeta'))\right],
\end{equation}
它本质上是我们在阿贝尔理论中找到的单瞬子修正的线性叠加。更高阶修正涉及诸 $\degen(\gamma';u)$ 的多重线性式，并带有某种 $R$ 依赖性，这种依赖性表明它们是多瞬子贡献。

我们在附录 C 中的论证，与 \cite{Cecotti:1993rm} 在有限维 $tt^*$ 语境中所给的论证密切相关。事实上，即使在有限维情形，我们的做法也使渐近分析得到简化；因此我们在附录 C 中也重新分析了那一情形。% TODO 合稿改 \ref{app:integral-asymptotics}（两处）

\subsubsection*{整体性问题}

通过求解黎曼--希尔伯特问题，我们得到了一个映射 $\CX:\widetilde\CM_u\to\tilde T_u$，它依赖于局部二次加细 $\sigma(\gamma)$ 的选取。这一选取通过不连续性 $\CK_\gamma$ 的定义影响黎曼--希尔伯特问题。然而，解 $\CX$ 对 $\sigma$ 的依赖方式很简单。回忆：对任意两个细化 $\sigma,\sigma'$，都存在某个 $c(\sigma,\sigma')\in\clat_u^*/2\clat_u^*$，使得 $\sigma(\gamma)\sigma'(\gamma)=(-1)^{\gamma\cdot c(\sigma,\sigma')}$。给定细化 $\sigma$ 下 \eqref{eq:X-integral-mult} 的一个解 $\CX^{[\sigma]}$，就有细化 $\sigma'$ 下相应的解 $\CX^{[\sigma']}$：
\begin{equation}
\CX_\gamma^{[\sigma']}(u,\theta;\zeta) = (-1)^{\gamma \cdot c(\sigma,\sigma')} \CX_\gamma^{[\sigma]}(u,\theta + c \pi;\zeta).
\end{equation}
由此可见，若利用细化把 $\widetilde\CM_u\simeq\CM_u$ 与 $\tilde T_u\simeq T_u$ 等同起来，我们就得到与细化选取\emph{无关}的 $\CX:\CM_u\to T_u$。

% ch05b —— §5.4–5.7：辛形式的构造、微分方程、度量及其大 $R$ 渐近、与物理度量的比较（原文 wcf-paper.tex 2545–3001 行）
\subsection{辛形式的构造}
\label{sec:symplectic-form}

至此，我们已经解出黎曼--希尔伯特问题，得到了一个映射 $\CX:\CM_u\to T_u$，它满足渐近条件 式 (5.3)、跳变条件 式 (5.7) 与实条件 式 (5.9)。% TODO 合稿改 \eqref：式 (5.3)(5.7)(5.9) 属前面各块
现在让 $u$ 变动，便得到一个映射 $\CX:\CM\to T$。于是我们拉回 $T$ 上典范的逐纤维辛形式，在 $\CM$ 上构造出一个复 $2$-形式 $\hsymp(\zeta)$：
\begin{equation} \label{eq:omega-pullback}
\hsymp(\zeta) = \frac{1}{4 \pi^2 R} \CX^* \hsymp^T = \frac{1}{8 \pi^2 R} \epsilon_{ij} \frac{d\CX_{\gamma^i}}{\CX_{\gamma^i}} \wedge \frac{d\CX_{\gamma^j}}{\CX_{\gamma^j}}.
\end{equation}
由式 \eqref{eq:omega-pullback} 可以直接得到 $\hsymp(\zeta)$ 的几条性质：
\begin{itemize}
\item 尽管 $\CX$ 关于 $\zeta$ 只是分段解析的，$\hsymp(\zeta)$ 却是真正解析的（因为诸不连续量 $S_\gamma$ 都是辛同胚，也就是说它们保持 $\hsymp^T$）。

\item 利用式 \eqref{eq:reality-X}，有 $\hsymp(- 1 / \bar\zeta) = \overline {\hsymp(\zeta)}$。

\item 当 $\zeta \to 0,\infty$ 时，利用 $\CX$ 的渐近行为 式 (5.3) 与 $\hsymp^\sf(\zeta)$ 的显式表达式 式 (3.15)，可以确定 $\hsymp(\zeta)$ 的行为。% TODO 合稿改 \eqref：式 (5.3)、(3.15) 属前面各块
我们发现，两种情形下 $\hsymp(\zeta)$ 都只有单极点，留数为
\begin{equation} \label{eq:omega-residues}
\Res_{\zeta = 0} \hsymp(\zeta) = \frac{i}{8 \pi} \Upsilon_0^*\inprod{dZ, d\theta}, \quad \Res_{\zeta = \infty} \hsymp(\zeta) = -\frac{i}{8 \pi} \Upsilon_\infty^*\inprod{d\bar Z, d\theta}.
\end{equation}

\item 利用 $\lim_{R \to \infty} \CX = \CX^\sf$ 可知，当 $R$ 充分大时 $\hsymp(\zeta)$ 非退化（在全纯意义下）。
\end{itemize}
这些性质在我们构造超凯勒度量时都将起重要作用。

\subsection{微分方程}
\label{sec:differential-problem}

我们的黎曼--希尔伯特问题，是用一些普适的不连续因子（它们局部地与规范理论的一切参数无关）来表述的，与之配套的渐近行为则由诸函数 $\CX^\sf_\gamma$ 给出，而这些函数对参数的依赖方式极为简单。本节将依照一套标准做法说明：这意味着解 $\CX$ 满足一族微分方程。

正如我们将看到的，这些方程的物理意义相当透明。其中一组表达了这样一件事实：解黎曼--希尔伯特问题所得的诸函数 $\CX(\zeta)$，在复结构 $J^\pz$ 下关于 $\CM$ 全纯。这组方程对于构造超凯勒度量是本质性的。另一对则描述重整化群流与一个 $U(1)_R$ 对称性的作用。它们对于把度量与 KS 穿墙公式联系起来很重要。

极为相似的一族方程，曾在「$tt^*$ 几何」的故事中起过关键作用；那一套东西出现在有质量的 $\CN=(2,2)$ 二维理论中 \cite{Cecotti:1993rm,Cecotti:1992qh,Cecotti:1991me,Cecotti:1993vy}。

我们首先回顾一点：我们的黎曼--希尔伯特问题在 $\IC\IP^1$ 上的解 $\CX$ 只是分段解析的，沿 $\IC\IP^1$ 中的各条射线 $\ell$，它带有 $\CX \to \CX S_\ell$ 形式的跳变。因此转而考虑\footnote{这是标准记号，但在我们的语境中它多少有些暗示意味，所以这里给出较长的说明。用 $\zeta \partial_\zeta$ 作用，使映射 $\CX$ 作无穷小变动，得到 $T$ 上的一个向量场，我们称之为 $\zeta \partial_\zeta \CX$。然后借助 $\CX$ 把它拉回，得到 $\CM_u$ 上的向量场 $\CA_\zeta$。我们把这一拉回运算记作 $\CX^{-1}$；由于我们采用的复合约定是非标准的，这个 $\CX^{-1}$ 出现在右边而不是左边，这使得我们的方程与黎曼--希尔伯特问题的通常形式一致；事实上，正是为了这一致性，我们在第 5 节才采用非标准的约定。% TODO 合稿改 \ref（Section \ref{sec:constructing-moduli}）
在局部坐标下，人们会写成
\begin{equation}
\CA_\zeta =  \frac{\partial \CX^i}{\partial \zeta}
\left[ \left(\frac{\partial \CX}{\partial \theta}\right)^{-1}
\right]^i_j \frac{\partial}{\partial \theta^j}.
\end{equation}
}
\begin{equation}
\CA_\zeta := \zeta \partial_\zeta \CX \CX^{-1}.
\end{equation}
$\CX$ 沿 BPS 射线的诸不连续量在 $\CA_\zeta$ 中相互抵消，因此 $\CA_\zeta$ 关于 $\zeta$ 是真正解析的——可能除了 $\zeta = 0, \infty$ 之外，那里 $\CX$ 变得奇异。于是我们可以把 $\CX$ 看作关于 $\zeta$ 的一个常微分方程的解，
\begin{equation} \label{eq:zeta-eq-general}
\zeta \partial_\zeta \CX = \CA_\zeta \CX.
\end{equation}
借助关于 $\CX$ 的渐近信息，可以把这一方程描述得相当具体。首先注意 $\CX^\sf$ 满足同一形式的方程。为写出它，先引入 $\CM_u$ 上的两个向量场，
\begin{equation}
\CA_\zeta^{(-1),\sf} := i \pi Z \cdot \partial_\theta, \quad \CA_\zeta^{(1),\sf} := i \pi \bar{Z} \cdot \partial_\theta.
\end{equation}
于是有
\begin{equation} \label{eq:zeta-eq-sf}
\zeta \partial_\zeta \CX^\sf = \CA^\sf_\zeta \CX^\sf,
\end{equation}
其中
\begin{equation}
\CA^\sf_\zeta = \frac{1}{\zeta} \CA_\zeta^{(-1),\sf} + \zeta \CA_\zeta^{(1),\sf}.
\end{equation}
要点在于，$\CA^\sf_\zeta$ 对 $\zeta$ 的依赖非常简单：只是在 $\zeta = 0,\infty$ 处各有一个单极点。我们可以把这一信息转换成关于 $\CA_\zeta$ 的信息，因为由式 (5.3) 知道 $\Upsilon = \CX (\CX^\sf)^{-1}$ 在 $\zeta = 0, \infty$ 两处都保持有限。% TODO 合稿改 \eqref：式 (5.3) 属前面各块
这说明 $\CA_\zeta$ 在 $\zeta = 0, \infty$ 处同样只有单极点，甚至还能定出留数，
\begin{equation}
\CA_\zeta = \frac{1}{\zeta}\CA_\zeta^{(-1)} + \CA_\zeta^{(0)} +  \zeta \CA_\zeta^{(1)},
\end{equation}
其中
\begin{equation}
\CA_\zeta^{(-1)} = \Upsilon_0 \CA_\zeta^{(-1),\sf} \Upsilon_0^{-1}, \quad \CA_\zeta^{(1)} = \Upsilon_\infty \CA_\zeta^{(1),\sf} \Upsilon_\infty^{-1}.
\end{equation}
因此我们看到，式 \eqref{eq:zeta-eq-general} 在 $\IC\IP^1$ 上定义了一个平坦联络 $\zeta \partial_\zeta - \CA_\zeta$，其取值在 $\CM_u$ 上向量场构成的无穷维代数中，并在 $\zeta = 0, \infty$ 处有秩 $1$ 的非正则奇点。$\CX$ 是这个联络的平坦截面。

于是，我们关于黎曼--希尔伯特问题的解，直接导致 $\IC\IP^1$ 上一个平坦联络的构造。事实上，这在常微分方程理论中是标准手法。我们得到的联络在 $\zeta = 0$ 与 $\zeta = \infty$ 处有非正则奇点，因而呈现斯托克斯现象。黎曼--希尔伯特构造的优点之一是：斯托克斯因子很容易确定——它们就是进入黎曼--希尔伯特问题的那些不连续量 $S_\ell$。

上面的讨论有一个重要的推广。我们面对的并非单个黎曼--希尔伯特问题，而是一整族这样的问题，它们随一些额外参数而变动。这些参数包括 $\vm$ 上的坐标 $u^i$，以及标度 $\Lambda$、$S^1$ 的半径 $R$，也许还有一些裸规范耦合 $\tau^0$。（眼下我们\emph{不}引入质量参数；但见下文第 6 节。）% TODO 合稿改 \ref（Section \ref{sec:masses}）
我们引入通用记号 $t^n$ 来统指这一切参数。

重要的是，定义黎曼--希尔伯特问题的诸不连续量 $S_\ell$ 不依赖于任何 $t^n$。因此，正如上面处理对 $\zeta$ 的依赖时那样，我们考虑
\begin{equation}
\CA_n := \partial_{t^n} \CX \CX^{-1}.
\end{equation}
与前面一样，$\CX$ 的诸不连续量相互抵消，所以 $\CA_n$ 关于 $\zeta$ 解析，只在 $\zeta = 0, \infty$ 处例外。同样与前面一样，我们可以先检查 $\CA^\sf_n := \partial_{t^n} \CX^\sf (\CX^\sf)^{-1}$ 的行为，以此控制奇点附近的行为。对我们的所有 $t^n$ 都有
\begin{equation}
\CA^\sf_n = \frac{1}{\zeta} \CA^{(-1),\sf}_n + \zeta \CA^{(1),\sf}_n
\end{equation}
其中 $\CA^{(\pm 1),\sf}_n$ 是些简单的向量场；然后和前面一样，利用 $\Upsilon$ 当 $\zeta \to 0,\infty$ 时有限这一事实，得到
\begin{equation}
\CA_n = \frac{1}{\zeta} \CA_n^{(-1)} + \CA_n^{(0)} + \zeta \CA_n^{(1)},
\end{equation}
其中
\begin{equation}
\CA_n^{(-1)} = \Upsilon_0 \CA_n^{(-1),\sf} \Upsilon_0^{-1}, \quad \CA_n^{(1)} = \Upsilon_\infty \CA_n^{(1),\sf} \Upsilon_\infty^{-1}.
\end{equation}
再把式 \eqref{eq:zeta-eq-general} 也算进来，我们得到的完整方程组是
\begin{align}
\partial_{u^j} \CX &= \left( \frac{1}{\zeta} {\CA^{(-1)}_{u^j}} + \CA^{(0)}_{u^j} \right) \CX, \label{eq:conn-holomorphy-1} \\
\partial_{\bar{u}^{\bar j}} \CX &= \left( {\CA}^{(0)}_{\bar{u}^{\bar j}} + \zeta {\CA}^{(1)}_{\bar{u}^{\bar j}} \right) \CX, \label{eq:conn-holomorphy-2} \\
\Lambda \partial_\Lambda \CX &= \left( \frac{1}{\zeta} {\CA_\Lambda^{(-1)}} + \CA_\Lambda^{(0)} \right) \CX, \label{eq:lambda-equation} \\
\bar\Lambda \partial_{\bar\Lambda} \CX &= \left( \CA_{\bar \Lambda}^{(0)} + \zeta \CA_{\bar \Lambda}^{(1)} \right) \CX, \label{eq:lambda-bar-equation} \\
R \partial_R \CX &= \left( \frac{1}{\zeta} {\CA_R^{(-1)}} + \CA_R^{(0)} + \zeta \CA_R^{(1)} \right) \CX, \label{eq:R-equation} \\
\zeta \partial_\zeta \CX &= \left( \frac{1}{\zeta} {\CA_\zeta^{(-1)}} + \CA_\zeta^{(0)} + \zeta \CA_\zeta^{(1)} \right) \CX. \label{eq:zeta-equation}
\end{align}
人们还会得到一些额外的关系
\begin{equation} \label{eq:M-relations}
\CA_R^{(-1)} = - \CA_\zeta^{(-1)}, \quad \CA_R^{(1)} = \CA_\zeta^{(1)},
\end{equation}
它们来自这样一件事实：$\CX^\sf$ 在 $\zeta \to 0$ 时被 $\zeta \partial_\zeta + R \partial_R$ 零化，在 $\zeta \to \infty$ 时被 $\zeta \partial_\zeta - R \partial_R$ 零化。

我们的方程到这里就构造完了。附录 D 讨论了在给定 $\CX$ 在 $\zeta = 0$ 附近的渐近展开时，如何把这些方程写得更具体。% TODO 合稿改 \ref（Appendix \ref{app:diff-asymptotics}）
本节最后给出几点评论：

\begin{itemize}
\item
既然辛形式 $\hsymp(\zeta)$ 是由 $\CX$ 构造出来的，式 \eqref{eq:zeta-equation} 与式 \eqref{eq:R-equation} 就平凡地蕴含了 $\hsymp(\zeta)$ 对 $\zeta$ 与 $R$ 的依赖所满足的方程，其形式为
\begin{align}
\zeta \partial_\zeta \hsymp &= \left( \frac{1}{\zeta} \CL_{\CA_\zeta^{(-1)}} +
\CL_{\CA_\zeta^{(0)}} + \zeta \CL_{\CA_\zeta^{(1)}} \right)  \hsymp,
\\ R \partial_R (R\hsymp) &= \left( \frac{1}{\zeta} \CL_{\CA_R^{(-1)}} +
\CL_{\CA_R^{(0)}} + \zeta \CL_{\CA_R^{(1)}} \right) (R \hsymp). \label{eq:omega-R}
\end{align}
回想式 (3.4) 中 $\hsymp(\zeta) = - \frac{i}{2 \zeta}\omega^+ + \omega_3 - \frac{i}{2} \zeta \omega_-$，把这些方程按 $\zeta$ 的幂展开，便可导出超凯勒形式 $\vec\omega$ 的一些有趣的微分方程。% TODO 合稿改 \eqref：式 (3.4) 属前面各块

\item
回想一下：黎曼--希尔伯特问题的解 $\CX$ 带有 $\CX \to b \CX$ 的任意性，其中 $b$ 是 $\CM_u$ 的任一微分同胚。这一任意性导致联络 $\CA$ 的、与 $\zeta$ 无关的规范变换。有几种特别方便的规范。其一是使 $\CA_R^{(0)} = 0$ 的规范。由式 \eqref{eq:omega-R} 可知，在这一规范下，$R \omega_3$ 到每个 $\CM_u$ 上的限制与 $R$ 无关（因而等于其 $R \to \infty$ 极限，即 $- \frac{1}{8 \pi^2} \inprod{d\theta, d\theta}$）。我们的积分方程 式 (5.11) 所选取的规范是否正是这一规范，将会很有意思。% TODO 合稿改 \eqref：式 (5.11) 属前面各块
如果我们允许 $b$ 是\emph{复化}的微分同胚，那么至少在形式上还可以选取使 $\Upsilon_0 = 1$ 的规范，从而 $\CA^{(-1)} = \CA^{(-1),\sf}$；这是 \cite{Cecotti:1993rm} 的「拓扑规范」的类比（对偶地，$\Upsilon_\infty = 1$ 将是「反拓扑规范」）。

\item
我们这些方程的两个线性组合有简单的物理意义：它们表达了理论在整体标度变换与 $R$ 对称变换下的不变性。为此首先注意
\begin{equation}
\left(a^I \partial_{a^I} + \Lambda \partial_\Lambda\right) Z_\gamma = Z_\gamma
\end{equation}
对一切 $\gamma \in \clat$ 成立。由此可得
\begin{align}
\left(R \partial_R - a^I \partial_{a^I} - \bar{a}^I \partial_{\bar{a}^I} - \Lambda \partial_\Lambda - \bar{\Lambda} \partial_{\bar{\Lambda}} \right) \CX^\sf &= 0, \\
\left(\zeta \partial_\zeta  + a^I \partial_{a^I} - \bar{a}^I \partial_{\bar{a}^I} + \Lambda \partial_\Lambda - \bar{\Lambda} \partial_{\bar{\Lambda}} \right) \CX^\sf &= 0.
\end{align}
这些方程可以解释为半平坦几何的（反常）标度不变性与 $R$ 对称不变性。它们在诸 $\CA^\sf_n$ 之间蕴含一些关系（只需作替换 $\partial \to \CA$ 即可），这些关系又给出诸 $\CA_n^{(\pm 1)}$ 之间的关系：我们发现
\begin{align}
\left(R \partial_R - a^I \partial_{a^I} - \bar{a}^I \partial_{\bar{a}^I} - \Lambda \partial_\Lambda - \bar{\Lambda} \partial_{\bar{\Lambda}} \right) \CX &= \Delta \CX, \label{eq:rescaling-1} \\
\left(\zeta \partial_\zeta  + a^I \partial_{a^I} - \bar{a}^I \partial_{\bar{a}^I} + \Lambda \partial_\Lambda - \bar{\Lambda} \partial_{\bar{\Lambda}} \right) \CX &= \Delta' \CX, \label{eq:rescaling-2}
\end{align}
其中 $\Delta$、$\Delta'$ 是 $\CM_u$ 上与 $\zeta$ 无关的向量场。通过规范变换可以令 $\Delta = 0$、$\Delta' = 0$。确实，我们的积分方程自动选取了合适的规范：附录 C 给出的 $R$ 很大时的递推解，逐项满足式 \eqref{eq:rescaling-1} 与式 \eqref{eq:rescaling-2}，且 $\Delta = \Delta' = 0$。% TODO 合稿改 \ref（Appendix \ref{app:integral-asymptotics}）
所以，在某种意义下，标度不变性与 $R$ 对称不变性在瞬子修正之后仍然存活。

\item
式 \eqref{eq:R-equation} 与式 \eqref{eq:zeta-equation} 之间的相容性，连同关系 \eqref{eq:M-relations}，蕴含四元组 $(\CA_\zeta^{(\pm 1)}, \CA_\zeta^{(0)}, \CA_R^{(0)})$ 对 $R$ 的依赖满足一组非线性微分方程。正如附录 D 所解释的，这些方程是 Nahm 方程的一个形变；% TODO 合稿改 \ref（Appendix \ref{app:diff-asymptotics}）
该四元组的大 $R$ 展开，可以直接通过迭代求解这些方程得到。它们是研究我们的构造在\emph{小} $R$ 处行为的一种可能工具。Nahm 方程的一般解会在某个有限的 $R$ 处变奇异，我们预期我们的问题也是如此。不过我们预期，这里所描述的这些特解——它们由 $\CN=2$、$d=4$ 场论中的 BPS 简并度所决定——实际上对一切 $R$ 都\emph{是}正则的。可能这会对 $\CN=2$ 理论可能的 BPS 谱与红外预势给出一个有趣的限制。\cite{Cecotti:1993rm} 中曾用极为相似的策略来限制 $d=2$ 理论的性质。

\item
本节的讨论为穿墙公式的作用提供了一个新的视角。方程组 \eqref{eq:conn-holomorphy-1}–\eqref{eq:zeta-equation} 描述了 $\IC\IP^1 \times \CP$ 上的一个平坦联络，其中 $\CP$ 是以诸 $t^n$ 为坐标的参数空间。这个平坦联络可以等价地看作 $\IC\IP^1$ 上联络的一个\emph{等单值形变族}，在 $\zeta=0, \infty$ 处有秩 $1$ 的非正则奇点。在每个 $t \in \CP$ 处，$\IC\IP^1$ 上该联络的斯托克斯数据由 KS 因子给出。利用沿 $\CP$ 的平行移动可以证明：非正则奇点处的斯托克斯数据在适当的意义下是「不变的」。确切地说：在 $\zeta$ 平面中任取一个凸扇形 $\CV$，则乘积
\begin{equation}
A_\CV = \prod_{\ell \subset \CV}^\ccwarrow S_\ell
\end{equation}
在使任何斯托克斯线 $\ell$ 都不进入或离开 $\CV$ 的任何 $t \in \CP$ 的变动下保持不变。把这一论断用到 $u$ 的变动上，就重新得到穿墙公式。
\end{itemize}

\subsection{度量的构造及其大 $R$ 渐近}
\label{sec:metric-construction}

至此我们已经构造了 $\CM$ 上的一族函数 $\CX_\gamma(\zeta)$，以及相应的全纯辛形式 $\hsymp(\zeta)$。正如第 3.2 节所讨论的，给定一个满足第 \ref{sec:symplectic-form} 节所列性质的 $\hsymp(\zeta)$，以及满足形如式 \eqref{eq:conn-holomorphy-1}、\eqref{eq:conn-holomorphy-2} 的「柯西--黎曼」方程的 $\CX_\gamma$，就存在 $\CM$ 上相应的超凯勒度量 $g$。% TODO 合稿改 \ref（第 3.2 节 = Section \ref{sec:twistor-general}）
这就是我们对 $g$ 的构造。

给定解出式 (5.13) 的精确函数 $\CX_\gamma(\zeta)$，我们可以把 $g$ 写成如下的封闭形式。% TODO 合稿改 \eqref：式 (5.13) 属前面各块
利用式 (5.13) 中核的展开（当 $|\zeta'/\zeta| < 1$ 时），得到 $\zeta \to 0$ 的渐近展开，% TODO 合稿改 \eqref：式 (5.13) 属前面各块
\begin{equation}
\log \CX_{\gamma} = \frac{1}{\zeta} F_{-1}^\gamma + F_0^\gamma +
\zeta F_1^\gamma + \CO(\zeta^2),
\end{equation}
其中
\begin{equation}
\begin{split}
F_{-1}^{\gamma} & = \pi R Z_\gamma, \\
F_0^\gamma & = i \theta_\gamma -  \frac{1}{4 \pi i} \sum_{\gamma'}
\Omega(\gamma';u) \langle \gamma,\gamma' \rangle
\int_{\ell_{\gamma'}} \frac{d\zeta'}{\zeta'}   \log (1-
\sigma(\gamma') \CX_{\gamma'}(\zeta')),\\
F_1^\gamma & = \pi R \bar Z_\gamma -  \frac{1}{2 \pi i}
\sum_{\gamma'} \Omega(\gamma';u) \langle \gamma,\gamma' \rangle
\int_{\ell_{\gamma'}} \frac{d\zeta'}{\zeta'^2}   \log (1-
\sigma(\gamma') \CX_{\gamma'}(\zeta')).
\end{split}
\end{equation}
然后代入 $\varpi$，即可提取出
\begin{equation}
\begin{split}
\omega_+ & = \frac{i}{2\pi^2 R} \epsilon_{ij} dF_{-1}^{\gamma^i}
\wedge d F_0^{\gamma^j}, \\
\omega_3 & = \frac{1}{8\pi^2 R} \epsilon_{ij} \left( 2
dF_1^{\gamma^i}\wedge d F_{-1}^{\gamma^j} + d F_0^{\gamma^i} \wedge
d
F_0^{\gamma^j} \right).\\
\end{split}
\end{equation}
由这些辛形式很容易得到 $g$。

现在来考虑 $g$ 在大 $R$ 处的行为。第 5.3 节讨论了诸 $\CX_\gamma$ 的大 $R$ 渐近，其中包括式 (5.14) 给出的首个 BPS 瞬子修正。% TODO 合稿改 \ref/\eqref：第 5.3 节 = Section \ref{sec:solving-rh}；式 (5.14) = \eqref{eq:X-first-approx}
现在把它翻译成对 $\hsymp(\zeta)$ 的修正。我们先算出 $\frac{d\CX_\gamma}{\CX_\gamma}$ 的修正：
\begin{equation} \label{eq:dX-correction}
\frac{d\CX_\gamma}{\CX_\gamma} = \frac{d\CX^\sf_\gamma}{\CX^\sf_\gamma} + \CI_\gamma + \cdots,
\end{equation}
其中
\begin{equation} \label{eq:correction-integral}
\CI_\gamma =  \frac{1}{4\pi i} \sum_{\gamma'} \degen(\gamma';u)
\langle \gamma,\gamma' \rangle \int_{\ell_{\gamma'}}
\frac{d\zeta'}{\zeta'} \frac{\zeta' + \zeta}{\zeta' - \zeta}
\frac{d\CX^\sf_{\gamma'}(\zeta')}{\CX^\sf_{\gamma'}(\zeta')}
\frac{\sigma(\gamma') \CX^\sf_{\gamma'}(\zeta')}{1- \sigma(\gamma')
\CX^\sf_{\gamma'}(\zeta')}.
\end{equation}
注意，正如所承诺的，$\CI_\gamma$ 当 $R \to \infty$ 时指数压低，因为在 $\ell$ 上，当 $R \to \infty$ 时 $\CX_{\gamma'}^\sf$ 指数地趋于零。式 \eqref{eq:dX-correction} 中的省略号表示多瞬子修正，它们被压得更低。因此 $\hsymp(\zeta)$ 的首个修正来自 $\frac{d\CX^\sf_\gamma}{\CX^\sf_\gamma}$ 与 $\CI_\gamma$ 之间的楔积。

为更显式地描述这一修正，方便的做法是逐个处理每个 $\gamma'$，并取一个辛基 $\{\gamma^1, \dots, \gamma^{2r}\}$，使 $\gamma' = q_{\gamma'} \gamma^1$。于是式 \eqref{eq:correction-integral} 中的积分与积分 式 (4.38) 本质上相同——后者在第 4.3 节给出了 $\CX_\magn$ 的瞬子修正。% TODO 合稿改 \eqref/\ref：式 (4.38) = \eqref{eq:ipm}；第 4.3 节 = Section \ref{sec:solving-xmagn}
像那里一样算出 $\hsymp(\zeta)$ 相应的修正，便得到
\begin{equation}
\hsymp(\zeta) = \hsymp^\sf(\zeta) + \sum_{\gamma' \in \clat} \hsymp^\inst_{\gamma'}(\zeta) + \cdots,
\end{equation}
其中
\begin{equation} \label{eq:Omega-inst-general}
\hsymp^\inst_{\gamma'}(\zeta) = - \degen(\gamma';u) \frac{1}{4 \pi^2 R}\frac{d \CX^\sf_{\gamma'}(\zeta)}{\CX^\sf_{\gamma'}(\zeta)} \left[ A^\inst_{\gamma'} +
\frac{1}{2} V^\inst_{\gamma'} \left(\frac{1}{\zeta} da_{\gamma'} - \zeta d\bar{a}_{\gamma'}\right)\right],
\end{equation}
而（对照式 (4.7)、式 (4.10)）% TODO 合稿改 \eqref：式 (4.7) = \eqref{eq:V-inst}，式 (4.10) = \eqref{eq:A-inst}
\begin{align} \label{eq:A-inst-general}
V^\inst_{\gamma'} &= \frac{R q_{\gamma'}^2}{2\pi} \sum_{n>0}
\sigma(n \gamma') e^{i n \theta_{\gamma'}} K_0(2 \pi R \abs{n
Z_{\gamma'}}), \\
A^\inst_{\gamma'} &= - \frac{R q_{\gamma'}^2}{4 \pi} \left( \frac{dZ_{\gamma'}}{Z_{\gamma'}} - \frac{d \bar Z_{\gamma'}}{\bar Z_{\gamma'}} \right) \sum_{n>0} \sigma(n \gamma') e^{i n \theta_{\gamma'}} \abs{Z_{\gamma'}} K_1(2 \pi R \abs{n Z_{\gamma'}}).
\end{align}
从这里出发，可以按 $\zeta$ 展开，以提取 $\omega_+$、$\omega_3$ 的首阶修正，从而得到 $g$ 的首阶修正。

\subsection{与物理度量的比较}
\label{sec:comparison}

我们已经在 $R$ 充分大时于 $\CM$ 上构造出超凯勒度量 $g$，现在把它的若干性质总结如下：

\begin{enumerate}
\item $g$ 连续；
\item 若一切 BPS 粒子都满足 $|Z| \to \infty$，则 $g$ 趋于半平坦度量 $g^\sf$；
\item $g$ 光滑，只有那些物理上预期的奇点例外——它们位于 $\vm$ 中诸奇异轨迹的上方；
\item $g$ 满足 $\vol(\CM_u) = \left(\frac{1}{R}\right)^{r}$；
\item 在复结构 $J_3$ 下，$(\CM, g)$ 可以等同于处于标准复结构的 SW 环面纤维化；在这一等同之下，全纯辛形式为 $\omega_+ = - \frac{1}{4 \pi} \inprod{dZ, d\theta}$。
\end{enumerate}

所有这些性质都与 $\IR^3 \times S^1$ 上物理度量的预期相符，后者在 \cite{Seiberg:1996nz} 中有描述。因此，最简单的自洽图像就是：我们构造的度量确实就是物理度量。（在秩 $1$ 情形，\cite{Seiberg:1996nz} 曾提出，借助丘成桐（Yau）定理的一个非紧类比，这些性质确实\emph{决定}了度量。更一般地，存在类似定理的可能性是相当大的。）

本节的余下部分将从我们的构造出发论证这些性质：
\begin{enumerate}
\item $g$ 的连续性由穿墙公式推出，这一点我们已经解释过。

\item 只需一看修正项的形式 \eqref{eq:Omega-inst-general} 即可：它们全都被 $R \abs{Z_{\gamma'}}$ 指数压低，因而当一切 $\abs{Z_{\gamma'}} \to \infty$ 时指数快地消失。

\item 在任何使 $R \abs{Z_\gamma} \to \infty$（对一切 $\gamma$）的极限下，瞬子贡献都指数压低，$g$ 趋于 $g^\sf$。这足以确立 $g$ 在大 $R$ 处的光滑性，\emph{除了}在某些奇异轨迹附近——那里有若干电荷为 $\gamma_i$ 的 BPS 粒子变成无质量（$Z_{\gamma_i} = 0$ 且 $\degen(\gamma_i;u) \neq 0$）。为理解这些点附近的行为，我们考虑一个标度极限：$R \to \infty$ 而保持 $R Z_{\gamma_i}$ 有限。在这一极限下，可以用一个只保留 BPS 射线 $\ell_{\gamma_i}$、丢掉其余一切射线的黎曼--希尔伯特问题来近似原来的问题。确实，其余一切不连续量都含有 $(1-\sigma(\gamma')\CX_{\gamma'})$ 形式的因子，而在这个标度极限下它们指数地趋近于 $1$。

在最简单的情形，即只有一个 $Z_\gamma = 0$ 时，我们总可以选取一个对偶标架，使 $\gamma$ 为电电荷。通过把某些角 $\theta$ 平移 $\pi$，还可以安排细化 $\sigma$ 在该标架下取标准形式 $\sigma = (-1)^{\gamma_e \cdot \gamma_m}$。于是我们就处在第 4 节所研究的情形中，那里找到了一个超凯勒度量，它除了一个周期性的 $q$ 个 $A_{q-1}$ 奇点阵列之外是光滑的。% TODO 合稿改 \ref（第 4 节 = Section \ref{sec:periodicnut}）
这与 $d=3$ 有效场论的预期一致：奇点出现在电荷为 $q$ 的超多重态 Kaluza--Klein 塔中某一个变成无质量的地方。

除了我们考察过的这些物理奇点——即一组相互局域的 BPS 粒子变成无质量之处——之外，还可能有超共形点，那里相互非局域的粒子同时变成无质量 \cite{Argyres:1995jj,Argyres:1995xn}。我们尚未分析这些奇点；尽管我们预期它们会很有意思，并且预期二次加细在其分析中将起重要作用。

\item 由于 $\CM_u$ 关于 $J_3$ 是一个复环面，其体积就是 $\frac{1}{r!} \int_{\CM_u} \omega_3^r$。另一方面，利用式 (3.4)，并注意到由式 \eqref{eq:omega-residues} 知 $\omega_{\pm}$ 在 $\CM_u$ 上的限制为零，便有% TODO 合稿改 \eqref：式 (3.4) 属前面各块
\begin{equation}
\vol(\CM_u) = \frac{1}{r!} \int_{\CM_u} \hsymp^r(\zeta) = \frac{1}{(4 \pi^2 R)^r r!} \int_{\CX(\CM_u)} (\hsymp^T)^r = \left(\frac{1}{R}\right)^r,
\end{equation}
正如所愿。

\item 复结构 $J_3$ 可以由 $\omega_1$ 与 $\omega_2$ 定出，只需 $J_3 = \omega_1^{-1} \omega_2$。但这一信息又由 $\hsymp(\zeta)$ 在 $\zeta = 0$ 处的留数给出；回想式 (3.4)，$\omega_+ = \omega_1 + i \omega_2$ 正是由% TODO 合稿改 \eqref：式 (3.4) 属前面各块
\begin{equation}
\omega_+ = 2i\,\Res_{\zeta = 0} \, \hsymp(\zeta).
\end{equation}
给出的。我们关于 $\hsymp(\zeta)$ 的渐近条件 式 (5.3) 恰好保证它与 $\hsymp^\sf(\zeta)$ 的留数相关：确实，我们正好有% TODO 合稿改 \eqref：式 (5.3) 属前面各块
\begin{equation}
\omega_+ = \Upsilon_0^* \omega_+^\sf.
\end{equation}
由此可知，只需作用以逐纤维微分同胚 $\Upsilon_0$，就能把 $(\CM, J_3)$ 等同于 $(\CM, J_3^\sf)$。正如第 3.3 节所解释的，$\CM$ 上的复结构 $J_3^\sf$ 正是 SW 环面纤维化的复结构。% TODO 合稿改 \ref（第 3.3 节 = Section \ref{sec:semiflat-twistor}）
此外，在这一等同之下，$\omega_+$ 被等同于式 (3.16) 给出的 $\omega_+^\sf$。% TODO 合稿改 \eqref：式 (3.16) 属前面各块
\end{enumerate}

% ch06 —— §6 Adding masses（原文 wcf-paper.tex 3002–3152 行）
\section{加入质量}
\label{sec:masses}

本节我们简要说明：若要纳入非平凡的质量参数，前几节的结果应当如何修改。

\subsection{含质量的单粒子修正}

第 4 节所考虑的 $U(1)$ 理论有一个简单的变体：我们可以考虑含有若干超多重态的 $U(1)$ 理论，这些超多重态带电，电荷为 $q_i$。% TODO 合稿改 \ref{sec:periodicnut}
粒子种类多于一种的理论会牵涉味荷，并且会非平凡地依赖于质量参数。四维中的质量参数是复数 $m_i$。紧致化到三维后，会多出一个实的周期质量参数 $m^3_i$，它本质上是味对称的一条 Wilson 线。我们记 $\psi_i := 2\pi R m^3_i$，其周期为 $2\pi$。

质量参数以非常简单的方式进入修正后的度量：每个粒子对 $V$ 与 $A$ 给出一个可加贡献，其形式与我们此前遇到的那个相仿，
\begin{equation} \label{eq:V-sum-masses}
V = \sum_i \frac{q_i^2 R}{4 \pi} \sum_{n=-\infty}^\infty \left(
\frac{1}{\sqrt{R^2 \abs{q_i a + m}^2 + (q \frac{\theta_\elec}{2 \pi} +
\frac{\psi_i}{2 \pi} + n)^2}} - \kappa_n \right)
\end{equation}

坐标 $\CX_\elec$ 不变：
\begin{equation}
\CX_\elec = \exp \left[ \pi R \frac{a}{\zeta} + i \theta_\elec +
\pi R \zeta \bar{a} \right].
\end{equation}
引入质量参数的类似组合也很有用：
\begin{equation}
\mu_i := \exp \left[ \pi R \frac{m_i}{\zeta} + i \psi_i + \pi R
\zeta \bar{m_i} \right].
\end{equation}
半平坦的 $\CX_\magn$ 含有把 $d=4$ 中所有粒子积分掉所产生的贡献：
\begin{equation}
\begin{aligned}
\CX_\magn^{\sfm}(\zeta) &= e^{i \theta_\magn} \times \\
&\quad \prod_i \exp\left[-i\frac{R q_i}{2
\zeta} \left((q_i a+m_i) \log\frac{q_i a+ m_i}{e \Lambda} \right) +
i \frac{\zeta R q_i}{2} \left((q_i \bar a + \bar
m_i)\log\frac{q_i \bar a+ \bar m_i}{e\bar \Lambda} \right)\right].
\end{aligned}
\end{equation}
$\CX_\magn^{\sfm}(\zeta)$ 绕 $q_i a+m_i =0$ 的单值变换为
\begin{equation}\label{eq:mono2-masses}
\CX_\magn^{\sfm} \to (-\mu_i)^{q_i} \CX_\elec^{q_i^2}
\CX_\magn^{\sfm}.
\end{equation}
完整的坐标同样含有来自所有这些粒子的瞬子贡献，
\begin{equation}\label{eq:FullX2-masses}
\begin{aligned}
\CX_\magn = \CX_\magn^{\sfm} \prod_i \exp\Biggl[ & \frac{i q_i}{4\pi}
\int_{\ell^i_+} \frac{d \zeta'}{\zeta'}
\frac{\zeta' + \zeta}{\zeta' - \zeta} \log[1-\mu_i \CX_\elec(\zeta')^{q_i}] \\
& - \frac{i q_i}{4\pi}\int_{\ell^i_-}
\frac{d \zeta'}{\zeta'} \frac{\zeta' + \zeta}{\zeta' - \zeta}
\log[1-\mu_i^{-1} \CX_\elec(\zeta')^{-q_i}] \Biggr],
\end{aligned}
\end{equation}
其中围道 $\ell^i_\pm$ 取为 $\zeta$ 平面中连接 $0$ 与 $\infty$、且落在下面两个半平面内的任意路径：
\begin{equation}
\CU^i_\pm = \left\{ \zeta: \pm {\rm Re} \frac{q_i a + m_i}{\zeta} < 0 \right\}.
\end{equation}

不连续性如今依赖于质量：
\begin{subequations} \label{eq:X2-disc-masses}
\begin{align}
(\CX_\magn)_{\ell^i_+}^{+} &= (\CX_\magn)_{\ell^i_+}^{-} (1 - \mu_i \CX_\elec^q)^{-q}, \\
(\CX_\magn)_{\ell^i_-}^{+} &= (\CX_\magn)_{\ell^i_-}^{-} (1 - \mu_i^{-1}
\CX_\elec^{-q})^q.
\end{align}
\end{subequations}

本节的所有公式也都可以沿着第 4.7 节的思路推广到更高秩。% TODO 合稿改 \ref{sec:higher-rank}

\subsection{含质量的多粒子修正}

现在我们已有条件理解质量参数在一般的黎曼--希尔伯特问题与微分问题中所起的作用。考虑一个具有 $n_f$ 个味对称性的规范理论。把味荷记为 $\gamma^f$，并由质量与味 Wilson 线以显然的方式构造出 $\mu_{\gamma^f}$：
\begin{equation}
\mu_{\gamma^f} := \exp \left[ \pi R \frac{m_{\gamma^f}}{\zeta} + i \psi_{\gamma^f} + \pi R
\zeta \bar{m}_{\gamma^f} \right].
\end{equation}
阿贝尔问题的不连续性提示出如下形式的广义 Kontsevich--Soibelman 因子：
\begin{equation}
\label{eq:KS-generalized-masses}
\CK_{\gamma,\gamma^f} := \CX_{\gamma'} \to \CX_{\gamma'} (1-
\sigma(\gamma) \mu_{\gamma^f} \CX_\gamma)^{\langle \gamma,\gamma'
\rangle}.
\end{equation}
于是我们可以定义一个与第 5.1 节相仿的黎曼--希尔伯特问题，它对每个电荷为 $\gamma$、味荷为 $\gamma^f$ 的粒子配上不连续性 $\CK_{\gamma,\gamma^f}$。% TODO 合稿改 \ref{sec:defining-rh}
假定推广到使用这些修改后的辛同胚之后，穿墙公式仍然给出正确的 BPS 简并度，我们就可以用这个黎曼--希尔伯特问题在 $\CM$ 上构造一个超凯勒度量；我们提出，它就是正确的那个度量。

超对称场论中有一个标准技巧：把质量参数视为某个扩充理论的矢量多重态标量的真空期望值，而在该扩充理论中，味对称性已被弱规范化 \cite{Seiberg:1993vc}。由于是弱规范化，带磁味荷的粒子非常重，在味规范耦合趋于零的极限下可以略去；此时剩下的味规范荷只有电性的。利用这个技巧，含质量的广义穿墙公式可以解释为标准穿墙公式在零耦合极限下的情形。

最后，我们希望把第 5.5 节的微分表述推广到能处理质量参数。% TODO 合稿改 \ref{sec:differential-problem}
在那一节中，我们依赖于诸因子 $S_\ell$ 与参数无关这一事实。在修改后的问题中，诸 $S_\ell$ 显式依赖于 $\mu_{\gamma^f}$，从而依赖于 $m_i$、$R$、$\zeta$。然而，所有 $S_\ell$ 确实都被 $\partial_{m_i} + i \frac{\pi R}{\zeta} \partial_{\psi_i}$ 零化，也被 $\partial_{\bar m_i} + i \pi R \zeta \partial_{\psi_i}$ 零化。于是把那一节的论证稍作修改，即可说明：黎曼--希尔伯特问题的解 $\CX$ 满足如下形式的微分方程
\begin{align}
\partial_{m_i} \CX &= \left(\frac{1}{\zeta} \CA^{(-1)}_{m_i} + \CA^{(0)}_{m_i} \right) \CX,
\label{eq:mass-equation} \\
\partial_{\bar{m}_i} \CX &= \left(
\CA^{(0)}_{\bar{m}_i} + \zeta \CA^{(1)}_{\bar{m}_i} \right) \CX, \label{eq:mass-bar-equation}
\end{align}
其中 $\CA^{(-1)}_{m_i}$ 不仅是 $\CM_u$ 上的矢量场，还包含算子 $-i \frac{\pi R}{\zeta} \partial_{\psi_i}$；对 $\CA^{(1)}_{\bar{m}_i}$ 也有类似的说法。

诸 $S_\ell$ 也被 R 对称性与标度不变性的算子零化，因此我们得到 \eqref{eq:rescaling-m1}、\eqref{eq:rescaling-m2} 的类似物（在转到适当的规范之后），
\begin{align}
\left(R \partial_R - a^I \partial_{a^I} - \bar{a}^I \partial_{\bar{a}^I} - \Lambda \partial_\Lambda - \bar{\Lambda} \partial_{\bar{\Lambda}} - m^i \partial_{m^i} - \bar{m}^i \partial_{\bar{m}^i} \right) \CX &= 0, \label{eq:rescaling-m1} \\
\left(\zeta \partial_\zeta + a^I \partial_{a^I} - \bar{a}^I \partial_{\bar{a}^I} + \Lambda \partial_\Lambda - \bar{\Lambda} \partial_{\bar{\Lambda}} + m^i \partial_{m^i} - \bar{m}^i \partial_{\bar{m}^i} \right) \CX &= 0. \label{eq:rescaling-m2}
\end{align}
利用这些方程，我们可以得到 $\CX$ 对 $R$ 与 $\zeta$ 的依赖所满足的标准形式 (5.34)、(5.35)，同样带有这样的修改：$\CA_R$ 与 $\CA_\zeta$ 现在包含对诸 $\psi_i$ 的导数。% TODO 合稿改 \eqref{eq:R-equation}、\eqref{eq:zeta-equation}

\begin{yizhu}
\begin{enumerate}[leftmargin=2em,itemsep=2pt,topsep=2pt]
\item 式 \eqref{eq:V-sum-masses} 中的 $\abs{q_i a + m}$ 与 $q\,\theta_\elec/(2\pi)$，原文分别印作 $m$ 与 $q$（不带下标 $i$），而下文同类式子一律带下标；译文公式照原文逐字保留，此处仅提示读者：这两处应按 $m_i$、$q_i$ 理解。
\item 式 \eqref{eq:KS-generalized-masses} 中的 $\sigma(\gamma)\in\IZ_2$ 是第 2.2 节引入的二次加细（quadratic refinement），满足 $\sigma(\gamma_1)\sigma(\gamma_2)=(-1)^{\langle\gamma_1,\gamma_2\rangle}\sigma(\gamma_1+\gamma_2)$；本块沿用它，不另作定义。% TODO 合稿改 \ref{sec:ReviewKS}
\end{enumerate}
\end{yizhu}

% ch07 —— §7 穿墙公式的一个证明 + 致谢（原文 wcf-paper.tex 3153–3239 行，PDF 第 50–51 页）
\section{穿墙公式的一个证明}
\label{sec:proof}

本文中，我们在 $\CM$ 上构造了一个 \hk 度量 $g$，并论证它与 $\IR^3\times S^1$ 上该规范理论的模空间所带的物理度量相符。Kontsevich--Soibelman 穿墙公式是作为一条自洽性条件出现的：没有它，我们的构造就无法给出光滑的度量。我们把这看作穿墙公式确实成立的有力旁证。

然而，这些构造还称不上是对穿墙公式的一个证明。要给出证明，我们必须直接从规范理论的物理出发，而不对度量应当采取何种形式作任何假设。在这种进路下，我们无法动用黎曼--希尔伯特构造这一工具（至少一开始不能）。我们改用第 5.5 节所描述的另一种视角。设 $\CM$ 是该规范理论在 $\IR^3\times S^1$ 上的模空间，并考虑映射 $\CX(\zeta)\colon \CM\to T$（$\zeta\neq0,\infty$）。我们的目标是为此类 $\CX(\zeta)$ 构造一组可积方程，其形式如 (5.30)–(5.35)，使得 $\IC\IP^1$ 上的联络 (5.35) 具有斯托克斯射线 $\ell$，其上承载斯托克斯因子
\begin{equation} \label{eq:sf-redux}
S_\ell = \prod_{\gamma \in (\clat_u)_\ell} \CK_\gamma^{\degen(\gamma;u)}.
\end{equation}
一旦构造出这样的微分方程，WCF 便是说：联络 (5.35) 经历等单值形变。% TODO 合稿改 \ref{sec:differential-problem}、\eqref{eq:conn-holomorphy-1}–\eqref{eq:zeta-equation}

下面我们说明如何直接从规范理论导出这些微分方程。正如我们在附录 B 中所说明的，(5.30)、(5.31) 有一个简单的几何意义：它们正是柯西--黎曼方程，表达 $\CX(\zeta)$ 关于复结构 $J^\pz$ 的全纯性。特别地，这些方程可以纯粹用约化理论的 $\CN=4$ 超对称来理解。其次注意到，(5.32)、(5.33) 与 (5.30)、(5.31) 的形式完全相同。事实上，只要把 $\Lambda$ 看作某个「背景」矢量多重态的标量分量，它们就会变得完全一样。这是证明非重整化定理的一个标准技巧，例如见 \cite{Seiberg:1993vc}；把它用在这里，应当就能得到所需的 (5.32)、(5.33)。若理论含有质量参数，我们还可以通过弱规范化味对称性，类似地证明 (6.11)、(6.12)。% TODO 合稿改 \ref{app:holomorphy}、\eqref{eq:conn-holomorphy-1}、\eqref{eq:conn-holomorphy-2}、\eqref{eq:lambda-equation}、\eqref{eq:lambda-bar-equation}、\eqref{eq:mass-equation}、\eqref{eq:mass-bar-equation}

最后，我们需要建立关键方程 (5.34)、(5.35)，它们给出 $\CX$ 对 $R$ 与 $\zeta$ 的依赖。这两个方程可从 $d=4$ 理论的反常 $U(1)_R$ 对称性与标度不变性推出，而这些对称性由 (5.42)、(5.43) 或 (6.13)、(6.14) 表达；把上述对称性与上文已经建立的方程结合起来，即可得到它们。诸函数 $\CX_\gamma$ 有一种物理解释，我们希望另文描述。可以把它们看作某个三维拓扑场论的手征环的元素，或看作四维理论中某些线算子的期望值。按这种眼光，(5.34)、(5.35) 就是反常 Ward 恒等式。% TODO 合稿改 \eqref{eq:R-equation}、\eqref{eq:zeta-equation}、\eqref{eq:rescaling-1}、\eqref{eq:rescaling-2}、\eqref{eq:rescaling-m1}、\eqref{eq:rescaling-m2}

为完成证明，我们还须说明：$\IC\IP^1$ 上该联络的斯托克斯因子确实由 \eqref{eq:sf-redux} 给出。为此我们利用这样一件事实：由 (5.34)，斯托克斯因子在 $R$ 变化下保持不变。于是我们可以取非常大的 $R$；在那里（远离诸墙处），度量的修正应当可以很好地近似为若干单瞬子修正的线性叠加，而这些单瞬子修正是我们在阿贝尔理论中已经知道的。从 $\IC\IP^1$ 上的联络转到相应的黎曼--希尔伯特问题，并重复我们在第 5 节用过的论证，便可说明这些修正直接对应于斯托克斯因子。至此，穿墙公式的证明完成——至少是在物理的严格性水平上。% TODO 合稿改 \eqref{eq:R-equation}、\ref{sec:constructing-moduli}

\section*{致谢}

我们要特别感谢 F. Denef、N. Nekrasov 与 X. Yin 在本项目早期阶段的合作，并感谢 B. Pioline 关于半平坦度量的重要讨论。我们也感谢 Wu-yen Chuang、E. Diaconescu、D. Jafferis、D. Joyce、M. Kontsevich、S. Lukyanov、N. Seiberg、Y. Soibelman、V. Toledano Laredo、C. Vafa、E. Witten 与 E. Zaslow 的有益讨论。GM 的工作受美国能源部（DOE）资助，资助号 DE-FG02-96ER40949。GM 还感谢阿斯彭物理中心（Aspen Center for Physics）与加州大学圣巴巴拉分校（UCSB）的卡弗里理论物理研究所（KITP）在本工作完成期间的热情接待；KITP 的部分经费来自美国国家科学基金会（NSF）的资助 PHY05-51164。AN 的工作部分由高等研究院（Institute for Advanced Study）的 Martin A. and Helen Chooljian 会员资格资助，部分由 NSF 在资助号 PHY-0503584 下资助。DG 的工作部分由 DOE 资助号 DE-FG02-90ER40542 资助。

\begin{yizhu}
本节中的 $T$ 指 $\CM$ 的环面纤维（$\CM$ 在拓扑上是 4d 模空间 $\vm$ 上的 \SW 环面丛），映射 $\CX(\zeta)$ 取值于其中；$\CK_\gamma$ 是第 2.2 节引入的 Kontsevich--Soibelman 辛同胚，$(\clat_u)_\ell$ 则是中心荷相位落在射线 $\ell$ 上的那些电荷所构成的子集。% TODO 合稿改 \ref{sec:ReviewKS}
\end{yizhu}


\appendix
% appA —— 附录 A（原文 wcf-paper.tex 3242–3392 行，PDF 51–55 页）
\section{验证若干 $SU(2)$ 规范理论的 KS 恒等式}
\label{app:ks-product}

有一个富有启发性的办法，可以证明涉及 $\CK_{1,0}$ 与 $\CK_{0,1}$ 的那个简单公式。考虑一列数 $x_n$，它们满足递推关系
\begin{equation} \label{eq:appA-five-recursion}
x_{n+1} x_{n-1} = 1 - x_n.
\end{equation}
令人惊讶的是，这个递推关系是周期的，周期为五：
\begin{equation} \label{eq:appA-five-period}
x_2 = \frac{1-x_1}{x_0}, \quad x_3 = \frac{x_0+ x_1-1}{x_0 x_1}, \quad x_4 = \frac{1-x_0}{x_1}, \quad x_5 = x_0, \quad x_6=x_1.
\end{equation}
现在令 $X_{1,0}=X_{1,0}^{(1)}= x_1^{-1}$，$X_{0,1}=X_{0,1}^{(1)}= x_0$。我们的策略是：取一列适当的 KS 变换 $\CK_n$，依次定义变换 $(X_{1,0}^{(n+1)},X_{0,1}^{(n+1)}) = \CK_n (X_{1,0}^{(n)},X_{0,1}^{(n)})$，直到在 $(X_{1,0}^{(1)},X_{0,1}^{(1)})\to (X_{1,0}^{(N)},X_{0,1}^{(N)})$ 处得到恒等变换为止（在第一个例子中 $N=6$，但在其余例子中将会是无穷大）。为避免记号累赘，下文中我们不再标出上标 ${}^{(n)}$。

若施加 $\CK_{1,0}$，它不改变 $X_{1,0} = x^{-1}_1$ 的值，而是把 $X_{0,1}$ 变为 $X_{0,1}=x_0(1-x^{-1}_1)^{-1} =- x_1 x^{-1}_2$。注意 $X_{1,1} = X_{1,0} X_{0,1} = -x_2^{-1}$。

接着可以施加 $\CK_{1,1}$：它使 $X_{1,1} = -x^{-1}_2$ 保持不变，而改变 $X_{1,0}$。结果现在 $X_{0,1}=- x_1 x^{-1}_2 (1-x^{-1}_2)^{-1}= x_3^{-1}$。

若施加 $\CK_{0,1}$，则 $X_{0,1} = x_3^{-1}$，而 $X_{1,0} = - x_3 x^{-1}_2 (1-x^{-1}_3)=x_4$。若施加 $\CK_{1,0}^{-1}$，则 $X_{1,0} = x_4$，$X_{0,1} =x^{-1}_3 (1-x_4) = x_5$。最后，若施加 $\CK_{0,1}^{-1}$，我们得到 $X_{0,1}=x_5=x_0$，$X_{1,0} = x_4 (1-x_5)^{-1} = x^{-1}_6=x^{-1}_1$。于是便导出了所要的
\begin{equation} \label{eq:appA-pentagon}
\CK_{0,1}^{-1}\CK_{1,0}^{-1}\CK_{0,1}\CK_{1,1}\CK_{1,0} = 1.
\end{equation}

对更有意思的公式而言，这是一个有用的热身练习。现在考虑另一个递推关系：
\begin{equation} \label{eq:appA-sq-recursion}
x_{n+1} x_{n-1} = (1- x_n)^2.
\end{equation}
这个递推关系一般不是周期的：它的通解为
\begin{equation} \label{eq:appA-sq-solution}
x_n = -\frac{\cosh^2 (a n +b)}{\sinh^2 a}.
\end{equation}
我们仍可把这个递推关系与 $\CK$ 因子的一个乘积联系起来。

我们仍从 $X_{0,1}=x_0$ 与 $X_{1,0} = x^{-1}_1$ 出发。若施加 $\CK_{1,0}^2$，结果是 $X_{1,0}=x^{-1}_1$，$X_{2,1} = x_1^{-2} x_0 (1- x^{-1}_1)^{-2}=x^{-1}_2$。若施加 $\CK_{2,1}^2$，结果是 $X_{2,1}=x^{-1}_2$，$X_{3,2}=x_2^{-2} x_1 (1-x^{-1}_2)^{-2} = x^{-1}_3$。我们可以对一切 $n$ 继续施加 $\CK_{n+1,n}^2$，循着递推关系走到任意大的 $n$。借助关系 $X_{n+1,n} = x^{-1}_{n+1}$ 与 $X_{n,n-1} = x^{-1}_{n}$ 在 $n\to\infty$ 时的极限，便可算出这个无穷乘积。若取 $a,b$ 的实部为正，则 $X_{1,1} = e^{-2 a}$，且 $X_{1,0} =- e^{-2 b} (1-e^{-2 a})^2$。

另一方面，我们也可以沿相反方向追随这个递推关系：$X_{0,1}=x_0$ 与 $X_{1,0} =x^{-1}_1$ 在 $\CK_{0,1}^2$ 之下变为 $X_{0,1}=x_0$ 与 $X_{1,0} = x^{-1}_1 (1-x_0)^2 = x_{-1}$。而 $\CK_{1,0}^2$ 把它送到 $X_{1,0} = x_{-1}$ 与 $X_{0,1} = x_0 (1-x_{-1})^{-2} = x^{-1}_{-2}$。

后一个关系，是 $X_{0,1} = x^{-1}_{-2}$ 与 $X_{1,2} = x_{-2}^{-2} x_{-1} (1-x^{-1}_{-2})^{-2} = x^{-1}_{-3}$ 在 $\CK_{0,1}^2$ 之下的像。现在我们可以对一切 $n$ 继续用 $\CK_{n,n+1}$ 的逆作用，从而再算出一个无穷乘积。关系 $X_{n,n+1} = x^{-1}_{-n-2}$ 与 $X_{n-1,n} = x^{-1}_{-n-1}$ 在 $n\to\infty$ 时的极限为 $X_{1,1} = e^{-2a}$ 与 $X_{1,0} =- e^{-2 b}(1-e^{-2 a})^{-2}$。

于是，沿整条递推关系从 $n= - \infty$ 追到 $n=\infty$，我们便得到无穷乘积的一个表达式
\begin{equation} \label{eq:appA-infinite-product}
\cdots \CK_{4,3}^2 \CK_{3,2}^2 \CK_{2,1}^2 \CK_{1,0}^2 \CK_{0,1}^{-2}
\CK_{1,0}^{-2} \CK_{0,1}^2 \CK_{1,2}^2 \CK_{2,3}^2 \cdots
\end{equation}
递推关系两端极限值之间的映射为 $(X_{1,1},X_{1,0}) \to (X_{1,1}, X_{1,0}(1-e^{-2 a})^4 = X_{1,0}(1-e^{-4 a})^4 (1+e^{-2 a})^{-4})$，这正是所预期的 $\CK_{2,2}^2 \CK_{1,1}^{-4}$！

在正文中，我们把这一公式与具有两种味的 $SU(2)$ \SW 理论的穿墙行为联系起来（其味对称性为 $SO(4)=SU(2)_A \times SU(2)_B$）。我们曾论证，应当有一个类似的关系成立，它携带着关于味荷的信息。该关系应当给出乘积 $\CK_{1,0;1,0} \CK_{1,0;-1,0} \CK_{0,1;0,1} \CK_{0,1;0,-1}$ 的重排，也就是说，给出这样一个理论的穿墙公式：它有一个由电荷为 $(1,0)$ 的粒子组成的 $SU(2)_A$ 二重态（在 $SU(2)_A$ 与 $SU(2)_B$ 的嘉当生成元之下，味荷为 $(\pm 1,0)$），还有一个由电荷为 $(0,1)$ 的粒子组成的 $SU(2)_B$ 二重态（在 $SU(2)_A$ 与 $SU(2)_B$ 的嘉当生成元之下，味荷为 $(0,\pm 1)$）。

基本的变换为
\begin{equation} \label{eq:appA-flavor-A}
\CK_{1,0;1,0} \CK_{1,0;-1,0}: (X_{1,0},X_{0,1}) \to (X_{1,0},X_{0,1}(1-
\mu_A X_{1,0})^{-1}(1- \mu_A^{-1} X_{1,0})^{-1})
\end{equation}
以及
\begin{equation} \label{eq:appA-flavor-B}
\CK_{0,1;0,1} \CK_{0,1;0,-1}: (X_{1,0},X_{0,1}) \to (X_{1,0}(1- \mu_B
X_{0,1})(1- \mu_B^{-1} X_{0,1}),X_{0,1}).
\end{equation}

对这个问题，我们需要把来自 $SU(2)_A$ 或 $SU(2)_B$ 二重态中粒子的那些因子交替相乘。取
\begin{equation} \label{eq:appA-flavor-solution}
\begin{aligned}
& x_n = -\frac{1}{2}\frac{\cosh u \cosh v}{\sinh^2 a} + (-1)^n
\frac{1}{2}\frac{\sinh u \sinh v}{\sinh^2 a}  \\
& - \frac{\sqrt{(\cosh 2
a + \cosh 2 u)(\cosh 2 a + \cosh 2 v)}}{\sinh^2 2 a} \cosh (2 a n +
2 b)
\end{aligned}
\end{equation}
其中 $u,v$ 待下文用 $\mu_A, \mu_B$ 确定。

它满足递推关系
\begin{equation} \label{eq:appA-flavor-recursion}
x_{n+1} x_{n-1} = (1 - e^{u + (-1)^n v} x_n)(1 - e^{-u - (-1)^n v}
x_n).
\end{equation}

我们仍可令递推关系的初值为 $X_{0,1}=x_0$ 与 $X_{1,0}
= x^{-1}_1$。若施加 $\CK_{1,0;1,0} \CK_{1,0;-1,0}$，结果仍是 $X_{1,0}=x^{-1}_1$、$X_{2,1} =x^{-1}_2$，只要我们令 $\mu_A = e^{u-v}$。然后再施加 $\CK_{2,1;0,1}
\CK_{2,1;0,-1}$，结果便是 $X_{2,1}=x_2^{-1}$、$X_{3,2}= x^{-1}_3$，只要我们令 $\mu_B = e^{u+v}$。我们可以对一切 $n$ 交替地施加 $\CK_{2n+1,2n;1,0} \CK_{2n+1,2n;-1,0}$ 与 $\CK_{2n+2,2n+1;0,1} \CK_{2n+2,2n+1;0,-1}$，循着递推关系走到任意大的 $n$。借助关系 $X_{n+1,n} = x^{-1}_{n+1}$ 与 $X_{n,n-1} = x^{-1}_{n}$ 在 $n\to\infty$ 时的极限，便可算出这个无穷乘积。若取 $a,b$ 的实部为正，则 $X_{1,1} = e^{-2 a}$，且
\begin{equation} \label{eq:appA-flavor-forward}
X_{1,0} =- e^{-2 b} (1-e^{-4 a})^2(1+e^{-2a
-2u})^{-1/2}(1+e^{-2a +2u})^{-1/2}(1+e^{-2a -2v})^{-1/2}(1+e^{-2a
+2v})^{-1/2}.
\end{equation}

另一方面，我们也可以沿相反方向追随这个递推关系。关系 $X_{n,n+1} = x^{-1}_{-n-2}$ 与 $X_{n-1,n} = x^{-1}_{-n-1}$ 在 $n\to\infty$ 时的极限为 $X_{1,1} = e^{-2a}$，以及
\begin{equation} \label{eq:appA-flavor-backward}
X_{1,0} =- e^{-2 b} (1-e^{-4 a})^{-2}(1+e^{-2a
-2u})^{1/2}(1+e^{-2a +2u})^{1/2}(1+e^{-2a -2v})^{1/2}(1+e^{-2a
+2v})^{1/2}.
\end{equation}

总的映射为
\begin{equation} \label{eq:appA-flavor-total}
\begin{aligned}
& (X_{1,1},X_{1,0}) \to  \\
& \left(X_{1,1},X_{1,0} \frac{(1-X_{1,1}^2)^{4}}{(1+ \mu_A \mu_B X_{1,1})(1+ \mu_A^{-1} \mu_B X_{1,1})
(1+ \mu_A \mu_B^{-1}X_{1,1})(1+ \mu_A^{-1} \mu_B^{-1}X_{1,1})}\right).
\end{aligned}
\end{equation}

我们认出了所预期的答案：一个电荷为 $(2,2)$、不带味荷的矢量多重态，以及一个电荷为 $(1,1)$ 的超多重态，后者处于味对称性的 $(2_A)\otimes (2_B)$ 表示中（即 $SO(4)$ 的矢量）。

\begin{yizhu}
\begin{enumerate}[leftmargin=2em,itemsep=2pt,topsep=2pt]
\item \textbf{记号}：本附录沿用正文的记号。$\CK_{m,n}$ 是电荷为 $(m,n)$ 的 KS 辛同胚（即 KS 因子），$\CK_{m,n}^{2}$ 是这个因子的平方；带味荷时写作 $\CK_{m,n;s_1,s_2}$，分号后的 $(s_1,s_2)$ 是 $SU(2)_A\times SU(2)_B$ 的味荷。诸 $X_{m,n}$ 是相应电荷的达布坐标，例如本附录用到的 $X_{1,1}=X_{1,0}X_{0,1}$。
\item \textbf{勘误}：本附录中「然后再施加 $\CK_{2,1;0,1}\CK_{2,1;0,-1}$」一处，原文印作 $X_{2,1}={}^{-1}x_2$，上标位置有误，应为 $x_2^{-1}$；译文径改。
\end{enumerate}
\end{yizhu}

% appB —— 附录 B $\CM$ 上的柯西--黎曼方程（原文 wcf-paper.tex 3395–3528 行，PDF 56–57 页；共 7 段、17 条编号公式 (B.1)–(B.17)）
% 记号纪律：$\CM$ = 三维模空间（总空间），$\vm$ = 四维库仑分支（底空间），全块不混。
% 交叉引用：第 5 节 = §5 Mutually non-local corrections（ch05a.tex \label{sec:constructing-moduli}）；式 (3.2) = 一般复结构 $J^{\pz}$（ch03.tex \label{eq:Jzeta}）。合稿时改回 \ref/\eqref。
\section{$\CM$ 上的柯西--黎曼方程}
\label{app:holomorphy}

本附录说明：$(\CM,g)$ 上关于复结构 $J^\pz$ 的柯西--黎曼方程，如何可以改写成 $\vm$ 上某个联络的平坦性方程，而该联络对 $\zeta$ 的依赖非常简单。我们并不假定 $g$ 来自我们在第 5 节所描述的构造；相反，我们只用一些一般事实，它们来自把 $(\CM,g)$ 认定为该规范理论在 $\IR^3\times S^1$ 上的模空间。% TODO 合稿改 \ref{sec:constructing-moduli}

下面，对每个 $\zeta\in\IC^\times$，我们考虑 $\CM$ 上关于复结构 $J^\pz$ 的柯西--黎曼方程
\begin{equation} \label{eq:cauchy-riemann}
\bar\partial f = 0
\end{equation}
我们将把这些方程改写成如下形式
\begin{align}
\partial_{u^i} f &= \CA_{u^i} f, \label{eq:cr-connection-1} \\
\partial_{\bar{u}^{\bar i}} f &= \CA_{{\bar u}^{\bar i}} f, \label{eq:cr-connection-2}
\end{align}
其中 $\CA_{u^i}$ 与 $\CA_{{\bar u}^{\bar i}}$ 是沿环面纤维作用的一阶微分算子（于是在 $\CM$ 的坐标 $(u,\bar u,\theta)$ 下，它们只涉及对 $\theta$ 的导数）；此外，它们还以简单的方式依赖于 $\zeta$：
\begin{align}
\CA_{u^i} &= \frac{1}{\zeta} \CA^{(-1)}_{u^i} + \CA^{(0)}_{u^i}, \\
\CA_{\bar u^{\bar i}} &= \CA_{\bar u^{\bar i}}^{(0)} + \zeta \CA_{\bar u^{\bar i}}^{(1)},
\end{align}
其中诸 $\CA^{(-1)}_{u^i}$ 在每一点都线性无关，诸 $\CA^{(1)}_{\bar{u}^{\bar{i}}}$ 亦然。

我们先把 \eqref{eq:cauchy-riemann} 改写成
\begin{equation} \label{eq:holomorphy}
(1 - i J^{(\zeta,\bar\zeta)}) df = 0.
\end{equation}
若把 $\zeta$ 与 $\bar\zeta$ 看作彼此独立的复变量，那么这个方程其实与 $\bar\zeta$ 无关。为看出这一点，只需在一个固定的 $\zeta$ 处论证即可，例如取 $\zeta=0$。将一般的复结构 (3.2) 特殊化到 $\zeta=0$，% TODO 合稿改 \eqref{eq:Jzeta}
\begin{equation}
J^{(\zeta=0,\bar\zeta)} = J_3 + i \bar\zeta J_+
\end{equation}
其中我们引入了 $J_+ = J_1 + i J_2$。其次注意 $J_+ J_3 = i J_+$，故 $J_+ (J_3 - i) = 0$，于是 $J_+$ 湮灭 $J_3$ 的 $-i$ 特征子空间。这样我们就证明了：$J^{(\zeta=0,\bar\zeta)}$ 与 $J_3$ 共享一个 $n$ 维的、特征值为 $-i$ 的特征子空间。为完成论证，我们还想知道：$J^{(\zeta=0,\bar\zeta)}$ 没有\emph{其他}特征值为 $-i$ 的特征向量。为此我们作一个类似的论证：固定 $\bar\zeta$ 而让 $\zeta$ 变动；这样得到 $J^{(\zeta=0,\bar\zeta)}$ 的 $n$ 个特征值为 $+i$ 的特征向量。于是由维数计数，已没有多余的余地。最终我们看到：$J^{(\zeta,\bar\zeta)}$ 的 $-i$ 特征子空间如所愿地与 $\bar\zeta$ 无关。

于是，在研究柯西--黎曼方程 \eqref{eq:holomorphy} 时，我们可以任取一个方便的 $\bar\zeta$。既然我们想弄清 \eqref{eq:holomorphy} 用 $\vm$ 上的 SW 环面纤维化来表达时是什么样子，自然就取 $\bar\zeta=0$；把它代入 \eqref{eq:holomorphy}，得
\begin{equation} \label{eq:hol-specialized}
(1 - i J_3 - \zeta J_-) df = 0.
\end{equation}
我们假定已给定复辛流形 $(\CM, J_3, \omega_+)$ 与 SW 环面纤维化 $(\CM, J_3^\sf, \omega_+^\sf = - \frac{1}{4 \pi} da^I \wedge dz_I)$ 之间的一个等同。（\cite{Seiberg:1996nz} 中论证过：借助规范理论与引力的弱耦合，这样的等同至少对 $J_3$ 应当存在；类似的论证表明，对 $\omega_+$ 也存在这样的等同。）然后，用与环面纤维相切的向量场 $\frac{\partial}{\partial \bar{z}_I}$ 去缩并 \eqref{eq:hol-specialized}，得到
\begin{equation}
2 \partial_{\bar{z}_I} f - \zeta \left( \partial_{\bar{z}_I} \cdot J_- df \right) = 0.
\end{equation}
为处理第二项，我们利用
\begin{equation}
J_- = g^{-1} \omega_- =  \frac{1}{4 \pi} g^{-1}(d\bar{a}^{\bar I}
\wedge d\bar{z}_{\bar I})
\end{equation}
再乘以 $4 \pi / \zeta$，得到
\begin{equation}
g^{-1}(df, d \bar{a}^I) =  \frac{8 \pi}{\zeta} \partial_{\bar{z}_I} f.
\end{equation}
这已经差不多是我们所要的形式 \eqref{eq:cr-connection-1}，但还不完全是：$g^{-1}(df, d \bar{a}^I)$ 是若干作用于 $f$ 的导数算子的混合。我们想作一次基变换，以便提取出关于 $\frac{\partial f}{\partial a^I}$ 的方程。为此，我们考虑 $g$ 在垂直于 $\CM_u$ 的水平子空间上的限制；把它写成 $g = h_{\bar{I} J} d\bar{a}^{\bar I} da^J$。再乘以 $h_{\bar{I} J}$，便得到
\begin{equation}
h_{\bar{I} J} g^{-1}(df, d \bar{a}^I) =  \frac{8 \pi}{\zeta}
h_{\bar{I} J} \partial_{\bar{z}_I} f.
\end{equation}
现在考虑 $f$ 只依赖于底空间坐标 $(a, \bar{a})$ 的特殊情形。在一点处取一组正交归一基，我们便看出 $g^{-1}(df, d \bar{a}^I) = (h^{-1})^{\bar{I} J} \partial_{a^J} f$。这意味着，对一般的 $f$，左端可以写成
\begin{equation}
\left( \partial_{a_J} - \CA^J \right) f
\end{equation}
其中 $\CA^J$ 是只在纤维方向上作用的微分算子。（更内蕴地说，完整的联络算子 $\partial_{a_J} - \CA^J$ 就是 $f$ 沿向量场 $\partial_{a_J}$ 从 $\vm$ 到 $\CM$ 的水平提升所作的导数。）

总之，我们已经得到
\begin{equation}
\frac{\partial}{\partial a_J} f = \frac{8 \pi}{\zeta} h_{\bar{I} J}
\frac{\partial f}{\partial \bar{z}_I} + \CA^J f,
\end{equation}
这正是所要的形式 \eqref{eq:cr-connection-1}。一个完全相同的论证（从 $\bar\zeta = \infty$ 而非 $\bar\zeta = 0$ 出发）给出共轭方程 \eqref{eq:cr-connection-2}。

我们在此发现的结构，与 \cite{Cecotti:1991me} 中引入的「改进的联络」非常相似。为最清楚地看出这种相似性，引入 $\vm$ 上的一个无穷维丛 $V$，使得 $V$ 在 $u \in \vm$ 处的纤维恰好是环面 $\CM_u$ 上实解析函数的空间，
\begin{equation}
V_u = C^\omega(\CM_u).
\end{equation}
于是，整个 $\CM$ 上的实解析复值函数 $f$，等价于 $V$ 在 $\vm$ 上的一个实解析截面。这样，上面所得到的结果就是：柯西--黎曼方程可以看作 $V$ 中一族具有特定形式的单参数联络的平坦性方程。这些联络的平坦性，是 $\CM$ 上诸复结构可积性的推论。在 \cite{Cecotti:1991me} 中，同样有一个模空间 $\vm$（它参数化 $d=2$ 的 $\CN=(2,2)$ 超对称场论），以及 $\vm$ 上的一个向量丛 $V$（即 Ramond 基态构成的丛）。人们在那里找到 $V$ 中一族由 $\zeta \in \IC^\times$ 参数化的平坦联络，其形式为
\begin{align}
\nabla_i &= \frac{1}{\zeta} C_i + D_i, \\
\nabla_{\bar i} &= \zeta \bar{C}_{\bar i} + \bar{D}_{\bar i},
\end{align}
其中 $D_i$ 是由耦合的绝热变化所提供的标准联络，诸 $C_i$ 则是「手征环」算符。这些联络的平坦性，是著名的 $tt^*$ 方程的推论。

在本文通篇，尤其是在第 5 节，许多构造——以及它们的物理解释——都与 $tt^*$ 故事中出现的那些相平行。% TODO 合稿改 \ref{sec:constructing-moduli}

% appC1 —— 附录 C 前段：积分方程的渐近（有限维情形）（原文 wcf-paper.tex 3529–3685 行；共 10 段、13 条编号公式 (C.1)–(C.13)）
\section{积分方程的渐近}
\label{app:integral-asymptotics}

本附录中，我们先说明如何修改 \cite{Cecotti:1993rm} 的渐近分析，使之适用于存在多条 BPS 射线的情形；然后把这个分析改造成适合我们问题的形式。

\subsection*{有限维情形}

文献 \cite{Cecotti:1993rm} 研究的是复 $x$ 平面上的一个黎曼--希尔伯特问题：关于 $m \times m$ 矩阵 $\Psi(x)$，其不连续性沿实轴分布：
\begin{align}
\Psi(y e^{- i \epsilon}) &= \Psi(y e^{ i \epsilon}) S \qquad \text{当 } y \in
\IR^+, \notag \\
\Psi(y e^{- i \epsilon}) &= \Psi(y e^{ i \epsilon}) S^t \qquad \text{当 } y \in
\IR^-. \label{eq:cv-discont}
\end{align}
（此处我们已回到算符复合的标准约定。）这里充当中心荷 $Z_\gamma$ 之对应物的是复数 $\Delta_{ij}$（$i,j = 1, \dots, m$），它们满足 $\Delta_{ij} = w_i-w_j$，其中诸 $w_i$ 为某些复数。矩阵 $S$ 是三角矩阵，对角元为 $1$，并且当 $\re \Delta_{ij} <0$ 时 $S_{ij}=0$。

$\Psi$ 的渐近行为由诸常数 $w_i$ 决定。若定义
\begin{equation} \label{eq:appC-Phi-def}
\Phi_{ij}(x) = \Psi_{ij}(x) e^{- \beta x w_j - \frac{\beta}{x} \bar w_j},
\end{equation}
则 $\Phi(x)$ 在 $x\to \infty$ 时趋于单位矩阵，在 $x\to 0$ 时趋于某个「度量」$g_{ij}$。

矩阵 $S$ 是该问题的「斯托克斯乘子」。把它重新表示成若干个更基本的「斯托克斯因子」之乘积，会比较方便。事实上，只要假定没有三个 $w_i$ 在复平面中共线，就存在唯一的分解
\begin{equation}\label{eq:appC-S-decomp}
S = \prod^\cwarrow_{(ij): \re \Delta_{ij} > 0} s_{(ij)}, \quad S^t = \prod^\cwarrow_{(ij): \re \Delta_{ij} < 0} s_{(ji)}
\end{equation}
其中两个乘积都按 $\arg \Delta_{ij}$ 递增的次序取，而每个 $s_{(ij)}$ 的对角元为 $1$，并且在位置 $(ji)$ 处有唯一的非零非对角元，其值为 $- \mu_{ij} = -\mu_{ji}$。

利用这一分解，我们可以引入黎曼--希尔伯特问题的「多射线」版本。具体地说，在 $x$ 平面上引入一组过原点的射线
\begin{equation}\label{eq:appC-rays}
\ell_{(ij)} = \{ x: x \Delta_{ij} \in \IR_+ \},
\end{equation}
并要求
\begin{equation}\label{eq:appC-jump}
\Psi(y e^{- i \epsilon}) = \Psi(y e^{ i \epsilon}) s_{(ij)} \qquad \text{当 } y \in
\ell_{(ij)}.
\end{equation}
重要的是，可以证明：无论采用单射线问题还是多射线问题，$g_{ij}$——它在文献 \cite{Cecotti:1993rm} 中正是主要的研究对象——都相同。

文献 \cite{Cecotti:1993rm} 中关于 $\Phi(x)$ 的积分方程 (4.17) 为（其中对 $\epsilon$ 的约定作了少许修改，以便后文使用）：
\begin{align}
\Phi_{ij}(x) = \delta_{ij} &+ \frac{1}{2 \pi i} \int_0^\infty
\frac{dy}{y-x} \sum_k \Phi_{ik}(e^{i \epsilon} y)(1-S)_{kj}e^{-
\beta y \Delta_{kj} - \beta/y \bar \Delta_{kj}}  \nonumber  \\
&+ \frac{1}{2 \pi i}
\int_{-\infty}^0 \frac{dy}{y-x} \sum_k \Phi_{ik}(e^{-i \epsilon}
y)(1-S^t)_{kj}e^{- \beta y \Delta_{kj} - \beta/y \bar \Delta_{kj}}. \label{eq:cv-integral}
\end{align}
该方程的一个解给出单射线黎曼--希尔伯特问题的一个解。现在我们表述一个新的积分方程，它与多射线问题等价：
\begin{equation} \label{eq:cv-multi-line}
\Phi(x)_{ij} = \delta_{ij} + \sum_{k} \frac{1}{2 \pi i}
\int_{\ell_{(kj)}} \frac{dy}{y-x}\Phi_{ik}(y)\mu_{kj}e^{- \beta y
\Delta_{kj} - \frac{\beta}{y} \bar \Delta_{kj}}
\end{equation}
这里的积分应理解为主值积分。注意：当没有三个 $w_i$ 共线时，矩阵元 $\Phi_{ik}$ 在穿过所有射线 $\ell_{(kj)}$ 时都是连续的。

\eqref{eq:cv-multi-line} 的递推解取一个简单的形式：
\begin{equation}\label{eq:appC-Phi-series}
\Phi(x)_{ij} = \delta_{ij} + \Phi(x)_{ij}^{(1)} + \Phi(x)_{ij}^{(2)} + \cdots,
\end{equation}
其中
\begin{align}
\Phi(x)_{ij}^{(1)} &= \frac{1}{2 \pi i} \mu_{ij}
\int_{\ell_{(ij)}} \frac{dy}{y-x} e^{ -\beta y
\Delta_{ij} - \frac{\beta}{y} \bar \Delta_{ij} }, \label{eq:appC-Phi12}\\
\Phi(x)_{ij}^{(2)} &= \frac{1}{(2 \pi i)^2} \sum_{i_2} \mu_{i i_2} \mu_{i_2 j}
\int_{\ell_{(i i_2)}} \int_{\ell_{(i_2 j)}} \frac{dy_1}{y_1-y_2}
\frac{dy_2}{y_2-x}e^{ -\beta
y_1 \Delta_{i i_2} - \frac{\beta}{y_1} \bar \Delta_{i i_2} }e^{ -\beta
y_2 \Delta_{i_2 j} - \frac{\beta}{y_2} \bar \Delta_{i_2 j} }, \notag
\end{align}
一般地，$\Phi(x)^{(n)}$ 涉及沿 $n$ 条射线 $\ell_{(i_k i_{k+1})}$ 的所有链的积分，其中 $i_1 = i$、$i_{n+1} = j$。
沿 $\ell_{(i_k i_{k+1})}$ 的每个积分都含有因子 $e^{ - \beta y \Delta_{i_k i_{k+1}} - \frac{\beta}{y} \bar \Delta_{i_k i_{k+1}} }$。在这条射线上，该指数是实的且为负，并在 $y = \exp (-i \arg \Delta_{i_k i_{k+1}})$ 处有唯一的峰。当 $\beta$ 取得很大时，该积分便由鞍点法很好地近似：把被积函数的其余部分换成它在峰点处的值。于是，固定 $x$ 时的大 $\beta$ 渐近式就简单地是：\footnote{这些渐近式成立，除非 $x$ 恰好落在对 $y_n$ 积分的鞍点上，即 $x = \exp [-i \arg \Delta_{i_n j}]$。在这一点上，我们得到类似的大 $\beta$ 渐近式，只是诸 $\sqrt{\beta}$ 压低因子中少了一个。}
\begin{equation}\label{eq:appC-asym1}
\Phi(x)^{(1)}_{ij}  \sim \frac{1}{2 \sqrt{\pi \beta |\Delta_{ij}|} i}
\mu_{ij}\frac{1}{\exp[-i \arg \Delta_{ij}]-x} e^{- 2 \beta
|\Delta_{ij}|},
\end{equation}
以及
\begin{equation}\label{eq:appC-asym2}
\begin{aligned}
\Phi(x)^{(2)}_{ij} &\sim \sum_{i_2} \mu_{i i_2} \mu_{i_2 j} \frac{1}{2 \sqrt{\pi \beta
|\Delta_{i i_2}|} i}\frac{1}{2 \sqrt{\pi \beta |\Delta_{i_2 j}|} i} \times \\
&\qquad \frac{1}{\exp[-i \arg \Delta_{i_2 j}]-x}\frac{1}{\exp[-i \arg
\Delta_{i i_2}]-\exp[-i \arg \Delta_{i_2 j}]}e^{- 2 \beta |\Delta_{i
i_2}|}e^{- 2 \beta |\Delta_{i_2 j}|},
\end{aligned}
\end{equation}
更高阶的 $\Phi^{(n)}$ 有类似的估计。由三角不等式可见，在大 $\beta$ 极限下 $\Phi^{(2)}$ 相对于 $\Phi^{(1)}$ 是指数压低的；同样，$\Phi^{(n+1)}$ 相对于 $\Phi^{(n)}$ 也是压低的。$\Phi^{(n)}$ 所具有的指数压低，恰好就是 $n$-瞬子修正所应期待的量级。

我们注意到，这一渐近分析比 \eqref{eq:cv-integral} 的相应分析要简单得多；在那种情形，必须变形积分围道使其穿过相应的鞍点，而途中会遇到割线与极点，其贡献必须仔细追踪。

为结束本节，我们简要讨论这一展开的解析性质。如上所见，$\Phi$ 的第 $n$ 阶修正可以用某些迭代积分表示：
\begin{equation}\label{eq:appC-Fn}
\begin{aligned}
\CF^{(n)}[x,\beta;\Delta_{i_1 i_2},\dots, \Delta_{i_{n} i_{n+1}}] ={}& \\
&\prod_{k=1}^n \left[ \int_{\ell_{(i_k i_{k+1})}} \frac{dy_k}{2 \pi i} e^{ -\beta y_k \Delta_{i_k i_{k+1}} -
\frac{\beta}{y_k} \bar \Delta_{i_k i_{k+1}} } \right] \\
&\times \prod_{k=1}^n
\frac{1}{y_k - y_{k+1}} \Bigg\vert_{y_{n+1} = x}.
\end{aligned}
\end{equation}
这个 $\CF^{(n)}$ 在射线 $\arg x = -\arg \Delta_{i_n i_{n+1}}$ 上有明显的不连续性，该不连续性等于
\begin{equation}\label{eq:appC-disc}
\CF^{(n-1)}[x,\beta;\Delta_{i_1 i_2}, \dots, \Delta_{i_{n-1}
i_{n}}]e^{ - \beta x \Delta_{i_{n} i_{n+1}} - \frac{\beta}{x} \bar
\Delta_{i_{n} i_{n+1}} }.
\end{equation}
注意，$\CF^{(n)}$ 对自变量 $d_k = \Delta_{i_k i_{k+1}}$ 的一般取值是有意义的，不必受 $\Delta_{ij} = w_i-w_j$ 这一限制。（事实上，甚至可以把 $\Delta_{ij}$ 与 $\bar \Delta_{ij}$ 取作彼此独立的复参数，而把诸射线放在 $x^2 \Delta_{ij}/\bar \Delta_{ij} \in \IR^+$ 处。）只要两个相邻自变量 $d_k, d_{k+1}$ 的相位对齐，它就有割线。此时的不连续性就是取同样自变量的 $\CF^{(n-1)}$，只须作替换 $d_k, d_{k+1} \to d_k + d_{k+1}$。

在文献 \cite{Cecotti:1993rm} 的渐近分析中，出现过本质上相同的函数 $\CF^{(n)}$；我们的函数与它们只差在分支割线的位置。

% appC2 —— 附录 C 后段：Infinite-dimensional case（原文 wcf-paper.tex 3686–3822 行）
% 原文小节标题为 \subsection*{Infinite-dimensional case}（不编号；TOC 中亦无 C.1/C.2 条目），
% 故此处保持带星号，与原文排版及上一块的 \subsection*{有限维情形} 一致。
% 原文宏 \eps 不在本工程 preamble 中，此处本地补齐（与原 wcf-paper.tex 的定义相同）。
\providecommand{\eps}{\epsilon}

\subsection*{无穷维情形}

现在我们对乘性积分方程 式 (5.13) 给出类似的分析。% TODO 合稿改 \eqref{eq:X-integral-mult-explicit}
首先，对任意向量 $\gamma \in \clat$，我们用如下幂级数展开来定义一个向量 $f^\gamma \in \clat_\IQ$：
\begin{equation}
- \sum_{\gamma' \in \clat} \left( \degen(\gamma';u)
\log (1- \sigma(\gamma')
\CX_{\gamma'} ) \right) \gamma' = \sum_{\gamma' \in \clat} f^{\gamma'} \CX_{\gamma'}
\end{equation}
或者更显式地，
\begin{equation}
f^\gamma = \sum_{n \ge 1 \text{ s.t. } \gamma = n \gamma'} \frac{\sigma(\gamma')^n}{n} \degen(\gamma';u) \gamma'.
\end{equation}
这样定义的用意在于：此后 式 (5.13) 便取如下形式% TODO 合稿改 \eqref{eq:X-integral-mult-explicit}
\begin{equation} \label{eq:X-integral-mult-std}
\CX_\gamma(\zeta) = \CX^{\sfm}_\gamma(\zeta) \exp \left\langle \gamma,
\frac{1}{4 \pi i} \sum_{\gamma'} f^{\gamma'} \int_{\ell_{\gamma'}} \frac{d\zeta'}{\zeta'} \frac{\zeta' +
\zeta}{\zeta' - \zeta} \CX_{\gamma'}(\zeta') \right\rangle.
\end{equation}

我们希望把 式 \eqref{eq:X-integral-mult-std} 的解 $\CX$ 构造为逐次逼近 $\CX^{(n)}$ 的极限，或者构造为 式 (5.1) 中所定义的 $\Upsilon$ 的相应逼近 $\Upsilon^{(n)}$ 的极限。% TODO 合稿改 \eqref{eq:def-upsilon}
$\Upsilon$ 是从 $\CM_u$ 到其复化的一个映射；我们把它的分量写成函数，即 $\Upsilon_\gamma(\theta) := \gamma \cdot \Upsilon(\theta)$。
我们取 $\Upsilon^{(0)}_\gamma = \theta_\gamma$ 作为起点。
回忆 $\Upsilon$ 的定义方式使下式成立
\begin{equation}
\CX^{\sfm}(\Upsilon^{(n)}) = \CX^{(n)}(\theta),
\end{equation}
于是迭代步可写成
\begin{equation} \label{eq:recursion}
e^{i \Upsilon^{(n+1)}_\gamma} = e^{i \theta_\gamma} \exp \left\langle \gamma,
\frac{1}{4 \pi i} \sum_{\gamma'} f^{\gamma'} \int_{\ell_{\gamma'}}
\frac{d\zeta'}{\zeta'} \frac{\zeta' + \zeta}{\zeta' - \zeta}
\CX^{\sfm}_{\gamma'}(\Upsilon^{(n)},\zeta')\right\rangle
\end{equation}
该迭代的不动点 $\Upsilon^{(n+1)} = \Upsilon^{(n)}$ 将给出 式 (5.11) 的一个解。% TODO 合稿改 \eqref{eq:X-integral-mult}
因此，要看出解确实存在，我们应当验证这一迭代是压缩的，即
\begin{equation}
\max_{\zeta,\theta} \norm{\Upsilon^{(n+1)} - \Upsilon^{(n)}} < C \max_{\zeta,\theta} \norm{\Upsilon^{(n)} - \Upsilon^{(n-1)}}
\end{equation}
其中 $C < 1$ 为某常数。
更确切地说，我们将验证：当这一迭代作用在满足 $\max_{\zeta,\theta} \norm{\Upsilon^{(n)} - \theta} < \infty$、并且还满足一个附加条件的诸 $\Upsilon^{(n)}$ 上时，它是压缩的；该附加条件表达的正是「它们离实环面不太远」这一事实：我们要求
$\max_{\zeta,\theta} |e^{i \Upsilon^{(n)}_\gamma}|<e^{\eps \norm{\gamma}}$，其中常数 $\eps>0$ 稍后确定。

首先我们要看出，这一迭代保持我们的附加条件。
取 式 \eqref{eq:recursion} 的绝对值并作鞍点分析，当 $R$ 足够大时便得到估计\footnote{与上一小节一样，当 $\zeta$ 恰好落在鞍点上时，这一分析还须补上一段单独的讨论；但那只会使压低程度相差一个 $\sqrt{R}$，仍足以让我们建立 式 \eqref{eq:estimate}。}
\begin{equation} \label{eq:estimate}
|e^{i \Upsilon^{(n+1)}_\gamma}| <  \exp \left[ \sum_{\gamma'} \abs{\inprod{f^{\gamma'},\gamma}}
e^{- 2 \pi R |Z_{\gamma'}| + \eps \norm{\gamma'}} \right].
\end{equation}
现在假设 $\sum_{\gamma'} f^{\gamma'} \norm{\gamma'} e^{- 2 \pi R |Z_{\gamma'}|}$ 在 $R$ 足够大时收敛。（这相当于要求诸 $\degen(\gamma';u)$ 随 $\gamma'$ 增长得不太快；对场论而言这看来非常合理，但在引力情形几乎肯定会被破坏。）
我们还利用第 2.2 节所回顾的支撑性质（Support Property），用 $K \abs{Z_{\gamma'}}$ 给出 $\norm{\gamma'}$ 的上界。% TODO 合稿改 \ref{sec:ReviewKS}
于是当 $R$ 足够大时，我们可以选取 $\eps$，使右端小于 $e^{\eps \norm{\gamma}}$。
对 $\gamma$ 取遍 $\clat$ 的一组基逐一这样做，就得到所需的 $\eps$。（事实上，当 $R$ 很大时，可以取 $\epsilon \to 0$ 以指数速度变快。）

现在我们想估计
\begin{align}
\norm{\Upsilon^{(n+1)}-\Upsilon^{(n)}} &=
\frac{1}{4 \pi} \bigg\lVert \sum_{\gamma' \in \clat} \int_{\ell_{\gamma'}}
\frac{d\zeta'}{\zeta'} \frac{\zeta' + \zeta}{\zeta' - \zeta}
f^{\gamma'}\left(\CX^{\sfm}_{\gamma'}(\Upsilon^{(n)},\zeta')-\CX^{\sfm}_{\gamma'}(\Upsilon^{(n-1)},\zeta')\right) \bigg\rVert \\
&\le \frac{1}{4 \pi}\bigg\lVert \sum_{\gamma' \in \clat} \int_{\ell_{\gamma'}}
\frac{d\zeta'}{\zeta'} \frac{\zeta' + \zeta}{\zeta' - \zeta}
f^{\gamma'} \abs{\CX^{\sfm}_{\gamma'}(\theta, \zeta')} \left(e^{i \Upsilon^{(n)}_{
\gamma'}(\zeta')} - e^{i \Upsilon^{(n-1)}_{\gamma'}(\zeta')}\right) \bigg\rVert \\
&\le \frac{1}{4 \pi}\bigg\lVert \sum_{\gamma' \in \clat} \int_{\ell_{\gamma'}}
\frac{d\zeta'}{\zeta'} \frac{\zeta' + \zeta}{\zeta' - \zeta}
f^{\gamma'} \abs{\CX^{\sfm}_{\gamma'}(\theta, \zeta')}  \bigg\rVert e^{\eps \norm{\gamma'}} \norm{\gamma'} \max_{\zeta}  \norm{\Upsilon^{(n)} -
\Upsilon^{(n-1)}}.
\end{align}
随后，大 $R$ 下的鞍点分析给出
\begin{equation}
\norm{\Upsilon^{(n+1)}-\Upsilon^{(n)}} \le
\frac{1}{4\pi}\max_{\zeta} \norm{\Upsilon^{(n)}(\zeta) -
\Upsilon^{(n-1)}(\zeta)} \bigg\lVert \sum_{\gamma' \in \clat}
f^{\gamma'} \norm{\gamma'} e^{- 2 \pi R |Z_{\gamma'}| + \eps
\norm{\gamma'}} \bigg\rVert.
\end{equation}
当 $R$ 足够大时，在我们的收敛假设下，这就建立了压缩性质；事实上该迭代收敛得很快，其速度由最大的那个 $e^{- 2 \pi R |Z_{\gamma'}|}$ 决定。

我们可以用与有限维情形中 $\CF$ 同类的函数，写出 $\CX_\gamma(\zeta)$ 的显式表达式。
递归关系中指数因子的出现使事情稍微复杂一些：现在不再对链求和，而是对带装饰的有根树求和。
设 $\CT$ 是一棵有根树，其边用数对 $(i,j)$ 标记（其中 $i$ 是更靠近根节点的那个节点），每个节点则带一个装饰，即取定 $\clat$ 中的一个 $\gamma_i$。
此外，把根节点处的装饰记为 $\gamma_\CT$。
于是定义该树的权重为 $\clat_\IQ$ 中的一个元素，它由诸 $\degen(\gamma;u)$ 确定：
\begin{equation}
\CW_\CT = \frac{f^{\gamma_\CT}}{\abs{\mathrm{Aut}(\CT)}} \prod_{(i,j) \in {\mathrm {Edges}}(\CT)} \inprod{\gamma_i, f^{\gamma_j}}.
\end{equation}
于是 $\CX_\gamma(\zeta)$ 的迭代解取如下形式
\begin{equation} \label{eq:iterative-X}
\CX_\gamma(\zeta) = \CX^{\sfm}_\gamma (\zeta)\exp \left\langle \gamma, \sum_\CT
\CW_\CT \CG_\CT(\zeta) \right\rangle,
\end{equation}
其中 $\CG_\CT(\zeta)$ 为某些函数。关于 $\CX$ 的积分方程 式 \eqref{eq:X-integral-mult-std} 则变成一条公式，把每个 $\CG_\CT(\zeta)$ 用较小的树所对应的那些函数表示出来。
具体说来，从 $\CT$ 中删去根节点，会留下一组有根树 $\CT_a$；而 式 \eqref{eq:X-integral-mult-std} 将被满足，只要
\begin{equation}
\CG_\CT(\zeta) = \frac{1}{4 \pi i} \int_{\ell_{\gamma_\CT}} \frac{d \zeta'}{\zeta'} \frac{\zeta' + \zeta}{\zeta' - \zeta} \CX^{\sfm}_{\gamma_\CT}(\zeta') \prod_a \CG_{\CT_a}(\zeta').
\end{equation}
由此可见，与有限维情形中的 $\CF^{(n)}$ 一样，$\CG_\CT$ 沿 $\ell_{\gamma_\CT}$ 的不连续性由较低的那些 $\CG_{\CT_a}$ 之乘积决定。

由 式 \eqref{eq:iterative-X} 我们还能直接得到 $\Upsilon$：
\begin{equation}
e^{i \Upsilon_\gamma} = \exp \left\langle \gamma, \sum_\CT
\CW_\CT \CG_\CT(\zeta) \right\rangle.
\end{equation}
特别地，这使我们能够计算 $\Upsilon(\zeta = 0)$。
辛形式 $\hsymp(\zeta)$ 的展开也可用同样的方式分析，并组织成对树的求和。

% appD —— 附录 D 微分方程的渐近（原文 wcf-paper.tex 3823–3899 行，PDF 61–62 页；共 4 段、14 条编号公式）
% 原文宏 \plog 不在本工程 preamble 中，此处本地补齐（定义与原文 wcf-paper.tex 159 行相同；ch04 已 providecommand，重复无害）。
% 原文把 \sf 重定义为 \mathrm{sf}（wcf-paper.tex 157 行），本工程对应宏为 \sfm，故原文的 \sf 一律写作 \sfm。
\providecommand{\plog}[1]{#1\partial_{#1}}
\section{微分方程的渐近}
\label{app:diff-asymptotics}

在本附录中，我们希望弄清楚如何计算 $\CX$ 所满足的微分方程的系数。考虑 $\CX$ 在 $\zeta=0$ 附近的渐近展开：
\begin{equation} \label{eq:as-X}
\CX(\theta,\zeta)_\gamma \sim
\CX^{\sfm}_\gamma(\Upsilon_0(\theta),\zeta) \exp \sum_{n>0} \zeta^n
\gamma \cdot g_n(\theta).
\end{equation}
我们考虑第 5.5 节所定义的微分算子，% TODO 合稿改 \ref{sec:differential-problem}（第 5.5 节即原文 §5.5 ``Differential equations''）
\begin{equation}
\frac{1}{\CX_\gamma}\frac{\partial \CX_\gamma}{\partial t^n}  =
\frac{1}{\CX_\gamma} \CA_n \CX_\gamma.
\end{equation}
把展开式 \eqref{eq:as-X} 代入，只保留最初几项：
\begin{equation}
\frac{\partial \log \CX^{\sfm}_\gamma(\Upsilon_0(\theta))}{\partial
t^n} = \left(\frac{1}{\zeta}\CA^{(-1)}_n + \CA^{(0)}_n\right)
\Upsilon_0(\theta)_\gamma+ \CA_n^{(-1)} \gamma \cdot g_1(\theta).
\end{equation}
$\zeta$ 的领头部分是一条我们已经理解的陈述：
\begin{equation}
\pi R \frac{\partial Z_\gamma}{\partial t^n} = \CA^{(-1)}_n
\Upsilon_0(\theta)_\gamma.
\end{equation}
这意味着，$\CA^{(-1)}$ 是 $\CA^{(-1)}_{\sfm}$ 经 $\Upsilon_0$ 拉回所得。

展开中的下一项是
\begin{equation}
i \frac{\partial \Upsilon_0(\theta)_\gamma}{\partial t^n} =
\CA^{(0)}_n \Upsilon_0(\theta)_\gamma+ \CA_n^{(-1)} \gamma \cdot
g_1(\theta),
\end{equation}
在已知 $\Upsilon_0(\theta)$ 的前提下，它确定了 $\CA^{(0)}_n$。

还有另一种观点，它相当有用：考虑各个微分方程之间的相容性条件。例如，考虑方程 $[R\partial_R - \CA_R, \zeta \partial_\zeta - \CA_\zeta] = 0$，亦即
\begin{equation} \left[ \plog{R} + \frac{1}{\zeta} {\CA_{\zeta}^{(-1)}} - \CA_R^{(0)} -
\zeta \CA^{(1)}_{\zeta} \, , \, \plog{\zeta} -
\frac{1}{\zeta} {\CA_{\zeta}^{(-1)}} - \CA_{\zeta}^{(0)} - \zeta
\CA^{(1)}_{\zeta} \right] = 0,
\end{equation}
把它按 $\zeta$ 的幂次展开。这给出三个方程：
\begin{align}
\plog{R} \CA_{\zeta}^{(-1)} - [\CA_R^{(0)},\CA_{\zeta}^{(-1)}] &=
[\CA_{\zeta}^{(0)},\CA_{\zeta}^{(-1)}] + \CA_{\zeta}^{(-1)},\\
\plog{R} \CA_{\zeta}^{(0)} - [\CA_R^{(0)},\CA_{\zeta}^{(0)}] &=
2 [\CA_{\zeta}^{(1)},\CA_{\zeta}^{(-1)}],\\
\plog{R} \CA_{\zeta}^{(1)} - [\CA_R^{(0)},\CA_{\zeta}^{(1)}] &=
[\CA_{\zeta}^{(1)},\CA_{\zeta}^{(0)}] + \CA_{\zeta}^{(1)}.
\end{align}
这些方程与 Nahm 方程极为相像，与后者的差别仅在于右端多出的两个线性项。这两个额外项在大半径处占主导。要推导大 $R$ 渐近，另一种策略仍是围绕半平坦解对这三个方程作迭代求解。

另一组有趣的「等单值」方程，可以通过类似地展开 $[\partial_u - \CA_u, \zeta\partial_\zeta - \CA_\zeta] = 0$ 得到：
\begin{align}
0 &=[\CA^{(-1)}_u, \CA^{(-1)}_\zeta], \\
\frac{\partial}{\partial u} \CA_{\zeta}^{(-1)} -
[\CA_u^{(0)},\CA_{\zeta}^{(-1)}] &=
[\CA_u^{(-1)},\CA_{\zeta}^{(0)}] - \CA_{u}^{(-1)},\\
\frac{\partial}{\partial u} \CA_{\zeta}^{(0)} -
[\CA_u^{(0)},\CA_{\zeta}^{(0)}]  &=
[\CA_u^{(-1)},\CA_{\zeta}^{(1)}], \\
\frac{\partial}{\partial u} \CA_{\zeta}^{(1)} -
[\CA_u^{(0)},\CA_{\zeta}^{(1)}] &=
[\CA_u^{(1)},\CA_{\zeta}^{(0)}] + \CA_{u}^{(-1)}, \\
0 &= [\CA^{(1)}_u, \CA^{(1)}_\zeta].
\end{align}

% appE —— 附录 E 与热力学贝特拟设的关系（原文 wcf-paper.tex 3900–4041 行）
\section{与热力学贝特拟设的关系}
\label{app:TBA}

{\itshape 2009 年 11 月 20 日补注}：

不久前，A. Zamolodchikov 向我们指出：本文的核心结果之一，即式 (5.13)，实际上是热力学贝特拟设（Thermodynamic Bethe Ansatz，下称 TBA）的一个版本 \cite{Zamolodchikov:1989cf}。本附录解释这一评注。四维超杨--米尔斯（super Yang--Mills）理论与 TBA 之间的另一个联系，最近也已由 Nekrasov 与 Shatashvili 讨论过 \cite{Nekrasov:2009ui}。% TODO 合稿改 \eqref{eq:X-integral-mult-explicit}

考虑一个可积的粒子系统：粒子记作 $a$，质量为 $m_a$，处于逆温度 $\beta$ 下，散射矩阵 $S_{ab}(\theta-\theta')$ 可积，其中 $\theta$ 为快度（rapidity）。该系统的 TBA 方程为
%
\begin{equation}\label{eq:gen-TBA}
\epsilon_a(\theta) = m_a \beta \cosh \theta - \sum_b
\int_{-\infty}^{+\infty} \frac{d \theta'}{2\pi}
\phi_{ab}(\theta-\theta') \log (1 + e^{\beta \mu_b -
\epsilon_b(\theta')})
\end{equation}
%
其中 $\phi_{ab}(\theta) = - i\frac{\p}{\p \theta} \log
S_{ab}(\theta)$。这里散射矩阵是对角的，也就是说，孤子产生算子满足 $\Phi_a(\theta) \Phi_b(\theta') =
S_{ab}(\theta-\theta') \Phi_b(\theta') \Phi_a(\theta)$。

我们可以把式 (5.13) 的对数写成 \eqref{eq:gen-TBA} 的形式，做法如下。显然，粒子标号 $a, b, \dots$ 对应于 $\gamma, \gamma', \dots$。现在令% TODO 合稿改 \eqref{eq:X-integral-mult-explicit}（式 (5.13)）
$Z_\gamma =
e^{i \alpha_\gamma} \vert Z_\gamma \vert$，其中 $\alpha_\gamma$ 为实数，且只在模 $2\pi$ 的意义下有定义。对任意一个 $\gamma$，我们都可以作变量替换 $\zeta = - e^{i \alpha_\gamma+ \theta}$，使 BPS 射线 $\ell_\gamma$ 由 $-\infty < \theta <
\infty$ 扫出。在这一变量替换下，半平坦坐标式 (3.9) 变为% TODO 合稿改 \eqref{eq:X-integral-mult-explicit}（式 (5.13)）、\eqref{eq:X-sf}（式 (3.9)）
%
\begin{equation}
\log \CX_{\gamma}^{\sf} = - 2 \pi R \lvert Z_\gamma \rvert \cosh \theta
+ i \varphi_\gamma.
\end{equation}
%
（注意：为了避免与快度 $\theta$ 混淆，我们把环面上角坐标的记号，从本文其余部分所用的 $\theta_\gamma$ 改成了 $\varphi_\gamma$。）现在置
%
\begin{equation}\label{eq:defmugamma}
\beta \mu_\gamma := i \varphi_\gamma + \log(-\sigma(\gamma))  \quad \mod \ 2
\pi i.
\end{equation}
%
注意，$\beta \mu_\gamma$ 就是 $i \varphi_\gamma$，或者与之相差 $\pm i \pi$。特别地，它是纯虚数。我们按下式定义「准粒子（quasiparticle）能量」$\epsilon_{\gamma}(\theta)$：
%
\begin{equation} \label{eq:qparticle}
\CX_{\gamma}(\zeta = - e^{i \alpha_\gamma+ \theta}) :=
-\sigma(\gamma) e^{\beta \mu_\gamma - \epsilon_{\gamma}(\theta)}=
e^{i \varphi_\gamma - \epsilon_\gamma(\theta)}.
\end{equation}
%
更确切地说，上式所定义的是 BPS 射线 $\ell_\gamma$ 上的 $\eps_\gamma$，在那里 $\theta$ 取实值。对其他 $\theta$，我们用解析延拓来定义 $\eps_\gamma$——这与 $\CX_\gamma$ 恰成对照：后者沿着某些 $\Im\,\theta$ 为常数的直线（即 BPS 射线）带有不连续性。

我们选取 \eqref{eq:defmugamma}，正是为了使式 (5.13) 的对数读作% TODO 合稿改 \eqref{eq:X-integral-mult-explicit}（式 (5.13)）
%
\begin{equation}\label{eq:GMN-2}
\epsilon_{\gamma}(\theta) = 2\pi R \vert Z_\gamma \vert \cosh\theta
+ \sum_{\gamma'} \Omega(\gamma') \int_{-\infty}^{+\infty}
\frac{d\theta'}{2\pi} K_{\gamma,\gamma'}(\theta-\theta') \log( 1 +
e^{\beta \mu_{\gamma'} - \epsilon_{\gamma'}(\theta')} ),
\end{equation}
%
其中
%
\begin{equation}
K_{\gamma,\gamma'}(\theta-\theta') = \frac{i}{2} \langle \gamma,
\gamma'\rangle \frac{ e^{ \theta - \theta' + i \alpha_{\gamma} - i
\alpha_{\gamma'}} +1 }{e^{ \theta - \theta' + i \alpha_{\gamma} - i
\alpha_{\gamma'}} -1 }.
\end{equation}
%
这个核也可以写成
%
\begin{equation}
\begin{aligned}
K_{\gamma,\gamma'}(\theta-\theta') & = i \langle \gamma, \gamma'
\rangle \frac{\p}{\p \theta} \log \bigl[\sinh\left(\half (\theta -
\theta' + i \alpha_{\gamma} - i \alpha_{\gamma'})\right)\bigr]\\
& =\frac{i}{2}\langle \gamma, \gamma' \rangle  \coth\left( \frac{
\theta-\theta' + i \alpha_{\gamma} - i \alpha_{\gamma'}}{2} \right).
\end{aligned}
\end{equation}
%
对数的宗量并不是一个纯相位，因此一般说来，$K_{\gamma,\gamma'}(\theta-\theta')$ 并不对应于一个幺正的散射矩阵。

下面简略讨论「准粒子能量」的实性（reality）性质。扭量坐标满足如下实性关系
%
\begin{equation}
\overline{\CX_{\gamma}(\zeta)} = \CX_{-\gamma}(-1/\bar \zeta).
\end{equation}
%
由于 $Z_{-\gamma} = - Z_\gamma$，我们有
%
\begin{equation}
e^{i \alpha_{-\gamma}} = - e^{i \alpha_\gamma}.
\end{equation}
%
因此，若 $\zeta = - e^{i \alpha_{\gamma} + \theta}$ 且 $\theta$ 为实数，则 $-1/\bar \zeta = - e^{i \alpha_{-\gamma} - \theta}$。现在，利用 $\sigma(-\gamma) = \sigma(\gamma)$ 与 $\varphi_{-\gamma} = -
\varphi_{\gamma}$，便得到「准粒子能量」的实性条件
%
\begin{equation}\label{eq:reality}
\overline{\epsilon_{\gamma}(\theta)} = \epsilon_{-\gamma}(-\theta).
\end{equation}
%
积分方程式 (5.13) 与实性条件 \eqref{eq:reality} 相容，因为当 $\theta, \theta'$ 均为实数时% TODO 合稿改 \eqref{eq:X-integral-mult-explicit}（式 (5.13)）
%
\begin{equation}
\overline{K_{\gamma, \gamma'}(\theta-\theta')} = K_{- \gamma, -
\gamma'}(-\theta + \theta').
\end{equation}

\begin{yizhu}
\begin{enumerate}[leftmargin=2em,itemsep=2pt,topsep=2pt]
\item \textbf{关于 $\sigma(\gamma)$}：\eqref{eq:defmugamma} 中的 $\sigma(\gamma)$ 是第 2.2 节引入的「二次加细」，即满足 $\sigma(\gamma_1)\sigma(\gamma_2) = (-1)^{\langle \gamma_1, \gamma_2\rangle} \sigma(\gamma_1+\gamma_2)$ 的函数 $\sigma:\clat\to\IZ_2$，故取值 $\pm1$；于是 $\log(-\sigma(\gamma))$ 只能取 $0$ 或 $\pm i\pi$，这正是原文「$\beta\mu_\gamma$ 就是 $i\varphi_\gamma$，或者与之相差 $\pm i\pi$」那句话的来由。% TODO 合稿改 \ref{sec:KS-formula}（第 2.2 节）
\item \textbf{与 \eqref{eq:gen-TBA} 的对照}：把 \eqref{eq:gen-TBA} 与 \eqref{eq:GMN-2} 并排比较时应注意，两式的约定并不逐字相同：\eqref{eq:gen-TBA} 对粒子种类 $b$ 求和，且 $\phi$ 之前带负号；\eqref{eq:GMN-2} 则对电荷 $\gamma'$ 求和、以 BPS 简并度 $\Omega(\gamma')$ 为权重，且 $K$ 之前带正号。
\end{enumerate}
\end{yizhu}


% 由 wcf-paper.bbl 机械转换（去掉 bbl 包装，反斜杠N -> 反斜杠CN）
\begin{thebibliography}{10}

\bibitem{Cecotti:1992qh}
S.~Cecotti, P.~Fendley, K.~A. Intriligator, and C.~Vafa, ``A new supersymmetric
  index,'' {\em Nucl. Phys.} {\bf B386} (1992) 405--452,
\href{http://www.arXiv.org/abs/hep-th/9204102}{{\tt hep-th/9204102}}.
%%CITATION = HEP-TH 9204102;%%.

\bibitem{Seiberg:1994rs}
N.~Seiberg and E.~Witten, ``Electric-magnetic duality, monopole condensation,
  and confinement in {$\CN=2$} supersymmetric {Y}ang-{M}ills theory,'' {\em
  Nucl. Phys.} {\bf B426} (1994) 19--52,
\href{http://www.arXiv.org/abs/hep-th/9407087}{{\tt hep-th/9407087}}.
%%CITATION = NUPHA,B426,19;%%.

\bibitem{Seiberg:1994aj}
N.~Seiberg and E.~Witten, ``Monopoles, duality and chiral symmetry breaking in
  {$\CN=2$} supersymmetric {Q}{C}{D},'' {\em Nucl. Phys.} {\bf B431} (1994)
  484--550,
\href{http://www.arXiv.org/abs/hep-th/9408099}{{\tt hep-th/9408099}}.
%%CITATION = NUPHA,B431,484;%%.

\bibitem{Denef:2007vg}
F.~Denef and G.~W. Moore, ``Split states, entropy enigmas, holes and halos,''
\href{http://www.arXiv.org/abs/hep-th/0702146}{{\tt hep-th/0702146}}.
%%CITATION = HEP-TH 0702146;%%.

\bibitem{Denef:2000nb}
F.~Denef, ``{Supergravity flows and D-brane stability},'' {\em JHEP} {\bf 08}
  (2000) 050,
\href{http://www.arXiv.org/abs/hep-th/0005049}{{\tt hep-th/0005049}}.
%%CITATION = HEP-TH/0005049;%%.

\bibitem{Denef:2000ar}
F.~Denef, ``{On the correspondence between D-branes and stationary supergravity
  solutions of type II Calabi-Yau compactifications},''
\href{http://www.arXiv.org/abs/hep-th/0010222}{{\tt hep-th/0010222}}.
%%CITATION = HEP-TH/0010222;%%.

\bibitem{ks1}
M.~Kontsevich and Y.~Soibelman, ``{S}tability structures, motivic
  {D}onaldson-{T}homas invariants and cluster transformations,''
  \href{http://www.arXiv.org/abs/0811.2435}{{\tt 0811.2435}}.

\bibitem{GarciaEtxebarria:2007zv}
I.~Garcia-Etxebarria and A.~M. Uranga, ``Non-perturbative superpotentials
  across lines of marginal stability,''
\href{http://www.arXiv.org/abs/0711.1430}{{\tt 0711.1430}}.
%%CITATION = ARXIV:0711.1430;%%.

\bibitem{Ooguri:1996me}
H.~Ooguri and C.~Vafa, ``{S}umming up {D}-instantons,'' {\em Phys. Rev. Lett.}
  {\bf 77} (1996) 3296--3298,
\href{http://www.arXiv.org/abs/hep-th/9608079}{{\tt hep-th/9608079}}.
%%CITATION = HEP-TH 9608079;%%.

\bibitem{Seiberg:1996ns}
N.~Seiberg and S.~H. Shenker, ``Hypermultiplet moduli space and string
  compactification to three dimensions,'' {\em Phys. Lett.} {\bf B388} (1996)
  521--523,
\href{http://www.arXiv.org/abs/hep-th/9608086}{{\tt hep-th/9608086}}.
%%CITATION = HEP-TH 9608086;%%.

\bibitem{Cecotti:1993rm}
S.~Cecotti and C.~Vafa, ``On classification of {$\CN=2$} supersymmetric
  theories,'' {\em Commun. Math. Phys.} {\bf 158} (1993) 569--644,
\href{http://www.arXiv.org/abs/hep-th/9211097}{{\tt hep-th/9211097}}.
%%CITATION = HEP-TH 9211097;%%.

\bibitem{Seiberg:1996nz}
N.~Seiberg and E.~Witten, ``Gauge dynamics and compactification to three
  dimensions,''
\href{http://www.arXiv.org/abs/hep-th/9607163}{{\tt hep-th/9607163}}.
%%CITATION = HEP-TH 9607163;%%.

\bibitem{Cecotti:1991me}
S.~Cecotti and C.~Vafa, ``Topological-antitopological fusion,'' {\em Nucl.
  Phys.} {\bf B367} (1991)
359--461.
%%CITATION = NUPHA,B367,359;%%.

\bibitem{Karlhede:1984vr}
A.~Karlhede, U.~Lindstrom, and M.~Ro\v{c}ek, ``Selfinteracting tensor
  multiplets in {$\CN=2$} superspace,'' {\em Phys. Lett.} {\bf B147} (1984)
297.
%%CITATION = PHLTA,B147,297;%%.

\bibitem{Lindstrom:1987ks}
U.~Lindstrom and M.~Rocek, ``New hyperkahler metrics and new supermultiplets,''
  {\em Commun. Math. Phys.} {\bf 115} (1988)
21.
%%CITATION = CMPHA,115,21;%%.

\bibitem{Lindstrom:2008gs}
U.~Lindstrom and M.~Rocek, ``{Properties of hyperkahler manifolds and their
  twistor spaces},''
\href{http://www.arXiv.org/abs/0807.1366}{{\tt 0807.1366}}.
%%CITATION = 0807.1366;%%.

\bibitem{MR2181810}
M.~Kontsevich and Y.~Soibelman, ``Affine structures and non-{A}rchimedean
  analytic spaces,'' in {\em The unity of mathematics}, vol.~244 of {\em Progr.
  Math.}, pp.~321--385.
\newblock Birkh\"auser Boston, Boston, MA, 2006.
\newblock \href{http://www.arXiv.org/abs/math.ag/0406564}{{\tt
  math.ag/0406564}}.

\bibitem{Joyce:2006pf}
D.~Joyce, ``{H}olomorphic generating functions for invariants counting coherent
  sheaves on {C}alabi-{Y}au 3-folds,''
\href{http://www.arXiv.org/abs/hep-th/0607039}{{\tt hep-th/0607039}}.
%%CITATION = HEP-TH 0607039;%%.

\bibitem{stab-stokes}
T.~Bridgeland and V.~Toledano~Laredo, ``{S}tability conditions and {S}tokes
  factors,'' \href{http://www.arXiv.org/abs/0801.3974}{{\tt 0801.3974}}.

\bibitem{MR2373143}
T.~Bridgeland, ``Stability conditions on triangulated categories,'' {\em Ann.
  of Math. (2)} {\bf 166} (2007), no.~2, 317--345.

\bibitem{Neitzke:2007ke}
A.~Neitzke, B.~Pioline, and S.~Vandoren, ``Twistors and black holes,'' {\em
  JHEP} {\bf 04} (2007) 038,
\href{http://www.arXiv.org/abs/hep-th/0701214}{{\tt hep-th/0701214}}.
%%CITATION = HEP-TH 0701214;%%.

\bibitem{Losev:1997tp}
A.~Losev, N.~A. Nekrasov, and S.~L. Shatashvili, ``Issues in topological gauge
  theory,'' {\em Nucl. Phys.} {\bf B534} (1998) 549--611,
\href{http://www.arXiv.org/abs/hep-th/9711108}{{\tt hep-th/9711108}}.
%%CITATION = HEP-TH 9711108;%%.

\bibitem{Losev:1999nt}
A.~Losev, N.~Nekrasov, and S.~L. Shatashvili, ``{The freckled instantons},''
\href{http://www.arXiv.org/abs/hep-th/9908204}{{\tt hep-th/9908204}}.
%%CITATION = HEP-TH/9908204;%%.

\bibitem{Ooguri:2005vr}
H.~Ooguri, C.~Vafa, and E.~P. Verlinde, ``{H}artle-{H}awking wave-function for
  flux compactifications,'' {\em Lett. Math. Phys.} {\bf 74} (2005) 311--342,
\href{http://www.arXiv.org/abs/hep-th/0502211}{{\tt hep-th/0502211}}.
%%CITATION = HEP-TH 0502211;%%.

\bibitem{Dorey:1999zk}
N.~Dorey, T.~J. Hollowood, and D.~Tong, ``{The BPS spectra of gauge theories in
  two and four dimensions},'' {\em JHEP} {\bf 05} (1999) 006,
\href{http://www.arXiv.org/abs/hep-th/9902134}{{\tt hep-th/9902134}}.
%%CITATION = HEP-TH/9902134;%%.

\bibitem{Ritz:2006rt}
A.~Ritz, ``{Superconformal R-charges and dyon multiplicities in N = 2 gauge
  theories},'' {\em Phys. Rev.} {\bf D75} (2007) 085008,
\href{http://www.arXiv.org/abs/hep-th/0612077}{{\tt hep-th/0612077}}.
%%CITATION = HEP-TH/0612077;%%.

\bibitem{Alexandrov:2008ds}
S.~Alexandrov, B.~Pioline, F.~Saueressig, and S.~Vandoren, ``{Linear
  perturbations of quaternionic metrics - {I}. The Hyperkahler case},''
\href{http://www.arXiv.org/abs/0806.4620}{{\tt 0806.4620}}.
%%CITATION = 0806.4620;%%.

\bibitem{Lerche:1997sm}
W.~Lerche, ``{Introduction to Seiberg-Witten theory and its stringy origin}.''
\newblock Prepared for CERN-Santiago de Compostela-La Plata Meeting on Trends
  in Theoretical Physics, CERN-Santiago de Compostela-La Plata, Argentina, 28
  Apr - 6 May 1997.

\bibitem{Freed:1997dp}
D.~S. Freed, ``{Special Kaehler manifolds},'' {\em Commun. Math. Phys.} {\bf
  203} (1999) 31--52,
\href{http://www.arXiv.org/abs/hep-th/9712042}{{\tt hep-th/9712042}}.
%%CITATION = HEP-TH/9712042;%%.

\bibitem{ias-volumes}
E.~Witten, ``Dynamics of quantum field theory,'' in {\em Quantum fields and
  strings: a course for mathematicians, Vol. 1, 2 (Princeton, NJ, 1996/1997)},
  pp.~1119--1424.
\newblock Amer. Math. Soc., Providence, RI, 1999.

\bibitem{Bilal:1996sk}
A.~Bilal and F.~Ferrari, ``{Curves of Marginal Stability and Weak and
  Strong-Coupling BPS Spectra in {$N=2$} Supersymmetric {QCD}},'' {\em Nucl.
  Phys.} {\bf B480} (1996) 589--622,
\href{http://www.arXiv.org/abs/hep-th/9605101}{{\tt hep-th/9605101}}.
%%CITATION = HEP-TH/9605101;%%.

\bibitem{Ferrari:1996sv}
A.~Bilal and F.~Ferrari, ``{The Strong-Coupling Spectrum of the Seiberg-Witten
  Theory},'' {\em Nucl. Phys.} {\bf B469} (1996) 387--402,
\href{http://www.arXiv.org/abs/hep-th/9602082}{{\tt hep-th/9602082}}.
%%CITATION = HEP-TH/9602082;%%.

\bibitem{Witten:1996hc}
E.~Witten, ``{Five-brane effective action in M-theory},'' {\em J. Geom. Phys.}
  {\bf 22} (1997) 103--133,
\href{http://www.arXiv.org/abs/hep-th/9610234}{{\tt hep-th/9610234}}.
%%CITATION = HEP-TH/9610234;%%.

\bibitem{Freed:2000ta}
D.~S. Freed, ``{Dirac charge quantization and generalized differential
  cohomology},''
\href{http://www.arXiv.org/abs/hep-th/0011220}{{\tt hep-th/0011220}}.
%%CITATION = HEP-TH/0011220;%%.

\bibitem{Hopkins:2002rd}
M.~J. Hopkins and I.~M. Singer, ``{Quadratic functions in geometry, topology,
  and M-theory},'' {\em J. Diff. Geom.} {\bf 70} (2005) 329--452,
\href{http://www.arXiv.org/abs/math/0211216}{{\tt math/0211216}}.
%%CITATION = MATH/0211216;%%.

\bibitem{Belov:2006jd}
D.~Belov and G.~W. Moore, ``{Holographic action for the self-dual field},''
\href{http://www.arXiv.org/abs/hep-th/0605038}{{\tt hep-th/0605038}}.
%%CITATION = HEP-TH/0605038;%%.

\bibitem{Hitchin:1986ea}
N.~J. Hitchin, A.~Karlhede, U.~Lindstrom, and M.~Ro\v{c}ek, ``Hyperk{\"a}hler
  metrics and supersymmetry,'' {\em Commun. Math. Phys.} {\bf 108} (1987)
535.
%%CITATION = CMPHA,108,535;%%.

\bibitem{MR1206066}
N.~Hitchin, ``Hyper-{K}\"ahler manifolds,'' {\em Ast\'erisque} (1992), no.~206,
  Exp.\ No.\ 748, 3, 137--166. S\'eminaire Bourbaki, Vol.\ 1991/92.

\bibitem{MR1213301}
B.~Dubrovin, ``Geometry and integrability of topological-antitopological
  fusion,'' {\em Commun. Math. Phys.} {\bf 152} (1993), no.~3, 539--564,
  \href{http://www.arXiv.org/abs/hep-th/9206037}{{\tt hep-th/9206037}}.

\bibitem{Argyres:1996eh}
P.~C. Argyres, M.~R. Plesser, and N.~Seiberg, ``{The Moduli Space of N=2 SUSY
  {QCD} and Duality in N=1 SUSY {QCD}},'' {\em Nucl. Phys.} {\bf B471} (1996)
  159--194,
\href{http://www.arXiv.org/abs/hep-th/9603042}{{\tt hep-th/9603042}}.
%%CITATION = HEP-TH/9603042;%%.

\bibitem{Cecotti:1993vy}
S.~Cecotti and C.~Vafa, ``Ising model and {$\CN=2$} supersymmetric theories,''
  {\em Commun. Math. Phys.} {\bf 157} (1993) 139--178,
\href{http://www.arXiv.org/abs/hep-th/9209085}{{\tt hep-th/9209085}}.
%%CITATION = HEP-TH 9209085;%%.

\bibitem{Argyres:1995jj}
P.~C. Argyres and M.~R. Douglas, ``New phenomena in {SU(3)} supersymmetric
  gauge theory,'' {\em Nucl. Phys.} {\bf B448} (1995) 93--126,
\href{http://www.arXiv.org/abs/hep-th/9505062}{{\tt hep-th/9505062}}.
%%CITATION = NUPHA,B448,93;%%.

\bibitem{Argyres:1995xn}
P.~C. Argyres, M.~Ronen~Plesser, N.~Seiberg, and E.~Witten, ``{New {$\CN=2$}
  Superconformal Field Theories in Four Dimensions},'' {\em Nucl. Phys.} {\bf
  B461} (1996) 71--84,
\href{http://www.arXiv.org/abs/hep-th/9511154}{{\tt hep-th/9511154}}.
%%CITATION = HEP-TH/9511154;%%.

\bibitem{Seiberg:1993vc}
N.~Seiberg, ``{Naturalness versus supersymmetric nonrenormalization
  theorems},'' {\em Phys. Lett.} {\bf B318} (1993) 469--475,
\href{http://www.arXiv.org/abs/hep-ph/9309335}{{\tt hep-ph/9309335}}.
%%CITATION = HEP-PH/9309335;%%.

\bibitem{Zamolodchikov:1989cf}
A.~B. Zamolodchikov, ``{Thermodynamic Bethe ansatz in relativistic models.
  Scaling three state Potts and Lee-Yang models},'' {\em Nucl. Phys.} {\bf
  B342} (1990)
695--720.
%%CITATION = NUPHA,B342,695;%%.

\bibitem{Nekrasov:2009ui}
N.~A. Nekrasov and S.~L. Shatashvili, ``{Quantum integrability and
  supersymmetric vacua},''
\href{http://www.arXiv.org/abs/0901.4748}{{\tt 0901.4748}}.
%%CITATION = 0901.4748;%%.

\end{thebibliography}


\end{document}
